% Lesson 40 — Speaking: Exchanges information on participating in training activities/courses on democracy — Anglais Tle A — Cours signé OnBuch+
% Programme officiel MINESEC. Compiler avec : tectonic EN40-speaking-exchanges-information-on-participating-in-trai.tex
\newcommand{\DOCMATIERE}{Anglais}
\newcommand{\DOCNIVEAU}{Tle A}
\newcommand{\DOCMODULE}{Chapter 4 : Citizenship and Human Rights}
\newcommand{\DOCLECON}{EN40}
\newcommand{\DOCTITRE}{Lesson 40 — Speaking: Exchanges information on participating in training activities/courses on democracy}
% OnBuch+ — préambule « cours signés » ENRICHI (compilé avec tectonic / XeLaTeX)
% Système visuel « Pop » d'OnBuch+ : fond crème (jamais blanc), bordures encre
% épaisses, ombres « dures » décalées, titres Archivo Black en capitales.
% Version MATHÉMATIQUES : graphes (pgfplots/tikz), boîtes pédagogiques
% (définition, propriété, méthode, attention, le savais-tu, exercice résolu,
% exercice, TP, bilan), page de garde et signature OnBuch+ ancrée en bas.
%
% Macros attendues dans main.tex AVANT \input{preamble} :
%   \DOCMATIERE  \DOCNIVEAU  \DOCMODULE  \DOCLECON (numéro)  \DOCTITRE
\documentclass[11pt]{article}

\usepackage{fontspec}
\usepackage[english,french]{babel} % français = langue principale ; l'anglais via \eng{..} / textebox
\usepackage[a4paper,margin=2cm,top=3.4cm,bottom=2.8cm,headheight=1.4cm,footskip=1.3cm]{geometry}
\usepackage{amsmath,amssymb,amsfonts}
\usepackage{mathtools} % \xmapsto, \coloneqq... souvent utilisés par les modèles
\usepackage{cancel} % simplifications barrées (bilans de forces)
% \abs{z} (module, valeur absolue) et \norm{u} (norme), utilisés par les modèles
\providecommand{\abs}{}\let\abs\relax\DeclarePairedDelimiter{\abs}{\lvert}{\rvert}
\providecommand{\norm}{}\let\norm\relax\DeclarePairedDelimiter{\norm}{\lVert}{\rVert}
\usepackage[table]{xcolor} % [table] : fournit aussi \rowcolors (alternance de lignes)
\usepackage{pagecolor}
\usepackage{tikz}
\usetikzlibrary{arrows.meta,positioning,calc,shapes.geometric,shapes.misc,shapes.arrows,decorations.pathmorphing,decorations.pathreplacing,decorations.markings,patterns,shadows,mindmap,trees,fit,backgrounds,matrix,intersections,angles,quotes}
\usepackage{pgfplots}
\usepackage[version=4]{mhchem} % équations nucléaires \ce{^{235}_{92}U}
\usepackage[european,siunitx]{circuitikz} % schémas électriques
\pgfplotsset{compat=1.18}
\usepgfplotslibrary{fillbetween,groupplots} % aires entre courbes (calcul intégral)
\usepackage{enumitem}
\usepackage{fancyhdr}
% [most] charge toutes les bibliothèques tcolorbox (listings, minted, raster,
% documentation...) et gonfle inutilement la mémoire TeX sur un document long
% (~300 boîtes) : on ne charge que ce qui est réellement utilisé.
\usepackage[skins,breakable]{tcolorbox}
\usepackage{graphicx}
\usepackage{colortbl}
\usepackage{booktabs}
\usepackage{tabularx}
\usepackage{diagbox} % case d'en-tête diagonale des tableaux à double entrée
% Colonnes extensibles centrée / gauche / droite (C, L, R), souvent utilisées par les modèles
\newcolumntype{C}{>{\centering\arraybackslash}X}
\newcolumntype{L}{>{\raggedright\arraybackslash}X}
\newcolumntype{R}{>{\raggedleft\arraybackslash}X}
\usepackage{array}
\usepackage{multicol}
\usepackage{multirow}
\usepackage{siunitx}
% parse-numbers=false : accepte aussi \SI{2,5 \times 10^{-3}}{...}
\sisetup{locale=FR,output-decimal-marker={,},group-minimum-digits=4,parse-numbers=false}
\DeclareSIUnit\ohm{\ensuremath{\symup{\Omega}}} % U+2126 absent de la police
\DeclareSIUnit\division{div} % réglages d'oscilloscope (ms/div, V/div)
\DeclareSIUnit\tr{tr} % tours (vitesse de rotation, ex. \SI{1500}{\tr\per\minute})
\DeclareSIUnit\molar{M} % concentration molaire (ex. \SI{3}{\molar}, \SI{1}{\micro\molar})
\DeclareSIUnit\an{an} % année (le modèle confond parfois avec \year, un compteur TeX) — ex. \SI{5}{\kilo\meter\per\an}
\DeclareSIUnit\FCFA{FCFA} % franc CFA (ex. \SI{500}{\FCFA\per\kilo\gram})
\DeclareSIUnit\degreeBrix{\textdegree Brix} % teneur en sucre d'un sirop/jus (réfractométrie)
\DeclareSIUnit\month{mois} % le modèle utilise parfois \month, un compteur TeX, comme unité de durée
\usepackage{qrcode}
% \degree : commande fréquemment utilisée par les modèles pour un angle ou une
% température en dehors de \SI{}{} (non fournie par les paquets chargés).
\providecommand{\degree}{^{\circ}}
% \qed (fin de démonstration), utilisé par les modèles sans amsthm
\providecommand{\qed}{\hfill$\square$}
% \barre : double barre des valeurs interdites dans un tableau de variations (tkz-tab)
\providecommand{\barre}{\ensuremath{\|}}
% environnement proof (amsthm non chargé) utilisé par les modèles
\ifdefined\proof\else\newenvironment{proof}[1][Démonstration]{\par\noindent\textit{#1.}\ }{\qed\par}\fi
\usepackage{lastpage}
\usepackage{hyperref}
\hypersetup{hidelinks}

% ============================================================
% Palette « Pop » OnBuch+ — mêmes hex que la classe P (plus_ui.dart)
% ============================================================
\definecolor{poporange}{HTML}{F59321}
\definecolor{popdark}{HTML}{E07A0C}
\definecolor{popcream}{HTML}{FFF6E8}
\definecolor{popink}{HTML}{1C1714}
\definecolor{poppurple}{HTML}{7A5AE0}
\definecolor{popgreen}{HTML}{2E9E6A}
\definecolor{poppink}{HTML}{E0457B}
\definecolor{popblue}{HTML}{2F7DE1}
\definecolor{popgold}{HTML}{C98A12}
\definecolor{popmuted}{HTML}{8A7F76}
\definecolor{popline}{HTML}{EADFD2}
\definecolor{popgray}{HTML}{8A7F76}
\colorlet{popgrey}{popgray}
% Alias tolérants (le modèle nomme parfois les couleurs différemment)
\colorlet{popwhite}{white}
\colorlet{popblack}{popink}
\colorlet{popred}{poppink}
\colorlet{popyellow}{popgold}
% Teintes claires (fonds de tableaux/encadrés) — le modèle nomme parfois
% les couleurs avec un suffixe L (ex. popblueL) sans les avoir définies.
\colorlet{poporangeL}{poporange!20}
\colorlet{popdarkL}{popdark!20}
\colorlet{poppurpleL}{poppurple!20}
\colorlet{popgreenL}{popgreen!20}
\colorlet{poppinkL}{poppink!20}
\colorlet{popblueL}{popblue!20}
\colorlet{popgoldL}{popgold!20}
\colorlet{popredL}{popred!20}
\colorlet{popgrayL}{popgray!20}
% Teintes claires pour fonds de tableaux / zones
\colorlet{popblueL}{popblue!10!white}
\colorlet{poporangeL}{poporange!14!white}
\colorlet{popgreenL}{popgreen!12!white}
\colorlet{poppurpleL}{poppurple!12!white}
\colorlet{poppinkL}{poppink!12!white}
\colorlet{popgoldL}{popgold!14!white}

\pagecolor{popcream}

% ============================================================
% Typographie
% ============================================================
\defaultfontfeatures{Ligatures=TeX}
\setmainfont{PlusJakartaSans-Medium.ttf}[Path=../../../fonts/,BoldFont=PlusJakartaSans-Bold.ttf]
% Symboles, API et emojis absents de la police principale : repli sur DejaVu Sans (caractères actifs)
\newfontface\symfont{DejaVuSans.ttf}[Path=../../../fonts/]
\makeatletter
\newcommand\symactive[1]{\catcode#1=13 \begingroup\lccode`\~=#1 \lowercase{\endgroup\def~}{{\symfont\char#1}}}
\newcommand\symblank[1]{\catcode#1=13 \begingroup\lccode`\~=#1 \lowercase{\endgroup\def~}{}}
\newcommand\symhyph[1]{\catcode#1=13 \begingroup\lccode`\~=#1 \lowercase{\endgroup\def~}{\mbox{-}}}
\newcount\symc
\newcommand\symrange[2]{\symc=#1 \@whilenum\symc<#2 \do{\expandafter\symactive\expandafter{\the\symc}\advance\symc by 1 }}
\newcommand\symblankrange[2]{\symc=#1 \@whilenum\symc<#2 \do{\expandafter\symblank\expandafter{\the\symc}\advance\symc by 1 }}
\makeatother
\symrange{"0250}{"02FF}\symactive{"2713}\symactive{"2714}\symactive{"2717}\symactive{"2718}\symactive{"2610}\symactive{"2611}\symactive{"2612}
\symhyph{"2011}\symhyph{"2E1A}\symblankrange{"1F300}{"1FAFF}\symblank{"FE0F}
% Maths en sans-serif (Fira Math) avec chiffres et lettres droites de Plus
% Jakarta, pour que les formules se fondent dans le texte.
\usepackage[math-style=ISO,bold-style=ISO]{unicode-math}
\setmathfont{FiraMath-Regular.otf}[Path=../../../fonts/]
\setmathfont{PlusJakartaSans-Medium.ttf}[Path=../../../fonts/,range={up/num,up/{Latin,latin},\mathup}]
\usepackage{icomma} % virgule décimale sans espace en mode math
% Caractères mathématiques Unicode absents des polices de texte (Plus Jakarta,
% Archivo Black) : rendus par leur équivalent mathématique, en texte comme en
% formule — sinon ils s'affichent en carré vide (ex. « Arithmétique dans ℤ »).
\usepackage{newunicodechar}
\newunicodechar{ℝ}{\ensuremath{\mathbb{R}}}
\newunicodechar{ℤ}{\ensuremath{\mathbb{Z}}}
\newunicodechar{ℂ}{\ensuremath{\mathbb{C}}}
\newunicodechar{ℕ}{\ensuremath{\mathbb{N}}}
\newunicodechar{ℚ}{\ensuremath{\mathbb{Q}}}
\newunicodechar{𝔻}{\ensuremath{\mathbb{D}}}
\newunicodechar{α}{\ensuremath{\alpha}}
\newunicodechar{β}{\ensuremath{\beta}}
\newunicodechar{γ}{\ensuremath{\gamma}}
\newunicodechar{δ}{\ensuremath{\delta}}
\newunicodechar{ε}{\ensuremath{\varepsilon}}
\newunicodechar{θ}{\ensuremath{\theta}}
\newunicodechar{λ}{\ensuremath{\lambda}}
\newunicodechar{μ}{\ensuremath{\mu}}
\newunicodechar{σ}{\ensuremath{\sigma}}
\newunicodechar{τ}{\ensuremath{\tau}}
\newunicodechar{φ}{\ensuremath{\varphi}}
\newunicodechar{ψ}{\ensuremath{\psi}}
\newunicodechar{ω}{\ensuremath{\omega}}
\newunicodechar{Σ}{\ensuremath{\Sigma}}
% majuscules produites par \MakeUppercase dans les titres (α-aminés -> Α)
\newunicodechar{Α}{\ensuremath{\alpha}}
\newunicodechar{Β}{\ensuremath{\beta}}
\newunicodechar{Γ}{\ensuremath{\Gamma}}
\newunicodechar{Δ}{\ensuremath{\Delta}}
\newunicodechar{Ε}{\ensuremath{\varepsilon}}
\newunicodechar{Θ}{\ensuremath{\Theta}}
\newunicodechar{Λ}{\ensuremath{\Lambda}}
\newunicodechar{Μ}{\ensuremath{\mu}}
\newunicodechar{Τ}{\ensuremath{\tau}}
\newunicodechar{Φ}{\ensuremath{\Phi}}
\newunicodechar{Ψ}{\ensuremath{\Psi}}
\newunicodechar{Ω}{\ensuremath{\Omega}}
\newunicodechar{↦}{\ensuremath{\mapsto}}
\newunicodechar{∈}{\ensuremath{\in}}
\newunicodechar{∉}{\ensuremath{\notin}}
\newunicodechar{∧}{\ensuremath{\wedge}}
\newunicodechar{⇔}{\ensuremath{\Leftrightarrow}}
\newunicodechar{⟺}{\ensuremath{\Longleftrightarrow}}
\newunicodechar{⇒}{\ensuremath{\Rightarrow}}
\newunicodechar{⟹}{\ensuremath{\Longrightarrow}}
\newunicodechar{∘}{\ensuremath{\circ}}
\newunicodechar{∪}{\ensuremath{\cup}}
\newunicodechar{∩}{\ensuremath{\cap}}
\newunicodechar{≡}{\ensuremath{\equiv}}
\newunicodechar{⊂}{\ensuremath{\subset}}
\newunicodechar{⊥}{\ensuremath{\perp}}
\newunicodechar{∃}{\ensuremath{\exists}}
\newunicodechar{∀}{\ensuremath{\forall}}
\newunicodechar{∛}{\ensuremath{\sqrt[3]{\,}}}
\newunicodechar{∜}{\ensuremath{\sqrt[4]{\,}}}
\newunicodechar{⃗}{\ensuremath{{}^{\to}}}
\newunicodechar{ⁿ}{\ifmmode^{n}\else\textsuperscript{\ensuremath{n}}\fi}
\newunicodechar{ˣ}{\ifmmode^{x}\else\textsuperscript{\ensuremath{x}}\fi}
\newunicodechar{ᵇ}{\ifmmode^{b}\else\textsuperscript{\ensuremath{b}}\fi}
\newunicodechar{ʳ}{\ifmmode^{r}\else\textsuperscript{\ensuremath{r}}\fi}
\newunicodechar{ᵘ}{\ifmmode^{u}\else\textsuperscript{\ensuremath{u}}\fi}
\newunicodechar{ʸ}{\ifmmode^{y}\else\textsuperscript{\ensuremath{y}}\fi}
\newunicodechar{ᵃ}{\ifmmode^{a}\else\textsuperscript{\ensuremath{a}}\fi}
\newunicodechar{ᵉ}{\ifmmode^{e}\else\textsuperscript{\ensuremath{e}}\fi}
\newunicodechar{ᵏ}{\ifmmode^{k}\else\textsuperscript{\ensuremath{k}}\fi}
\newunicodechar{ᵖ}{\ifmmode^{p}\else\textsuperscript{\ensuremath{p}}\fi}
\newunicodechar{ᴺ}{\ifmmode^{N}\else\textsuperscript{\ensuremath{N}}\fi}
\newunicodechar{ᶜ}{\ifmmode^{c}\else\textsuperscript{\ensuremath{c}}\fi}
\newunicodechar{ʰ}{\ifmmode^{h}\else\textsuperscript{\ensuremath{h}}\fi}
\newunicodechar{ᵠ}{\ifmmode^{q}\else\textsuperscript{\ensuremath{q}}\fi}
\newunicodechar{ₙ}{\ifmmode_{n}\else\textsubscript{\ensuremath{n}}\fi}
\newunicodechar{ₐ}{\ifmmode_{a}\else\textsubscript{\ensuremath{a}}\fi}
\newunicodechar{ᵢ}{\ifmmode_{i}\else\textsubscript{\ensuremath{i}}\fi}
\newunicodechar{ᵣ}{\ifmmode_{r}\else\textsubscript{\ensuremath{r}}\fi}
\newunicodechar{ₖ}{\ifmmode_{k}\else\textsubscript{\ensuremath{k}}\fi}
\newunicodechar{ᵦ}{\ifmmode_{\beta}\else\textsubscript{\ensuremath{\beta}}\fi}
\newunicodechar{ₓ}{\ifmmode_{x}\else\textsubscript{\ensuremath{x}}\fi}
\newfontfamily\popdisplay{ArchivoBlack-Regular.ttf}[Path=../../../fonts/]
\newfontfamily\poplogo{SpaceGrotesk-Bold.ttf}[Path=../../../fonts/]
\newfontfamily\popsemi{PlusJakartaSans-SemiBold.ttf}[Path=../../../fonts/]
\newfontfamily\popxbold{PlusJakartaSans-ExtraBold.ttf}[Path=../../../fonts/]
\newfontfamily\popxboldls{PlusJakartaSans-ExtraBold.ttf}[Path=../../../fonts/,LetterSpace=3.0]

\setlist[enumerate]{leftmargin=1.7em,itemsep=5pt,topsep=4pt}
\setlist[itemize]{leftmargin=1.7em,itemsep=4pt,topsep=4pt,label={\color{poporange}\textbullet}}
\setlength{\parindent}{0pt}
\setlength{\parskip}{5pt}
\renewcommand{\arraystretch}{1.25}

% pgfplots : style Pop par défaut
\pgfplotsset{
  every axis/.append style={axis line style={popink,line width=1pt},tick style={popink},
    grid style={popline,line width=0.5pt},label style={font=\small},tick label style={font=\footnotesize}},
  popcurve/.style={poporange,line width=1.8pt,no markers,smooth},
}
% Même style disponible dans un \draw TikZ simple (hors pgfplots)
\tikzset{popcurve/.style={poporange,line width=1.8pt}}
\tikzset{popgrid/.style={popline,very thin}}
\colorlet{popcurve}{poporange} % le nom du style est parfois utilisé comme couleur

% ============================================================
% Pill / squiggle
% ============================================================
\newcommand{\poppill}[3]{%
  \tikz[baseline=-0.55ex]\node[draw=popink,line width=1.2pt,rounded corners=2.6pt,%
    fill=#2,inner xsep=6.5pt,inner ysep=3pt]{%
    \popxboldls\fontsize{7}{7}\selectfont\color{#3}\MakeUppercase{#1}};%
}
\newcommand{\popsquiggle}[1]{%
  \tikz[baseline=-0.4ex]{
    \draw[#1,line width=2.6pt,line cap=round]
      plot [smooth] coordinates {(0,0.09) (0.34,0.01) (0.68,0.21) (1.02,0.01) (1.36,0.10)};
  }%
}
% Mot-clé surligné (marqueur orange sous le texte)
\newcommand{\cle}[1]{{\popxbold\color{popdark}#1}}

% ============================================================
% En-tête / pied
% ============================================================
\pagestyle{fancy}
\fancyhf{}
\renewcommand{\headrulewidth}{0pt}
\renewcommand{\footrulewidth}{0pt}
\lhead{\raisebox{-0.15ex}{\poplogo\fontsize{13}{13}\selectfont\color{popink}OnBuch\color{poporange}+}}
\rhead{\poppill{\DOCMATIERE\ \textbullet\ \DOCNIVEAU\ \textbullet\ Leçon \DOCLECON}{white}{popink}}
\lfoot{\popxbold\fontsize{7.6}{7.6}\selectfont\color{popmuted}\MakeUppercase{Cours signé OnBuch+}\ \textbullet\ {\popsemi\DOCTITRE}}
\rfoot{\tikz[baseline=-0.6ex]\node[rounded corners=8.5pt,fill=popink,minimum height=17pt,minimum width=17pt,inner xsep=5pt,inner sep=0]{\popxbold\fontsize{8}{8}\selectfont\color{popcream}\thepage\,/\,\pageref*{LastPage}};}

% ============================================================
% Titres
% ============================================================
\newcommand{\chaptitre}[2]{%
  \begin{tcolorbox}[enhanced,colback=poporange,colframe=popink,boxrule=2.6pt,arc=11pt,
      drop shadow=popink,left=15pt,right=15pt,top=12pt,bottom=11pt,before skip=3pt,after skip=18pt]
    \poppill{#2}{popink}{popcream}\par\vspace{9pt}
    {\raggedright\hyphenpenalty=10000\exhyphenpenalty=10000\popdisplay\fontsize{22}{24}\selectfont\color{popcream}\MakeUppercase{#1}\par}
  \end{tcolorbox}
}
% Section numérotée : pastille orange avec le numéro + titre Archivo
\newcounter{popsec}
\newcommand{\coursec}[1]{%
  \stepcounter{popsec}\setcounter{popsub}{0}%
  \par\vspace{16pt}\noindent
  \begin{minipage}{\linewidth}
  \tikz[baseline=-0.7ex]\node[draw=popink,line width=1.6pt,fill=poporange,rounded corners=4pt,minimum size=22pt,inner sep=0]{\popdisplay\fontsize{12}{12}\selectfont\color{popcream}\thepopsec};%
  \hspace{8pt}{\popdisplay\fontsize{15}{17}\selectfont\color{popink}\MakeUppercase{#1}}\par
  \vspace{4pt}\noindent{\color{poporange}\rule{36pt}{3.6pt}}
  \end{minipage}\par\nopagebreak\vspace{8pt}
}
\newcounter{popsub}
\newcommand{\courssub}[1]{%
  \stepcounter{popsub}%
  \par\vspace{9pt}\noindent{\popxbold\fontsize{11.5}{13}\selectfont\color{popink}{\color{poporange}\thepopsec.\thepopsub}\hspace{6pt}#1}\par\nopagebreak\vspace{4pt}
}

% ============================================================
% Boîtes pédagogiques — même recette : fond blanc, bordure épaisse colorée,
% ombre dure encre, badge plein. Titre optionnel [#1].
% ============================================================
\tcbset{popbase/.style={breakable,enhanced,colback=white,boxrule=2.3pt,arc=9pt,
  shadow={3pt}{-3pt}{0pt}{popink},left=14pt,right=14pt,top=11pt,bottom=11pt,before skip=11pt,after skip=15pt,
  fontupper=\popsemi\fontsize{10.6}{14.5}\selectfont\color{popink},
  fontlower=\popsemi\fontsize{10.6}{14.5}\selectfont\color{popink}}}
% Coche dessinée (le glyphe ✓ n'existe pas dans les polices du document)
\newcommand{\popcheck}[1]{\tikz[baseline=-0.5ex]\draw[#1,line width=1.6pt,line cap=round,line join=round](-0.13,0.0)--(-0.04,-0.09)--(0.13,0.1);}
\providecommand{\coche}{\popcheck{popgreen}}
\providecommand{\resultat}[1]{\par\smallskip\noindent\fcolorbox{popgreen}{popgreenL}{\parbox{\dimexpr\linewidth-2\fboxsep-2\fboxrule\relax}{\textbf{Résultat.} #1}}\par\smallskip}
\newcommand{\popboxhead}[3]{\poppill{#1}{#2}{white}\ifx\relax#3\relax\else\hspace{6pt}{\popxbold\fontsize{10.6}{12}\selectfont\color{popink}#3}\fi\par\vspace{7pt}}

\newtcolorbox{aretenir}[1][]{popbase,colframe=popgreen,before upper={\popboxhead{À retenir}{popgreen}{#1}}}
% Texte en anglais (document de lecture, dialogue, extrait) : cadre
% d'étude. Un titre optionnel (source, auteur) s'affiche dans l'en-tête.
\newtcolorbox{textebox}[1][]{popbase,colframe=popblue,before upper={\popboxhead{Text}{popblue}{#1}\begin{otherlanguage}{english}},after upper={\end{otherlanguage}}}
% Marques ✓ / ✗ (forme correcte / fautive) souvent utilisées par les modèles
\providecommand{\cross}{\textcolor{poppink}{\ensuremath{\times}}}
\providecommand{\xmark}{\cross}\providecommand{\crossmark}{\cross}\providecommand{\cmark}{\ensuremath{\checkmark}}
% Ligne de réponse / justification à compléter (inventée par les modèles)
\providecommand{\justification}[1]{\par\noindent\textit{Justification~:} \dotfill\par}
% \textipa (package tipa non chargé) : la police ne couvre pas l'API -> simple passage
\providecommand{\textipa}[1]{#1}
% Soulignement (ulem non chargé) : \uline, \ul
\providecommand{\uline}[1]{\underline{#1}}\providecommand{\ul}[1]{\underline{#1}}
% \glossaire{\item ...} inventée par les modèles : liste de glossaire
\providecommand{\glossaire}[1]{\begin{itemize}#1\end{itemize}}
% \alert (beamer) inventée par les modèles : texte mis en évidence
\providecommand{\alert}[1]{\textbf{\textcolor{poppink}{#1}}}
% variantes ASCII d'ordinaux écrites en macro
\providecommand{\iere}{\textsuperscript{re}}\providecommand{\ieme}{\textsuperscript{e}}\providecommand{\ere}{\textsuperscript{re}}\providecommand{\eme}{\textsuperscript{e}}\providecommand{\er}{\textsuperscript{er}}\providecommand{\ie}{\textsuperscript{e}}
% Ordinaux français écrits en macro par les modèles : 1\ière, 3\ième
\newcommand{\ière}{\textsuperscript{re}}\newcommand{\ième}{\textsuperscript{e}}
% \exemple (inventée par les modèles) : étiquette « Exemple : »
\providecommand{\exemple}{\textbf{Exemple~:}\ }
% \square utilisable aussi hors mode math (cases à cocher)
\AtBeginDocument{\let\squareMath\square\renewcommand{\square}{\ensuremath{\squareMath}}}
% Mot ou phrase en anglais à l'intérieur d'un paragraphe français
\newcommand{\eng}[1]{\ifmmode\text{\textit{#1}}\else\begingroup\selectlanguage{english}\itshape #1\endgroup\fi}
\let\esp\eng % alias toléré
% flèches d'intonation inventées par les modèles
\providecommand{\textsearrow}{}\renewcommand{\textsearrow}{\ensuremath{\searrow}}\providecommand{\textnearrow}{}\renewcommand{\textnearrow}{\ensuremath{\nearrow}}
% \fr{..} : mot français (inventé par les modèles)
\providecommand{\fr}[1]{\textit{#1}}
\newtcolorbox{exemplebox}[1][]{popbase,colframe=poporange,before upper={\popboxhead{Exemple}{poporange}{#1}}}
\newtcolorbox{definition}[1][]{popbase,colframe=popblue,before upper={\popboxhead{Définition}{popblue}{#1}}}
\newtcolorbox{propriete}[1][]{popbase,colframe=poppurple,before upper={\popboxhead{Propriété}{poppurple}{#1}}}
\newtcolorbox{methode}[1][]{popbase,colframe=popgold,before upper={\popboxhead{Méthode}{popgold}{#1}}}
\newtcolorbox{attention}[1][]{popbase,colframe=poppink,before upper={\popboxhead{Attention}{poppink}{#1}}}
\newtcolorbox{experience}[1][]{popbase,colframe=popblue,colback=popblueL,before upper={\popboxhead{Expérience}{popblue}{#1}}}
% « Le savais-tu ? » : carte violette pleine (contexte, culture, Cameroun)
\newtcolorbox{savaistu}[1][]{popbase,colback=poppurple,colframe=popink,
  fontupper=\popsemi\fontsize{10.4}{14.2}\selectfont\color{white},
  before upper={\poppill{Le savais-tu\,?}{white}{poppurple}\ifx\relax#1\relax\else\hspace{6pt}{\popxbold\color{white}#1}\fi\par\vspace{7pt}}}
% Exercice résolu : énoncé en haut, \tcblower puis la solution sur fond crème
\newcounter{popexo}\newcounter{popexr}
\newtcolorbox{exoresolu}[1][]{popbase,colframe=poporange,colbacklower=popcream,
  segmentation style={popink,line width=1.2pt,dashed},
  before upper={\stepcounter{popexr}\popboxhead{Exercice résolu \thepopexr}{poporange}{#1}},
  before lower={\poppill{Solution}{popink}{popcream}\par\vspace{6pt}}}
% Exercice (non résolu en place) — la correction est dans la partie Corrigés
\newtcolorbox{exercice}[1][]{popbase,colframe=popink,
  before upper={\stepcounter{popexo}\popboxhead{Exercice \thepopexo}{popink}{#1}}}
% Corrigé
\newtcolorbox{corrige}[1][]{popbase,colframe=popgreen,colback=popgreenL,
  before upper={\popboxhead{Corrigé}{popgreen}{#1}}}
% Objectifs / prérequis / bilan
\newtcolorbox{objectifs}{popbase,colframe=popink,colback=poporangeL,
  before upper={\popboxhead{Objectifs de la leçon}{popink}{}}}
\newtcolorbox{prerequis}{popbase,colframe=popink,colback=popblueL,
  before upper={\popboxhead{Prérequis}{popblue}{}}}
\newtcolorbox{bilan}{popbase,colframe=popink,colback=white,boxrule=2.8pt,
  before upper={\popboxhead{Fiche bilan}{poporange}{L'essentiel à connaître pour l'examen}}}
% Figure encadrée avec légende
\newtcolorbox{popfigure}[1][]{enhanced,breakable=false,colback=white,colframe=popline,boxrule=1.4pt,arc=8pt,
  left=10pt,right=10pt,top=10pt,bottom=8pt,before skip=10pt,after skip=12pt,halign=center,
  lower separated=false,halign lower=center,
  fontlower=\popsemi\fontsize{9}{11.5}\selectfont\color{popmuted},
  #1}
\newcommand{\legende}[1]{\par\vspace{6pt}{\popsemi\fontsize{9}{11.5}\selectfont\color{popmuted}#1}}

% ============================================================
% Signature OnBuch+
% ============================================================
\newcommand{\poplogoinline}[1]{{\poplogo\fontsize{#1}{#1}\selectfont\color{popink}OnBuch\color{poporange}+}}
\newcommand{\signatureonbuch}{%
  \begin{tcolorbox}[enhanced,colback=popink,colframe=popink,boxrule=2.2pt,arc=10pt,
      drop shadow=poporange,left=14pt,right=14pt,top=10pt,bottom=10pt,before skip=0pt,after skip=0pt]
    \tikz[baseline=-0.7ex]\node[circle,fill=popgreen,draw=popcream,line width=1.3pt,minimum size=22pt,inner sep=0]{\popcheck{white}};%
    \hspace{9pt}\begin{minipage}[c]{0.62\linewidth}
      {\popxbold\fontsize{12}{13}\selectfont\color{popcream}Signé\ \textbullet\ {\poplogo OnBuch\color{poporange}+}}\par\vspace{2pt}
      {\raggedright\popsemi\fontsize{8}{10.5}\selectfont\color{popline}\MakeUppercase{Équipe pédagogique OnBuch+}\\\MakeUppercase{Conforme au programme officiel MINESEC}\par}
    \end{minipage}\hfill
    {\poplogo\fontsize{22}{22}\selectfont\color{popcream}OnBuch\color{poporange}+}
  \end{tcolorbox}%
}

% ============================================================
% Page de garde
% #1 titre · #2 matière · #3 niveau · #4 module · #5 n° leçon · #6 durée ·
% #7 description / objectifs (texte ou itemize)
% ============================================================
\newcommand{\pagegarde}[7]{%
  \thispagestyle{empty}
  \newgeometry{a4paper,margin=1.7cm,top=1.6cm,bottom=1.4cm}%
  \begin{tikzpicture}[remember picture,overlay]
    \fill[poporange!18] ([xshift=-3.2cm,yshift=-3.6cm]current page.north east) circle (4.6cm);
    \fill[poppurple!14] ([xshift=2.4cm,yshift=7.6cm]current page.south west) circle (3.8cm);
  \end{tikzpicture}%
  {\poplogo\fontsize{30}{30}\selectfont\color{popink}OnBuch\color{poporange}+}\hfill
  \raisebox{0.6ex}{\poppill{Cours signé \textbullet\ Premium}{popink}{popcream}}\par
  \vspace{1pt}\popsquiggle{poppurple}\par
  \vspace{8pt}
  \begin{tcolorbox}[enhanced,colback=poporange,colframe=popink,boxrule=2.8pt,arc=13pt,
      drop shadow=popink,left=18pt,right=18pt,top=12pt,bottom=13pt,after skip=10pt]
    \poppill{#2 \textbullet\ #3}{popink}{popcream}\hspace{5pt}\poppill{#4}{white}{popink}\par\vspace{8pt}
    {\popdisplay\fontsize{14}{14}\selectfont\color{popink}LEÇON #5}\par\vspace{4pt}
    {\raggedright\hyphenpenalty=10000\exhyphenpenalty=10000\popdisplay\fontsize{25}{27}\selectfont\color{popcream}\MakeUppercase{#1}\par}
    \vspace{8pt}
    {\popxbold\fontsize{9.5}{9.5}\selectfont\color{popink}\MakeUppercase{Durée conseillée : #6}}
  \end{tcolorbox}
  \begin{tcolorbox}[enhanced,colback=white,colframe=popink,boxrule=2.2pt,arc=10pt,
      drop shadow=popink,left=16pt,right=16pt,top=10pt,bottom=10pt,after skip=8pt]
    \poppill{Dans cette leçon}{poppurple}{white}\par\vspace{5pt}
    \popsemi\fontsize{9.3}{12.6}\selectfont\color{popink}#7
  \end{tcolorbox}
  \noindent\begin{minipage}[t]{0.487\linewidth}
    \begin{tcolorbox}[enhanced,colback=popgreen,colframe=popink,boxrule=2.2pt,arc=10pt,
        drop shadow=popink,left=11pt,right=11pt,top=10pt,bottom=10pt,height=3.1cm,valign=center]
      \tcbox[colback=white,colframe=popink,boxrule=1.3pt,arc=5pt,boxsep=0pt,nobeforeafter,
        left=3pt,right=3pt,top=3pt,bottom=3pt,valign=center]{\qrcode[height=1.9cm]{https://play.google.com/store/apps/details?id=com.luvvix.onbuch&hl=fr}}%
      \hspace{8pt}\begin{minipage}[c]{0.5\linewidth}
        {\popdisplay\fontsize{11}{12.5}\selectfont\color{white}\MakeUppercase{Télécharge OnBuch}}\par\vspace{4pt}
        {\raggedright\popsemi\fontsize{8.4}{11}\selectfont\color{white}Gratuit \textbullet\ Google Play\\Tuteur IA, annales, concours, résultats\par}
      \end{minipage}
    \end{tcolorbox}
  \end{minipage}\hfill
  \begin{minipage}[t]{0.487\linewidth}
    \begin{tcolorbox}[enhanced,colback=poppurple,colframe=popink,boxrule=2.2pt,arc=10pt,
        drop shadow=popink,left=11pt,right=11pt,top=10pt,bottom=10pt,height=3.1cm,valign=center]
      \tikz[baseline=-0.7ex]\node[circle,fill=white,draw=popink,line width=1.3pt,minimum size=1.6cm,inner sep=0]{\popdisplay\fontsize{24}{24}\selectfont\color{poppurple}+};%
      \hspace{8pt}\begin{minipage}[c]{0.6\linewidth}
        {\popdisplay\fontsize{11}{12.5}\selectfont\color{white}\MakeUppercase{Passe à OnBuch+}}\par\vspace{4pt}
        {\raggedright\popsemi\fontsize{8.4}{11}\selectfont\color{white}Tous les cours signés, Tuteur IA illimité, fiches PDF — onglet OnBuch+ de l'app\par}
      \end{minipage}
    \end{tcolorbox}
  \end{minipage}
  \vspace{8pt}
  \popcontenu
  \vfill
  \signatureonbuch
  \restoregeometry
  \newpage
}

% Rangée « Ce cours contient » (page de garde)
\newcommand{\poptile}[3]{% #1 couleur #2 gros texte #3 légende
  \begin{tcolorbox}[enhanced,colback=#1,colframe=popink,boxrule=2pt,arc=9pt,shadow={2.5pt}{-2.5pt}{0pt}{popink},
      left=6pt,right=6pt,top=9pt,bottom=9pt,halign=center,valign=center,height=1.75cm,width=\linewidth,nobeforeafter]
    {\popdisplay\fontsize{12.5}{13}\selectfont\color{white}\MakeUppercase{#2}}\par\vspace{3pt}
    {\centering\popsemi\fontsize{7.8}{9.5}\selectfont\color{white}#3\par}
  \end{tcolorbox}}
\newcommand{\popcontenu}{%
  \par\noindent{\popxboldls\fontsize{7.5}{7.5}\selectfont\color{popmuted}CE COURS SIGNÉ CONTIENT}\par\vspace{7pt}
  \noindent\begin{minipage}[t]{0.235\linewidth}\poptile{popblue}{Cours}{complet, illustré, pas à pas}\end{minipage}\hfill
  \begin{minipage}[t]{0.235\linewidth}\poptile{poporange}{Méthodes}{et exercices résolus}\end{minipage}\hfill
  \begin{minipage}[t]{0.235\linewidth}\poptile{poppink}{Exercices}{type examen + corrigés}\end{minipage}\hfill
  \begin{minipage}[t]{0.235\linewidth}\poptile{popgreen}{Fiche bilan}{l'essentiel à retenir}\end{minipage}\par}

% Sommaire
\newtcolorbox{sommaire}{enhanced,colback=white,colframe=popink,boxrule=2.2pt,arc=10pt,
  drop shadow=popink,left=18pt,right=18pt,top=13pt,bottom=13pt,before skip=6pt,after skip=20pt,
  before upper={\poppill{Sommaire}{poppurple}{white}\par\vspace{9pt}\popsemi\fontsize{10.6}{14.5}\selectfont\color{popink}}}

% Signature de fin de cours — poussée tout en bas de la dernière page
\newcommand{\findecours}{%
  \par\vfill\nopagebreak
  \begin{center}\popsquiggle{poporange}\end{center}\vspace{4pt}
  \signatureonbuch
}
\AtBeginDocument{\def\llbracket{\mathopen{[\mkern-3.5mu[}}\def\rrbracket{\mathclose{]\mkern-3.5mu]}}} % intervalles d'entiers (glyphe ⟦ absent de la police)
% Glyphes absents de Fira Math : redéfinis à partir de glyphes disponibles
\AtBeginDocument{%
  \def\setminus{\mathbin{\backslash}}\let\smallsetminus\setminus
  \def\vdots{\mathord{\vbox{\baselineskip3.2pt\lineskiplimit0pt\kern4pt\hbox{.}\hbox{.}\hbox{.}}}}%
  \def\ddots{\mathinner{\mkern1mu\raise6.5pt\hbox{.}\mkern2mu\raise3.6pt\hbox{.}\mkern2mu\raise0.7pt\hbox{.}\mkern1mu}}%
  \def\triangle{\mathord{\scalebox{0.9}{$\symup{\Delta}$}}}%
  \def\nabla{\mathord{\raisebox{\depth}{\scalebox{1}[-1]{$\symup{\Delta}$}}}}%
  \def\checkmark{\popcheck{popdark}}%
  \def\star{\mathbin{\ast}}\def\bigstar{\mathord{\ast}}%
  \def\diamond{\mathbin{\scalebox{0.6}{$\Diamond$}}}%
  \def\bigcup{\mathop{\mathchoice{\raisebox{-0.3em}{\scalebox{1.7}{$\cup$}}}{\scalebox{1.25}{$\cup$}}{\cup}{\cup}}\displaylimits}%
  \def\bigcap{\mathop{\mathchoice{\raisebox{-0.3em}{\scalebox{1.7}{$\cap$}}}{\scalebox{1.25}{$\cap$}}{\cap}{\cap}}\displaylimits}%
  \def\lesseqqgtr{\mathrel{\lessgtr}}\def\bigoplus{\mathop{\oplus}\displaylimits}%
}
\providecommand{\ssi}{si et seulement si} % abréviation fréquente
\AtBeginDocument{\providecommand{\wideparen}{\overparen}} % arc de cercle

%% FIGFIX-BEGIN v5 — correctifs globaux des figures (docs/onbuch-generation/tools/figfix)
\usepackage{adjustbox}\usepackage{etoolbox}\tcbuselibrary{raster}
% toute figure TikZ/pgfplots plus large que la ligne est réduite à la largeur disponible
\BeforeBeginEnvironment{tikzpicture}{\begin{adjustbox}{max width=\linewidth}}
\AfterEndEnvironment{tikzpicture}{\end{adjustbox}}
% légendes pgfplots lisibles par-dessus les courbes
\pgfplotsset{every axis/.append style={legend style={fill=white,fill opacity=0.92,text opacity=1,draw=black!15,inner sep=3pt}}}
% étiquettes d'arêtes (syntaxe edge["..."]) avec fond blanc
\tikzset{every edge quotes/.append style={fill=white,inner sep=1.2pt}}
%% FIGFIX-END
\begin{document}

\pagegarde{Lesson 40 — Speaking: Exchanges information on participating in training activities/courses on democracy}{Anglais}{Tle A}{Chapter 4 : Citizenship and Human Rights}{EN40}{2 h}{Cette leçon d'expression orale apprend aux élèves de Terminale A à échanger des informations sur leur participation à des formations ou cours sur la démocratie. À travers des dialogues modèles, des formules fonctionnelles et des jeux de rôle guidés, les élèves apprennent à demander, donner et relancer une information avec une prononciation et une intonation correctes.
\vspace{4pt}
\begin{multicols}{2}\small
\begin{itemize}
\item À la fin de cette leçon, je sais utiliser le vocabulaire essentiel lié aux formations et cours sur la démocratie (training course, workshop, participant, certificate...).
\item À la fin de cette leçon, je sais demander des informations sur une formation (qui, quoi, où, quand, combien de temps, comment s'inscrire).
\item À la fin de cette leçon, je sais donner des informations sur ma participation à une formation de façon claire et organisée.
\item À la fin de cette leçon, je sais relancer la conversation et réagir aux propos de mon interlocuteur (Really? That sounds interesting! Tell me more...).
\item À la fin de cette leçon, je sais ouvrir et conclure poliment un dialogue en anglais.
\item À la fin de cette leçon, je sais prononcer correctement les mots-clés du thème avec la bonne accentuation et la bonne intonation (montante pour les questions yes/no, descendante pour les questions en wh-).
\end{itemize}\end{multicols}}

\begin{sommaire}
\begin{multicols}{2}
\begin{enumerate}
\item Mise en situation et vocabulaire du thème
\item Demander et donner une information
\item Relancer, réagir, ouvrir et conclure
\item Dialogues modèles commentés
\item Prononciation, accentuation et intonation
\item Jeux de rôle guidés et production orale
\item Activité d'intégration
\item Méthodes et exercices résolus
\item Exercices
\item Corrigés des exercices
\item Fiche bilan
\end{enumerate}
\end{multicols}
\end{sommaire}

\begin{objectifs}
\begin{itemize}
\item À la fin de cette leçon, je sais utiliser le vocabulaire essentiel lié aux formations et cours sur la démocratie (training course, workshop, participant, certificate...).
\item À la fin de cette leçon, je sais demander des informations sur une formation (qui, quoi, où, quand, combien de temps, comment s'inscrire).
\item À la fin de cette leçon, je sais donner des informations sur ma participation à une formation de façon claire et organisée.
\item À la fin de cette leçon, je sais relancer la conversation et réagir aux propos de mon interlocuteur (Really? That sounds interesting! Tell me more...).
\item À la fin de cette leçon, je sais ouvrir et conclure poliment un dialogue en anglais.
\item À la fin de cette leçon, je sais prononcer correctement les mots-clés du thème avec la bonne accentuation et la bonne intonation (montante pour les questions yes/no, descendante pour les questions en wh-).
\item À la fin de cette leçon, je sais mener un court dialogue complet et naturel sur la participation à une activité de formation civique.
\item À la fin de cette leçon, je sais adapter mon langage au contexte (dialogue entre amis ou entretien plus formel avec un organisateur).
\end{itemize}
\end{objectifs}

\begin{prerequis}
\begin{itemize}
\item Les questions en wh- et les questions yes/no (Première)
\item Le present perfect pour parler d'expériences (I have attended..., Première)
\item Le past simple pour raconter des événements terminés (Première)
\item Le vocabulaire de base de la citoyenneté et des droits de l'homme (Chapter 4, leçons précédentes)
\item Les formules de salutation et de politesse (Seconde/Première)
\end{itemize}
\end{prerequis}

\newpage

\coursec{Mise en situation et vocabulaire du thème}

\courssub{Brainstorming : qu'est-ce qu'une formation sur la démocratie ?}

Pendant les vacances, Amina, élève de Terminale à Yaoundé, a suivi une formation de trois jours sur la démocratie organisée par Elections Cameroon (ELECAM) et une ONG locale. De retour en classe, son camarade Tabi, qui vit à Bamenda, veut tout savoir : qui a organisé la formation, combien de temps elle a duré, ce qu'elle y a appris et comment s'inscrire à la prochaine session. Au Nigeria et au Ghana, des milliers de jeunes participent chaque année à des \eng{civic education workshops} et en parlent entre eux en anglais. Savez-vous raconter votre participation à une formation et poser les bonnes questions à quelqu'un qui y a assisté ? C'est exactement ce que cette leçon vous apprend à faire.

\textbf{Problématique.} Comment échanger, en anglais correct et naturel, des informations sur sa participation à une formation ou à un cours sur la démocratie ? Pour y répondre, il faut d'abord maîtriser le \cle{vocabulaire} du thème : c'est l'objet de cette première partie.

\textbf{Questions de brainstorming (à discuter en classe).}
\begin{itemize}
\item What is a training course on democracy? (Qu'est-ce qu'une formation sur la démocratie ?)
\item Who organizes such courses in Cameroon? Think of ELECAM, NGOs, the Ministry of Youth (MINJEC), schools and churches.
\item Have you ever attended a workshop or a seminar? What did you learn?
\end{itemize}

\begin{definition}[Training on democracy]
A \cle{training course on democracy} is an organized learning activity (course, workshop, seminar) in which people — often young citizens — learn about \eng{elections}, \eng{voting rights}, \eng{good governance} and \eng{civic responsibility}. In Cameroon, such activities are organized by \eng{ELECAM} (Elections Cameroon), \eng{NGOs} (non-governmental organizations), the \eng{MINJEC} (Ministry of Youth and Civic Education), schools and churches.
\end{definition}

\courssub{Texte support : Amina's experience}

Lisons d'abord le témoignage d'Amina. Observez les mots en gras : ils appartiennent au champ lexical de la leçon.

\begin{textebox}[Amina's testimony — “My training course on democracy”]
Last holidays, I attended a \textbf{three-day training course} on democracy organized by a local \textbf{NGO} in partnership with ELECAM. I \textbf{registered} online two weeks before the course started. The training \textbf{was held} at the youth centre in Yaoundé and it \textbf{lasted} three days, from Monday to Wednesday.

There were about forty \textbf{participants}, mostly students like me. Our \textbf{trainer}, Mrs Ngo Bell, was an experienced \textbf{facilitator}. She did not lecture us all day; instead, she \textbf{ran practical workshops} where we discussed real problems. A \textbf{guest speaker} from ELECAM explained how elections are organized in Cameroon.

We talked about \textbf{voting rights}, \textbf{good governance}, \textbf{tolerance} and \textbf{free and fair elections}. I learned that voting is not only a right but also a \textbf{civic responsibility}. We also prepared a small \textbf{awareness campaign} for our neighbourhood, to encourage young people to register on the electoral roll.

At the end of the course, every participant \textbf{received a certificate}. I am very proud of it. Next month, I will \textbf{take part in} a \textbf{seminar} on youth leadership in Douala. My friend Tabi wants to come too, so I told him how to register. Training courses like this one really \textbf{raise awareness} and help young people become active citizens.
\end{textebox}

\textbf{Glossaire.}
\begin{center}
\begin{tabularx}{\linewidth}{|l|l|X|}
\hline
\rowcolor{popblueL} Mot anglais & Nature & Traduction française \\
\hline
training course & nom & formation (stage) \\
\hline
to register & verbe & s'inscrire, s'enregistrer \\
\hline
to be held & verbe (passif) & se tenir, avoir lieu \\
\hline
to last & verbe & durer \\
\hline
facilitator & nom & animateur / animatrice \\
\hline
guest speaker & nom & intervenant invité \\
\hline
awareness campaign & nom & campagne de sensibilisation \\
\hline
electoral roll & nom & liste électorale \\
\hline
certificate & nom & attestation, certificat \\
\hline
\end{tabularx}
\end{center}

\textbf{Questions de compréhension.}
\begin{enumerate}
\item Who organized the training course? \emph{(A local NGO in partnership with ELECAM.)}
\item Where was the training held? \emph{(At the youth centre in Yaoundé.)}
\item How long did it last? \emph{(Three days.)}
\item What topics did the participants discuss? \emph{(Voting rights, good governance, tolerance, free and fair elections.)}
\item What did every participant receive at the end? \emph{(A certificate.)}
\end{enumerate}

\courssub{Champ lexical 1 : les types de formations}

\begin{definition}[Workshop]
A \cle{workshop} is a short, practical training session where participants actively discuss and practise (un atelier). Unlike a lecture, a workshop is interactive: people work in groups, ask questions and do activities.
\end{definition}

\begin{center}
\begin{tabularx}{\linewidth}{|l|X|X|}
\hline
\rowcolor{popblueL} English & Français & Remarque \\
\hline
training course & formation, stage & le terme général le plus courant \\
\hline
workshop & atelier & pratique, interactif, en petits groupes \\
\hline
seminar & séminaire & plus formel, souvent avec des experts \\
\hline
civic education programme & programme d'éducation civique & sur une longue période \\
\hline
awareness campaign & campagne de sensibilisation & destinée au grand public \\
\hline
conference & conférence & grande réunion avec des discours \\
\hline
\end{tabularx}
\end{center}

\begin{attention}[Faux ami : « formation »]
« Formation » ne se dit \textbf{pas} \eng{formation} en anglais dans ce sens ! On dit \eng{training course} ou simplement \eng{training}. En anglais, \eng{formation} signifie « forme, structure, formation (géologique, militaire) ».
\begin{itemize}
\item \eng{I attended a training course on democracy.} « J'ai suivi une formation sur la démocratie. »
\item \eng{The formation of the government took two weeks.} « La formation (constitution) du gouvernement a pris deux semaines. »
\end{itemize}
De même, « assister à » se dit \eng{to attend} et non \eng{to assist} (qui signifie « aider »).
\end{attention}

\courssub{Champ lexical 2 : les acteurs}

\begin{center}
\begin{tabularx}{\linewidth}{|l|X|X|}
\hline
\rowcolor{popblueL} English & Français & Exemple \\
\hline
participant & participant(e) & \eng{Forty participants attended the course.} \\
\hline
trainer / facilitator & formateur / animateur & \eng{The facilitator asked many questions.} \\
\hline
organizer & organisateur & \eng{The organizers welcomed us warmly.} \\
\hline
NGO (non-governmental organization) & ONG & \eng{A local NGO funded the workshop.} \\
\hline
youth leader & leader de jeunesse, responsable de jeunes & \eng{Our youth leader gave a speech.} \\
\hline
guest speaker & intervenant invité & \eng{The guest speaker came from ELECAM.} \\
\hline
\end{tabularx}
\end{center}

Notez la prononciation : \eng{participant} se prononce « par-ti-ssi-pent » (accent sur la deuxième syllabe) ; \eng{facilitator} « fe-si-li-té-teur » ; \eng{organizer} « or-ga-naï-zeur ».

\courssub{Champ lexical 3 : le déroulement de la formation}

\begin{propriete}[Les verbes-clés du déroulement]
\begin{itemize}
\item \cle{to attend} + nom (sans préposition !) : assister à. \eng{I attended a seminar.} « J'ai assisté à un séminaire. »
\item \cle{to register (for)} : s'inscrire (à). \eng{She registered for the workshop.} « Elle s'est inscrite à l'atelier. »
\item \cle{to take part in} : participer à, prendre part à. \eng{We took part in the discussions.} « Nous avons participé aux discussions. »
\item \cle{to last} + durée : durer. \eng{The course lasted three days.} « La formation a duré trois jours. »
\item \cle{to be held} (passif de \eng{to hold}, irrégulier : hold — held — held) : se tenir, avoir lieu. \eng{The workshop was held in Buea.} « L'atelier s'est tenu à Buéa. »
\item \cle{to receive a certificate} : recevoir une attestation. \eng{Every participant received a certificate.} « Chaque participant a reçu une attestation. »
\end{itemize}
\end{propriete}

\begin{exemplebox}[Exemples traduits]
\begin{itemize}
\item \eng{I registered for the seminar last week.} « Je me suis inscrit au séminaire la semaine dernière. »
\item \eng{The workshop was held in Buea.} « L'atelier s'est tenu à Buéa. »
\item \eng{The training lasted three days.} « La formation a duré trois jours. »
\item \eng{Did you attend the civic education programme?} « As-tu assisté au programme d'éducation civique ? »
\item \eng{She received her certificate at the end of the course.} « Elle a reçu son attestation à la fin de la formation. »
\end{itemize}
\end{exemplebox}

\begin{attention}[Prépositions et passif]
\begin{itemize}
\item On dit \eng{to attend a course} \textbf{sans} \eng{to} : ne dites jamais \eng{I attended to the course}.
\item \eng{To take part} exige \eng{in} : \eng{to take part \textbf{in} a workshop}.
\item « Se tenir » se traduit par le passif : \eng{The seminar \textbf{was held} in Douala} (et non \eng{was taken place} — \eng{to take place} est déjà intransitif : \eng{The seminar took place in Douala}).
\end{itemize}
\end{attention}

\courssub{Champ lexical 4 : le contenu démocratique}

\begin{center}
\begin{tabularx}{\linewidth}{|l|X|}
\hline
\rowcolor{popblueL} English & Français \\
\hline
elections & élections \\
\hline
voting rights & droits de vote \\
\hline
good governance & bonne gouvernance \\
\hline
tolerance & tolérance \\
\hline
civic responsibility & responsabilité civique \\
\hline
free and fair elections & élections libres et équitables \\
\hline
peaceful elections & élections apaisées / pacifiques \\
\hline
democracy & démocratie \\
\hline
citizen / citizenship & citoyen / citoyenneté \\
\hline
\end{tabularx}
\end{center}

Prononciation : \eng{democracy} « dé-mo-kre-ssi » (accent sur la deuxième syllabe) ; \eng{governance} « go-veur-nens » ; \eng{fair} « fèr » (rime avec \eng{hair}).

\courssub{Collocations utiles}

Une \cle{collocation} est une association habituelle de mots. Les apprendre par blocs vous fera parler plus naturellement.

\begin{center}
\begin{tabularx}{\linewidth}{|X|X|}
\hline
\rowcolor{popblueL} Collocation anglaise & Traduction \\
\hline
attend a training course & suivre une formation \\
\hline
run a workshop & animer un atelier \\
\hline
gain knowledge & acquérir des connaissances \\
\hline
raise awareness & sensibiliser (litt. « élever la conscience ») \\
\hline
issue a certificate & délivrer une attestation \\
\hline
organize a seminar & organiser un séminaire \\
\hline
register for a course & s'inscrire à une formation \\
\hline
\end{tabularx}
\end{center}

\begin{exemplebox}[Les collocations en contexte]
\begin{itemize}
\item \eng{The NGO runs workshops in many schools.} « L'ONG anime des ateliers dans de nombreuses écoles. »
\item \eng{I gained a lot of knowledge about elections.} « J'ai acquis beaucoup de connaissances sur les élections. »
\item \eng{The campaign raised awareness among young voters.} « La campagne a sensibilisé les jeunes électeurs. »
\item \eng{ELECAM issued certificates to all participants.} « ELECAM a délivré des attestations à tous les participants. »
\end{itemize}
\end{exemplebox}

\courssub{Carte mentale du vocabulaire}

\begin{popfigure}
\begin{tikzpicture}[mindmap, grow cyclic, every node/.style={concept, text=white}, root concept/.append style={concept color=popblue, font=\small\bfseries}, level 1/.append style={level distance=3.6cm, sibling angle=90}, level 2/.append style={level distance=2.4cm, sibling angle=45, font=\tiny}, every child node/.style={concept}]
\node[root concept]{Training on democracy}
  child[concept color=poporange]{ node {Types}
    child { node {workshop} }
    child { node {seminar} }
    child { node {awareness campaign} }
  }
  child[concept color=popgreen]{ node {People}
    child { node {participant} }
    child { node {trainer} }
    child { node {guest speaker} }
  }
  child[concept color=poppurple]{ node {Actions}
    child { node {attend} }
    child { node {register} }
    child { node {receive a certificate} }
  }
  child[concept color=poppink]{ node {Topics}
    child { node {voting rights} }
    child { node {good governance} }
    child { node {free and fair elections} }
  };
\end{tikzpicture}
\legende{Carte mentale du champ lexical : quatre branches pour organiser le vocabulaire de la formation sur la démocratie.}
\end{popfigure}

\begin{methode}[Comment mémoriser ce vocabulaire]
\begin{enumerate}
\item Classez chaque nouveau mot dans l'une des quatre branches de la carte mentale (Types, People, Actions, Topics).
\item Apprenez les mots par \cle{collocations} (blocs), jamais isolément : \eng{attend a course}, pas seulement \eng{attend}.
\item Fabriquez une phrase vraie sur votre propre vie avec chaque collocation : \eng{I attended a workshop in my school last year.}
\item Testez-vous en racontant la formation d'Amina en cinq phrases, sans regarder le texte.
\end{enumerate}
\end{methode}

\begin{exoresolu}[Complétez avec le mot qui convient]
Énoncé. Complétez les phrases avec : \eng{held}, \eng{registered}, \eng{lasted}, \eng{certificate}, \eng{raise awareness}.
\begin{enumerate}
\item The seminar \ldots\ in Bamenda last month.
\item I \ldots\ for the workshop online.
\item The training course \ldots\ two days.
\item Every participant received a \ldots .
\item The campaign aims to \ldots\ about voting rights.
\end{enumerate}
\tcblower
Solution détaillée.
\begin{enumerate}
\item \eng{was held} — « se tenir » exige le passif de \eng{hold} (hold — held — held) : « Le séminaire s'est tenu à Bamenda le mois dernier. »
\item \eng{registered} — \eng{to register for} = s'inscrire à ; prétérit régulier : « Je me suis inscrit à l'atelier en ligne. »
\item \eng{lasted} — \eng{to last} + durée = durer : « La formation a duré deux jours. »
\item \eng{certificate} — collocation \eng{receive a certificate} : « Chaque participant a reçu une attestation. »
\item \eng{raise awareness} — collocation : « La campagne vise à sensibiliser aux droits de vote. »
\end{enumerate}
\end{exoresolu}

\begin{exercice}[À vous de jouer]
Traduisez en anglais :
\begin{enumerate}
\item J'ai assisté à un atelier sur la bonne gouvernance.
\item La formation a duré trois jours et s'est tenue à Douala.
\item L'ONG a délivré des attestations à tous les participants.
\item Nous voulons sensibiliser les jeunes aux élections libres et équitables.
\end{enumerate}
\end{exercice}

\begin{corrige}[Exercice « À vous de jouer »]
\begin{enumerate}
\item \eng{I attended a workshop on good governance.}
\item \eng{The training (course) lasted three days and was held in Douala.}
\item \eng{The NGO issued certificates to all the participants.}
\item \eng{We want to raise awareness among young people about free and fair elections.}
\end{enumerate}
\end{corrige}

\begin{savaistu}[Civic education in Ghana]
Au Ghana, le \eng{National Commission for Civic Education} (NCCE) organise chaque année des milliers de \eng{workshops} pour les jeunes électeurs, dans les écoles, les églises et les marchés. L'objectif : expliquer le droit de vote, encourager la tolérance et promouvoir des élections pacifiques. Au Cameroun, ELECAM et le MINJEC mènent des actions comparables d'\eng{civic education} auprès des jeunes.
\end{savaistu}

\begin{aretenir}[L'essentiel de la section]
\begin{itemize}
\item « Formation » = \eng{training course} (jamais \eng{formation} !).
\item Types : \eng{workshop} (atelier pratique), \eng{seminar}, \eng{awareness campaign}, \eng{civic education programme}.
\item Acteurs : \eng{participant}, \eng{trainer/facilitator}, \eng{organizer}, \eng{NGO}, \eng{guest speaker}.
\item Déroulement : \eng{attend} (sans \eng{to}), \eng{register for}, \eng{take part in}, \eng{last three days}, \eng{be held in Yaoundé}, \eng{receive a certificate}.
\item Contenu : \eng{voting rights}, \eng{good governance}, \eng{tolerance}, \eng{free and fair elections}.
\end{itemize}
\end{aretenir}

\coursec{Demander et donner une information}

Dans la section précédente, nous avons découvert le vocabulaire essentiel des formations sur la démocratie : \eng{training course}, \eng{workshop}, \eng{seminar}, \eng{participant}, \eng{facilitator}, \eng{certificate}. Mais connaître les mots ne suffit pas : pour \emph{échanger} des informations, il faut savoir \cle{poser les bonnes questions} et \cle{donner des réponses complètes et naturelles}. C'est l'objet de cette section : apprendre à interroger un camarade sur une formation à laquelle il a participé, et à raconter sa propre expérience de façon organisée.

\courssub{Observer avant de comprendre : un extrait de dialogue}

Lisons d'abord ce court échange entre deux élèves de Terminale, à Yaoundé, après les vacances.

\begin{textebox}[At the school gate — after the holidays]
Amina: Hi, Boris! I heard you attended a training course during the holidays. What was it about?

Boris: It was a course on democracy and elections. It was organised by ELECAM and a local NGO.

Amina: Really? Where did it take place?

Boris: At the community hall in Mvog-Ada, here in Yaoundé.

Amina: And how long did it last?

Boris: It was a three-day course, from the 12th to the 14th of August. About fifty participants attended.

Amina: Did you enjoy it?

Boris: Yes, I did! It was mostly group work and debates. In fact, I have never learnt so much in three days. And I received a certificate at the end.

Amina: That's great! Have you attended other courses like this one?

Boris: Yes, I have attended two workshops so far, but this one was the most useful.
\end{textebox}

\textbf{Glossaire}

\begin{center}
\begin{tabularx}{\linewidth}{|l|l|X|}
\hline
\rowcolor{popblueL} English & Nature & Français \\
\hline
to attend & verbe & assister à, suivre (une formation) \\
\hline
to take place & locution verbale & avoir lieu, se dérouler \\
\hline
community hall & nom & salle communautaire \\
\hline
to last & verbe & durer \\
\hline
group work & nom & travail de groupe \\
\hline
so far & locution adverbiale & jusqu'à présent \\
\hline
useful & adjectif & utile \\
\hline
\end{tabularx}
\end{center}

\textbf{Observons.} Quelles questions Amina pose-t-elle ? \eng{What was it about? Where did it take place? How long did it last? Did you enjoy it? Have you attended other courses?} On remarque deux grandes familles de questions : celles qui commencent par un \cle{mot interrogatif en wh-} (qui demandent une information précise) et celles qui commencent par un \cle{auxiliaire} (qui attendent \eng{yes} ou \eng{no}). Observons aussi les réponses de Boris : il dit \eng{I have attended two workshops} (expérience générale, sans date) mais \eng{It was a three-day course, from the 12th to the 14th of August} (détail daté). Ce n'est pas un hasard : c'est une règle fondamentale de l'anglais.

\courssub{Les questions en wh- : obtenir une information précise}

\begin{propriete}[Former une question en wh-]
La structure d'une question en wh- au passé est :

\textbf{Mot en wh- + auxiliaire (did / was / were) + sujet + verbe + complément ?}

\begin{itemize}
\item Avec un verbe ordinaire : \eng{What \textbf{did} the training \textbf{cover}?} — Que couvrait la formation ?
\item Avec \eng{be} (pas de \eng{did}) : \eng{Where \textbf{was} it \textbf{held}?} — Où s'est-elle tenue ?
\item Quand le mot en wh- est le \emph{sujet} du verbe, il n'y a pas d'auxiliaire : \eng{Who \textbf{organised} it?} — Qui l'a organisée ? (et non \emph{Who did organise it?})
\end{itemize}
\end{propriete}

Voici les questions en wh- indispensables pour s'informer sur une formation :

\begin{center}
\begin{tabularx}{\linewidth}{|X|X|}
\hline
\rowcolor{popblueL} Question en anglais & Traduction française \\
\hline
What was the training about? & De quoi traitait la formation ? \\
\hline
Where was it held? / Where did it take place? & Où s'est-elle déroulée ? \\
\hline
When did it take place? & Quand a-t-elle eu lieu ? \\
\hline
How long did it last? & Combien de temps a-t-elle duré ? \\
\hline
Who organised it? & Qui l'a organisée ? \\
\hline
How many people attended? & Combien de personnes y ont assisté ? \\
\hline
What did you learn? & Qu'as-tu appris ? \\
\hline
How much did it cost? & Combien a-t-elle coûté ? \\
\hline
\end{tabularx}
\end{center}

\begin{attention}[How long, pas « how many time »]
Ne dis jamais \emph{How many time did it last?} mais \eng{How long did it last?} — Combien de temps cela a-t-il duré ? Le mot \eng{time} au sens de « durée » est \cle{indénombrable} en anglais : on ne peut pas le compter. \eng{How many} ne s'emploie qu'avec des noms \emph{comptables au pluriel} : \eng{How many participants attended?} — Combien de participants ont assisté ? Compare : \eng{How many days did it last?} (correct, car \eng{days} est comptable) et \eng{How long did it last?} (correct, pour la durée).
\end{attention}

\courssub{Les questions yes/no et les réponses courtes}

\begin{propriete}[Questions fermées et réponses courtes]
Une question \cle{yes/no} commence par un auxiliaire : \eng{did}, \eng{was/were}, \eng{have/has}. La réponse courte reprend \textbf{le même auxiliaire}, jamais le verbe principal seul :

\begin{itemize}
\item \eng{Did you enjoy it? — Yes, I did. / No, I didn't.} — Tu as aimé ? — Oui. / Non.
\item \eng{Was it useful? — Yes, it was.} — C'était utile ? — Oui.
\item \eng{Did you receive a certificate? — No, I didn't, unfortunately.} — As-tu reçu un certificat ? — Non, malheureusement.
\item \eng{Have you ever attended such a course? — No, not yet.} — As-tu déjà suivi une telle formation ? — Pas encore.
\end{itemize}

En français on répond simplement « oui » ou « non » ; en anglais, la réponse courte avec l'auxiliaire (\eng{Yes, I did.}) est beaucoup plus naturelle et polie à l'oral.
\end{propriete}

\begin{exemplebox}[Réponses courtes naturelles]
\eng{Did you learn a lot? — Yes, I did.} — As-tu beaucoup appris ? — Oui.

\eng{Have you got your certificate yet? — No, not yet.} — As-tu déjà ton certificat ? — Pas encore.

\eng{How many people attended? — About fifty participants.} — Combien de personnes ont participé ? — Environ cinquante participants.

\eng{Was the course free? — Yes, it was, actually.} — La formation était-elle gratuite ? — Oui, en fait.
\end{exemplebox}

\courssub{Present perfect ou past simple ? Le choix décisif}

C'est le cœur grammatical de cette leçon. Observons à nouveau les phrases de Boris : \eng{I have attended two workshops so far} puis \eng{It was a three-day course, from the 12th to the 14th of August}.

\begin{propriete}[Expérience sans date = present perfect ; détail daté = past simple]
Pour parler d'une \cle{expérience} sans donner de date précise, on utilise le \cle{present perfect} : \eng{I have taken part in a seminar.} — J'ai participé à un séminaire (un jour, dans ma vie, on ne dit pas quand).

Dès qu'on donne un \cle{détail daté} (quand, où, combien de temps), on repasse obligatoirement au \cle{past simple} : \eng{It took place in July. It lasted three days.} — Elle a eu lieu en juillet. Elle a duré trois jours.

\textbf{Forme du present perfect :} sujet + \eng{have/has} + participe passé. \textbf{Forme du past simple :} sujet + verbe au passé (\eng{attended}, \eng{was}, \eng{took}, \eng{lasted}).

Marqueurs typiques du present perfect d'expérience : \eng{ever} (déjà, dans une question), \eng{never} (jamais), \eng{so far} (jusqu'à présent), \eng{twice, three times} (deux, trois fois). Marqueurs typiques du past simple : \eng{yesterday}, \eng{last year}, \eng{in July}, \eng{two weeks ago}, une date précise.
\end{propriete}

\begin{center}
\begin{tabularx}{\linewidth}{|X|X|X|}
\hline
\rowcolor{popblueL} Present perfect (expérience) & Past simple (détail daté) & Remarque \\
\hline
I have attended two workshops. & I attended a workshop in June. & Dès qu'une date apparaît, past simple. \\
\hline
Have you ever been to Buea? & I went there last year. & \eng{ever} = expérience ; \eng{last year} = date. \\
\hline
She has never received a certificate. & She received one in 2023. & \eng{never} = expérience ; \eng{in 2023} = date. \\
\hline
We have had three training sessions so far. & The last one ended on Friday. & \eng{so far} = bilan ; \eng{on Friday} = date. \\
\hline
\end{tabularx}
\end{center}

\begin{attention}[Le piège n° 1 des francophones]
En français, on dit « j'ai assisté à un séminaire en juillet » avec un passé composé. En anglais, c'est \textbf{impossible} : une date passée et terminée interdit le present perfect. On ne dit jamais \emph{I have attended a seminar in July}, mais \eng{I attended a seminar in July.} Règle pratique : \textbf{date précise = past simple, toujours}.
\end{attention}

\courssub{Donner une information complète et organisée}

Quand on te demande de raconter une formation, une réponse d'un seul mot (\eng{It was good.}) est pauvre. Un bon locuteur organise son information.

\begin{methode}[L'ordre WHAT → WHERE/WHEN → CONTENT → OPINION]
Pour donner une information complète sur une formation, suis toujours ces quatre étapes :

\begin{enumerate}
\item \textbf{WHAT} — annonce le sujet : \eng{It was a training course on democracy and citizenship.} (la formation)
\item \textbf{WHERE/WHEN} — situe dans l'espace et le temps : \eng{It took place at City Hall in Douala, last month.} (le lieu et la date)
\item \textbf{CONTENT} — décris le contenu : \eng{We discussed elections, voting rights and the role of young citizens. We worked in groups and did role plays.} (le contenu)
\item \textbf{OPINION} — donne ton appréciation personnelle : \eng{In my opinion, it was very useful. I really enjoyed the debates.} (ton avis)
\end{enumerate}

Ce plan en quatre points te garantit une réponse riche, claire et naturelle — exactement ce qu'attend l'examinateur à l'oral.
\end{methode}

\begin{exemplebox}[Une réponse modèle complète]
\eng{During the last holidays, I attended a three-day training on democracy. It was held at the youth centre in Bamenda, in August. We mainly studied the electoral process and human rights, and we did a lot of group work. Actually, it was one of the most useful courses I have ever attended.}

« Pendant les dernières vacances, j'ai suivi une formation de trois jours sur la démocratie. Elle s'est tenue au centre des jeunes de Bamenda, en août. Nous avons principalement étudié le processus électoral et les droits de l'homme, et nous avons fait beaucoup de travaux de groupe. En fait, c'était l'une des formations les plus utiles auxquelles j'aie jamais assisté. »

Observe : WHAT → WHERE/WHEN → CONTENT → OPINION, et le passage du past simple (détails) au present perfect (\eng{I have ever attended}, expérience).
\end{exemplebox}

\begin{attention}[« a three-day course », sans s à day]
En anglais, quand un nombre et un nom forment un adjectif composé devant un autre nom, le nom reste au \cle{singulier} et on met des traits d'union : \eng{a three-day course} — une formation de trois jours ; \eng{a two-hour workshop} — un atelier de deux heures ; \eng{a five-page report} — un rapport de cinq pages. On ne dit jamais \emph{a three-days training}. Le français dit « de trois jours » avec un pluriel : ne traduis pas mot à mot !
\end{attention}

\courssub{Nuancer ses propos : mainly, mostly, in fact, actually, as far as I know}

Pour sonner naturel et précis, un bon locuteur \cle{nuance} ses informations au lieu de tout affirmer de façon absolue.

\begin{center}
\begin{tabularx}{\linewidth}{|l|X|X|}
\hline
\rowcolor{popblueL} Expression & Sens & Exemple traduit \\
\hline
mainly & principalement & \eng{The course was mainly about voting rights.} — La formation portait principalement sur le droit de vote. \\
\hline
mostly & surtout, en majorité & \eng{The participants were mostly students.} — Les participants étaient surtout des étudiants. \\
\hline
in fact & en fait, d'ailleurs & \eng{In fact, it was free of charge.} — En fait, c'était gratuit. \\
\hline
actually & en réalité, en fait & \eng{It was actually quite interesting.} — C'était en fait assez intéressant. \\
\hline
as far as I know & pour autant que je sache & \eng{As far as I know, about fifty people attended.} — Pour autant que je sache, environ cinquante personnes ont assisté. \\
\hline
\end{tabularx}
\end{center}

\begin{attention}[Faux ami : actually ne veut pas dire « actuellement »]
\eng{Actually} signifie « en fait, en réalité » : \eng{I actually enjoyed it.} — En fait, j'ai aimé. Pour dire « actuellement », utilise \eng{currently} ou \eng{at the moment} : \eng{I am currently attending a course.} — Je suis actuellement une formation. De même, \eng{in fact} renforce ou corrige une information ; il ne s'emploie pas comme « en fait » au sens de « en pratique » dans tous les contextes.
\end{attention}

\courssub{Figure récapitulative : demander et répondre}

\begin{popfigure}
\begin{tikzpicture}
\node[draw=popblue, fill=popblueL, rounded corners, thick, minimum width=6.8cm, text width=6.4cm, align=left, anchor=north west] (ask) at (0,0) {
\textbf{\textcolor{popdark}{ASKING — poser les questions}}\\[2pt]
\eng{What was the training about?}\\
\eng{Where was it held?}\\
\eng{When did it take place?}\\
\eng{How long did it last?}\\
\eng{Who organised it?}\\
\eng{How many people attended?}\\
\eng{Did you enjoy it?}\\
\eng{Was it useful?}\\
\eng{Have you ever attended such a course?}
};
\node[draw=popgreen, fill=popgreenL, rounded corners, thick, minimum width=6.8cm, text width=6.4cm, align=left, anchor=north west] (ans) at (7.6,0) {
\textbf{\textcolor{popdark}{ANSWERING — donner l'information}}\\[2pt]
\eng{It was about democracy and elections.}\\
\eng{It took place at the youth centre.}\\
\eng{It was held in August, last year.}\\
\eng{It was a three-day course.}\\
\eng{It was organised by ELECAM.}\\
\eng{About fifty participants.}\\
\eng{Yes, I did. It was mostly group work.}\\
\eng{Yes, it was very useful, actually.}\\
\eng{Yes, I have attended two so far.}
};
\draw[-{Stealth}, very thick, poporange] (ask.east) -- (ans.west) node[midway, above, text=popdark, font=\small\bfseries] {échange};
\end{tikzpicture}
\legende{Les deux colonnes de l'échange : à chaque question (colonne de gauche) correspond une réponse naturelle et organisée (colonne de droite).}
\end{popfigure}

\courssub{Exercice résolu et entraînement}

\begin{exoresolu}[Compléter un dialogue]
Complète ce dialogue avec la question ou la réponse qui convient.

1. A : \ldots\ldots\ldots\ldots\ldots\ldots\ldots ? — B : It was about human rights and democracy.

2. A : How long did the course last? — B : \ldots\ldots\ldots\ldots\ldots\ldots\ldots (three days)

3. A : Did you receive a certificate? — B : \ldots\ldots\ldots\ldots\ldots\ldots\ldots (réponse courte affirmative)

4. A : Have you ever attended a workshop in Buea? — B : \ldots\ldots\ldots\ldots\ldots\ldots\ldots (réponse courte négative avec « pas encore »)

\tcblower
\textbf{Solution détaillée}

1. La réponse donne le \emph{sujet} de la formation : on demande donc « de quoi traitait la formation ? » → \eng{What was the training about?} (ou \eng{What was it about?}). Attention à l'ordre des mots : \eng{was} avant le sujet, sans \eng{did} car le verbe est \eng{be}.

2. On demande la durée : réponse au past simple (détail daté) → \eng{It lasted three days.} On peut aussi dire \eng{It was a three-day course.} — sans « s » à \eng{day}.

3. Réponse courte affirmative avec le même auxiliaire \eng{did} → \eng{Yes, I did.} On ne répond jamais \emph{Yes, I received} seul.

4. Réponse courte négative au present perfect, avec \eng{not yet} → \eng{No, not yet.} (ou \eng{No, I haven't.}). \eng{Yet} est le marqueur typique du present perfect dans les phrases négatives et les questions.
\end{exoresolu}

\begin{exercice}[À toi de jouer]
1. Mets les mots dans l'ordre pour former des questions : (a) about / what / training / the / was / ? (b) did / how long / last / it / ? (c) organised / who / it / ?

2. Choisis le bon temps : I (attend) \ldots\ a seminar on citizenship last March. I (never receive) \ldots\ a certificate before, so I was very proud.

3. Corrige les erreurs : (a) \emph{How many time did the workshop last?} (b) \emph{It was a five-days training.} (c) \emph{I have attended this course in 2023.}
\end{exercice}

\begin{corrige}[Exercice « À toi de jouer »]
1. (a) \eng{What was the training about?} (b) \eng{How long did it last?} (c) \eng{Who organised it?} — pas d'auxiliaire car \eng{who} est le sujet.

2. \eng{I \textbf{attended} a seminar on citizenship last March. I \textbf{have never received} a certificate before, so I was very proud.} — \eng{last March} impose le past simple ; \eng{never... before} (expérience de vie) impose le present perfect.

3. (a) \eng{How long did the workshop last?} (b) \eng{It was a five-day training.} (c) \eng{I attended this course in 2023.} — une date passée interdit le present perfect.
\end{corrige}

\begin{aretenir}[L'essentiel de la section]
\begin{itemize}
\item Questions de précision : \eng{What / Where / When / How long / Who / How many} + auxiliaire + sujet + verbe.
\item Questions fermées : auxiliaire en tête ; réponse courte avec le même auxiliaire : \eng{Yes, I did. / No, not yet.}
\item Expérience sans date → \cle{present perfect} (\eng{I have attended two workshops}) ; détail daté → \cle{past simple} (\eng{It lasted three days}).
\item Réponse complète = WHAT → WHERE/WHEN → CONTENT → OPINION.
\item Nuancer avec \eng{mainly, mostly, in fact, actually, as far as I know}.
\item Pièges : \eng{How long} (pas « how many time »), \eng{a three-day course} (sans « s »), jamais de present perfect avec une date passée.
\end{itemize}
\end{aretenir}

Tu sais désormais poser des questions précises et donner des réponses complètes et nuancées. Dans la section suivante, nous apprendrons à \emph{relancer} la conversation, à \emph{réagir} aux informations de ton partenaire et à \emph{ouvrir et conclure} poliment un échange — les gestes qui transforment une série de questions en une véritable conversation.

\coursec{Relancer, réagir, ouvrir et conclure}

Dans la section précédente, nous avons appris à \emph{demander} une information (\eng{Where was the training held?}) et à la \emph{donner} (\eng{It took place at the City Hall in Yaoundé.}). Mais un dialogue ne se résume pas à une suite de questions et de réponses : pour qu'un échange soit vivant et naturel, il faut savoir \cle{ouvrir} la conversation, \cle{réagir} à ce que dit le partenaire, \cle{relancer} quand le rythme retombe, et \cle{conclure} poliment. C'est exactement ce que nous allons apprendre ici, avec les formules toutes prêtes utilisées par les anglophones dans la vie quotidienne.

\courssub{Observer : un dialogue complet}

\begin{textebox}[At the school canteen — a dialogue between two friends]
\textbf{Martha:} Guess what! I attended a training course on democracy last month.\\
\textbf{Eyong:} Really? That sounds interesting! Tell me more.\\
\textbf{Martha:} Well, we learned about voting procedures and citizens' rights. The trainers were very friendly.\\
\textbf{Eyong:} Wow! And what happened next?\\
\textbf{Martha:} We organised a mock election in groups. It was great fun.\\
\textbf{Eyong:} Sorry, what do you mean by “mock election”?\\
\textbf{Martha:} A practice election, not a real one. We voted for class representatives.\\
\textbf{Eyong:} I see. And did you get a certificate?\\
\textbf{Martha:} Yes, I did! In my opinion, it was very useful.\\
\textbf{Eyong:} I wish I could attend the next one. We don't have such courses in my school.\\
\textbf{Martha:} Don't worry, there is another session in December.\\
\textbf{Eyong:} Thanks for telling me about it. I'll try to register for the next session.\\
\textbf{Martha:} Great! Let's keep in touch about it.
\end{textebox}

\textbf{Glossaire :} \emph{training course} (nom) : stage de formation ; \emph{voting procedures} (nom) : procédures de vote ; \emph{citizens' rights} (nom) : droits des citoyens ; \emph{mock election} (nom) : élection simulée ; \emph{certificate} (nom) : attestation ; \emph{to register} (verbe) : s'inscrire ; \emph{session} (nom) : session.

Observe bien la structure du dialogue : Martha \textbf{ouvre} avec \eng{Guess what!}, Eyong \textbf{réagit} (\eng{Really? That sounds interesting!}), \textbf{relance} (\eng{Tell me more. / What happened next?}), demande une \textbf{précision} (\eng{what do you mean by...?}), donne son \textbf{avis} (\eng{I wish I could attend...}), puis les deux amis \textbf{concluent} poliment. C'est le squelette de tout échange oral réussi.

\courssub{Ouvrir le dialogue}

Pour annoncer une nouvelle ou lancer un sujet, on utilise des formules d'\cle{ouverture} qui attirent l'attention du partenaire.

\begin{propriete}[Formules d'ouverture]
\begin{itemize}
\item \eng{Guess what!} — « Devine quoi ! » : pour annoncer une nouvelle surprenante (registre familier).
\item \eng{Have you heard about the civic education workshop?} — « As-tu entendu parler de l'atelier d'éducation civique ? » : pour introduire un sujet.
\item \eng{You know what? I joined a democracy club at school.} — « Tu sais quoi ? Je me suis inscrit au club de démocratie. »
\item \eng{By the way, did you attend the training last week?} — « Au fait, as-tu participé à la formation la semaine dernière ? »
\end{itemize}
\end{propriete}

\begin{attention}[Piège des francophones]
Ne traduis pas « Devine quoi ! » par \eng{Guess what thing!} : la formule figée est \eng{Guess what!}, sans \emph{thing}. De même, \eng{Have you heard about...?} exige la préposition \cle{about}. Ne confonds pas avec \eng{Have you heard of...?}, qui signifie « connais-tu, as-tu déjà entendu parler de » une personne ou une chose : \eng{Have you heard of this trainer?} (« Connais-tu ce formateur ? »).
\end{attention}

\courssub{Réagir et montrer son intérêt}

\begin{propriete}[Réactions courtes]
Après une information, l'anglophone réagit presque toujours avec une courte exclamation :
\begin{itemize}
\item \eng{Really?} — « Vraiment ? »
\item \eng{That's great!} — « C'est formidable ! »
\item \eng{Wow, that sounds interesting!} — « Ouah, ça a l'air intéressant ! »
\item \eng{How wonderful!} — « Comme c'est merveilleux ! »
\item \eng{I see.} — « Je vois. » (réaction neutre, pour accuser réception)
\item \eng{Uh-huh? / Oh yes?} — « Ah oui ? » (petites réactions d'écoute)
\end{itemize}
La réaction interrogative (\eng{Really?}) se dit avec une \textbf{intonation montante} ; la réaction admirative (\eng{That's great!}) avec une intonation descendante et enthousiaste.
\end{propriete}

\begin{savaistu}[La politesse de l'écoute]
Dans la culture anglophone, montrer son intérêt par de petites réactions (\eng{Really? Uh-huh? I see.}) est essentiel : un silence prolongé peut sembler impoli ou faire croire que tu n'écoutes pas. Ces petites « réponses d'écoute » (les \emph{back-channels}) ponctuent toutes les conversations, en face à face comme au téléphone. Entraîne-toi à les placer naturellement !
\end{savaistu}

\courssub{Relancer la conversation}

\begin{propriete}[Formules de relance]
Pour prolonger l'échange après une réponse :
\begin{itemize}
\item \eng{Tell me more about it.} — « Raconte-m'en davantage. »
\item \eng{What happened next?} — « Que s'est-il passé ensuite ? »
\item \eng{And then?} — « Et ensuite ? »
\item \eng{What about the trainers?} — « Et les formateurs ? »
\item \eng{What was it like?} — « C'était comment ? »
\end{itemize}
\end{propriete}

\begin{methode}[La technique « réaction + question »]
Pour relancer efficacement, procède en deux temps après chaque réponse de ton partenaire :
\begin{enumerate}
\item Utilise d'abord une \textbf{réaction courte} : \eng{Really?}, \eng{I see.}, \eng{That's great!}
\item Enchaîne avec une \textbf{nouvelle question en wh-} (\emph{what, where, when, who, how}) pour prolonger l'échange.
\end{enumerate}
Exemple : \eng{— We practised public speaking. — Really? How did it go?} (« — Nous avons pratiqué la prise de parole en public. — Vraiment ? Comment cela s'est-il passé ? ») Cette alternance réaction/question est le moteur d'un dialogue naturel.
\end{methode}

\begin{attention}[Tell me more, pas Continue !]
\eng{Tell me more} est bien plus naturel que \eng{Continue} ou \eng{Explain} entre amis : ces derniers sonnent comme des ordres de professeur ! Réserve \eng{Could you explain...?} aux contextes formels (avec un formateur, un fonctionnaire). De même, évite \eng{What else?} trop sec ; préfère \eng{What else did you learn?} (« Qu'as-tu appris d'autre ? »).
\end{attention}

\courssub{Demander une précision}

Quand tu ne comprends pas un mot ou une phrase, ne reste pas silencieux : demande poliment une explication.

\begin{exemplebox}[Demander une précision ou une répétition]
\eng{Sorry, what do you mean by “facilitator”?} — « Pardon, qu'entends-tu par “facilitator” (animateur) ? »\\
\eng{Could you repeat that, please?} — « Pourriez-vous répéter, s'il vous plaît ? »\\
\eng{What does “mock election” mean?} — « Que signifie “mock election” ? »\\
\eng{Sorry, I didn't catch the date.} — « Pardon, je n'ai pas saisi la date. »
\end{exemplebox}

Note la structure : \eng{What do you mean by + mot ?} (« que veux-tu dire par... ») et \eng{What does + mot + mean?} (« que signifie... »). Ne confonds pas les deux constructions !

\courssub{Donner son avis et comparer}

\begin{exemplebox}[Opinion et comparaison]
\eng{In my opinion, it was very useful.} — « À mon avis, c'était très utile. »\\
\eng{I wish I could attend the next one.} — « J'aimerais pouvoir assister à la prochaine session. »\\
\eng{We don't have such courses in my school.} — « Nous n'avons pas de tels cours dans mon école. »\\
\eng{I think every student should attend it.} — « Je pense que chaque élève devrait y participer. »
\end{exemplebox}

\begin{propriete}[La structure “I wish + could”]
\eng{I wish I could + base verbale} exprime un souhait irréel au présent : « j'aimerais pouvoir... ». Après \cle{wish}, on utilise le prétérit (ici \emph{could}, prétérit de \emph{can}) : \eng{I wish I could attend} (et non \eng{I wish I can attend}). Compare : \emph{je veux} = \eng{I want to attend} ; \emph{j'aimerais (mais je ne peux pas)} = \eng{I wish I could attend}.
\end{propriete}

\courssub{Conclure poliment}

\begin{propriete}[Formules de conclusion]
\begin{itemize}
\item \eng{Thanks for telling me about it.} — « Merci de m'en avoir parlé. »
\item \eng{I'll try to register for the next session.} — « Je vais essayer de m'inscrire à la prochaine session. »
\item \eng{Let's keep in touch about it.} — « Restons en contact à ce sujet. »
\item \eng{It was nice talking to you.} — « C'était un plaisir de discuter avec toi. »
\end{itemize}
Attention à la construction : \eng{Thanks for + -ing} (\eng{Thanks for telling me}), jamais \eng{Thanks for to tell me}.
\end{propriete}

\courssub{Tableau récapitulatif des fonctions}

\begin{center}
\begin{tabularx}{\linewidth}{|l|X|X|}
\hline
\rowcolor{popblueL} \textbf{Fonction} & \textbf{English} & \textbf{Français} \\
\hline
Ouvrir & Guess what! / Have you heard about...? & Devine quoi ! / As-tu entendu parler de... ? \\
\hline
Réagir & Really? / That's great! / I see. & Vraiment ? / C'est génial ! / Je vois. \\
\hline
Relancer & Tell me more. / What happened next? / And then? & Raconte-m'en plus. / Et ensuite ? \\
\hline
Préciser & What do you mean by...? / Could you repeat that? & Qu'entends-tu par... ? / Peux-tu répéter ? \\
\hline
Opinion & In my opinion... / I wish I could... & À mon avis... / J'aimerais pouvoir... \\
\hline
Conclure & Thanks for telling me. / Let's keep in touch. & Merci de m'en avoir parlé. / Restons en contact. \\
\hline
\end{tabularx}
\end{center}

\begin{popfigure}
\begin{center}
\begin{tikzpicture}[node distance=0.55cm and 0.9cm, every node/.style={font=\small}]
\node[draw=popblue, fill=popblueL, rounded corners, align=center, minimum width=3.1cm, minimum height=0.9cm] (open) {\textbf{Opening}\\ Guess what!};
\node[draw=poporange, fill=poporangeL, rounded corners, align=center, minimum width=3.1cm, minimum height=0.9cm, right=of open] (ask) {\textbf{Asking}\\ Have you heard about...?};
\node[draw=popgreen, fill=popgreenL, rounded corners, align=center, minimum width=3.1cm, minimum height=0.9cm, right=of ask] (answer) {\textbf{Answering}\\ Yes, I attended it.};
\node[draw=poppurple, fill=poppurpleL, rounded corners, align=center, minimum width=3.1cm, minimum height=0.9cm, below=of answer] (react) {\textbf{Reacting / Following up}\\ Really? Tell me more.};
\node[draw=poppink, fill=poppinkL, rounded corners, align=center, minimum width=3.1cm, minimum height=0.9cm, left=of react] (close) {\textbf{Closing}\\ Thanks for telling me.};
\draw[-{Stealth}, thick, popdark] (open) -- (ask);
\draw[-{Stealth}, thick, popdark] (ask) -- (answer);
\draw[-{Stealth}, thick, popdark] (answer) -- (react);
\draw[-{Stealth}, thick, popdark] (react) -- (close);
\draw[-{Stealth}, thick, popmuted, dashed] (react.north) to[bend right=25] (ask.south);
\end{tikzpicture}
\end{center}
\legende{Déroulement d'un dialogue : Opening $\rightarrow$ Asking $\rightarrow$ Answering $\rightarrow$ Reacting/Following up $\rightarrow$ Closing. La flèche en pointillés montre qu'on peut boucler entre questions et réactions avant de conclure.}
\end{popfigure}

\begin{exoresolu}[Remettre le dialogue dans l'ordre]
Énoncé — Les répliques suivantes sont mélangées. Remets-les dans l'ordre logique d'un dialogue complet.\\
a) \eng{Really? Tell me more about it.}\\
b) \eng{Thanks for telling me about it. Let's keep in touch!}\\
c) \eng{Guess what! I attended a course on democracy in Buea.}\\
d) \eng{Well, we discussed human rights and free elections.}\\
e) \eng{In my opinion, it was very useful. I wish you could attend the next one.}
\tcblower
Solution détaillée — On suit la structure Opening $\rightarrow$ Reacting $\rightarrow$ Answering $\rightarrow$ Opinion $\rightarrow$ Closing :
\begin{enumerate}
\item \textbf{c)} : ouverture avec \eng{Guess what!} (annonce de la nouvelle).
\item \textbf{a)} : réaction (\eng{Really?}) suivie d'une relance (\eng{Tell me more}).
\item \textbf{d)} : réponse à la relance (\eng{Well, we discussed...}).
\item \textbf{e)} : avis personnel (\eng{In my opinion...}) et souhait (\eng{I wish...}).
\item \textbf{b)} : conclusion polie (\eng{Thanks for telling me...}).
\end{enumerate}
Ordre correct : \textbf{c — a — d — e — b}.
\end{exoresolu}

\begin{exercice}[À toi de jouer]
Complète ce dialogue avec une formule adaptée à chaque fonction indiquée entre parenthèses.\\
\textbf{A:} \dotfill (ouvrir) I joined a civic education workshop in Douala!\\
\textbf{B:} \dotfill (réagir) \dotfill (relancer)\\
\textbf{A:} We learned about citizens' duties and elections.\\
\textbf{B:} Sorry, \dotfill (demander une précision sur “duties”)\\
\textbf{A:} It means our responsibilities as citizens.\\
\textbf{B:} \dotfill (donner ton avis)\\
\textbf{A:} There is another session next month.\\
\textbf{B:} \dotfill (conclure poliment)
\end{exercice}

\begin{corrige}[Exercice « À toi de jouer »]
Réponses possibles : \eng{Guess what!} (ouvrir) ; \eng{Really? That's great!} (réagir) ; \eng{Tell me more about it.} (relancer) ; \eng{what do you mean by “duties”?} (précision) ; \eng{In my opinion, it sounds very useful. I wish I could attend it.} (avis) ; \eng{Thanks for telling me about it. I'll try to register. Let's keep in touch!} (conclusion).
\end{corrige}

\begin{aretenir}[L'essentiel de la section]
Un dialogue vivant suit cinq étapes : \textbf{ouvrir} (\eng{Guess what!}), \textbf{réagir} (\eng{Really? That's great!}), \textbf{relancer} (\eng{Tell me more. / What happened next?}), \textbf{préciser} (\eng{What do you mean by...?}) et \textbf{conclure} (\eng{Thanks for telling me about it.}). Retiens la technique « réaction courte + question en wh- », la structure \eng{Thanks for + -ing}, et le souhait \eng{I wish I could...}. Ces formules seront tes outils pour les jeux de rôle de la fin de la leçon.
\end{aretenir}

\coursec{Dialogues modèles commentés}

Dans la section précédente, nous avons appris à \cle{ouvrir} une conversation, à \cle{relancer} avec des questions, à \cle{réagir} avec enthousiasme et à \cle{conclure} poliment. Il est temps maintenant de voir ces formules fonctionnelles \emph{en action}, dans deux dialogues complets et authentiques. Nous allons d'abord \textbf{observer} les dialogues, puis \textbf{repérer} les formules, et enfin \textbf{analyser} la différence de registre entre une conversation entre amis et un échange avec un organisateur.

\courssub{Dialogue 1 : entre amis (registre informel)}

\begin{textebox}[Dialogue 1 — Talking about a training course]
Tabi: Hi Amina! Long time no see.\\
Amina: Hi! Guess what? I attended a training course on democracy during the holidays.\\
Tabi: Really? Who organized it?\\
Amina: A local NGO, with ELECAM.\\
Tabi: How long did it last?\\
Amina: Three days, in Yaoundé.\\
Tabi: What did you learn?\\
Amina: Mainly voting rights and peaceful elections.\\
Tabi: That sounds great! Did you get a certificate?\\
Amina: Yes, I did.\\
Tabi: Wow! I'd love to attend the next one.\\
Amina: I'll tell you when registration opens.
\end{textebox}

\textbf{Glossaire :}

\begin{center}
\begin{tabularx}{\linewidth}{|X|X|X|}\hline
\rowcolor{popblueL} English & Nature & Français \\ \hline
long time no see & expression & ça fait longtemps (qu'on ne s'est pas vus) \\ \hline
guess what? & expression & devine quoi ? \\ \hline
to attend & verbe & assister à, suivre (une formation) \\ \hline
a training course & nom & une formation, un stage \\ \hline
an NGO & nom & une ONG \\ \hline
to last & verbe & durer \\ \hline
voting rights & nom & droits de vote \\ \hline
peaceful elections & nom & des élections pacifiques \\ \hline
a certificate & nom & une attestation, un certificat \\ \hline
registration & nom & l'inscription \\ \hline
\end{tabularx}
\end{center}

\begin{exemplebox}[Repérage des formules dans le Dialogue 1]
\begin{itemize}
\item \textbf{Ouverture :} \eng{Hi Amina! Long time no see.} — salutation chaleureuse entre amis.
\item \textbf{Accroche / annonce :} \eng{Guess what?} — pour créer la curiosité avant de donner l'information.
\item \textbf{Questions de relance :} \eng{Who organized it?}, \eng{How long did it last?}, \eng{What did you learn?}, \eng{Did you get a certificate?}
\item \textbf{Réactions :} \eng{Really?}, \eng{That sounds great!}, \eng{Wow!}
\item \textbf{Clôture :} \eng{I'll tell you when registration opens.} — promesse qui termine naturellement l'échange.
\end{itemize}
\end{exemplebox}

\begin{exemplebox}[Exemples traduits]
\eng{I'll tell you when registration opens.} « Je te dirai quand les inscriptions ouvriront. »\\
\eng{That sounds great!} « Ça a l'air génial ! »\\
\eng{I'd love to attend the next one.} « J'aimerais beaucoup assister à la prochaine session. »
\end{exemplebox}

\courssub{Dialogue 2 : avec un organisateur (registre semi-formel)}

\begin{textebox}[Dialogue 2 — Asking for information about a workshop]
Student: Good morning, Sir. I'd like some information about your next workshop on democracy.\\
Organizer: Certainly. It will be held on 15th March at the Community Hall.\\
Student: How can I register?\\
Organizer: Just fill in this form. It's free of charge.\\
Student: Thank you very much.\\
Organizer: You're welcome.
\end{textebox}

\textbf{Glossaire :}

\begin{center}
\begin{tabularx}{\linewidth}{|X|X|X|}\hline
\rowcolor{popblueL} English & Nature & Français \\ \hline
I'd like (= I would like) & expression & je voudrais (poli) \\ \hline
a workshop & nom & un atelier \\ \hline
to be held & verbe (passif) & avoir lieu, se tenir \\ \hline
the Community Hall & nom & la salle communale \\ \hline
to register & verbe & s'inscrire \\ \hline
to fill in a form & expression & remplir un formulaire \\ \hline
free of charge & expression & gratuit \\ \hline
you're welcome & expression & je vous en prie, de rien \\ \hline
\end{tabularx}
\end{center}

\begin{exemplebox}[Exemples traduits]
\eng{It's free of charge.} « C'est gratuit. »\\
\eng{It will be held on 15th March at the Community Hall.} « Il aura lieu le 15 mars à la salle communale. »\\
\eng{Just fill in this form.} « Remplissez simplement ce formulaire. »
\end{exemplebox}

\begin{propriete}[La date dans un dialogue : “on 15th March”]
Pour donner une date précise, on utilise la préposition \cle{on} : \eng{on 15th March}, \eng{on Monday}, \eng{on 15th March 2025}. À l'oral, on prononce la date avec l'article défini et l'ordinal : \eng{on the fifteenth of March} (prononciation approximative : « one ze fif-ti-nth of martch »). Attention : pas de préposition \eng{in} devant une date précise (\eng{in} s'emploie pour les mois et les années : \eng{in March}, \eng{in 2025}).
\end{propriete}

\courssub{Analyse guidée des deux dialogues}

\begin{methode}[Comment analyser un dialogue en trois couleurs]
\begin{enumerate}
\item \textbf{Souligne} toutes les questions en \cle{wh-} : elles servent à demander une information précise (\eng{Who organized it?}, \eng{How long did it last?}, \eng{What did you learn?}, \eng{How can I register?}).
\item \textbf{Encadre} les réactions : courtes exclamations qui montrent l'intérêt (\eng{Really?}, \eng{Wow!}, \eng{That sounds great!}, \eng{Certainly.}).
\item \textbf{Entoure} les formules de politesse : ouvertures, remerciements et clôtures (\eng{Hi!}, \eng{Good morning, Sir.}, \eng{I'd like...}, \eng{Thank you very much.}, \eng{You're welcome.}).
\end{enumerate}
\end{methode}

\begin{exoresolu}[Repérage des formules dans le Dialogue 2]
Énoncé : dans le Dialogue 2, (a) souligne la question en wh-, (b) encadre la réaction de l'organisateur, (c) entoure deux formules de politesse.
\tcblower
Solution détaillée :\\
(a) La question en wh- est \eng{How can I register?} — elle demande un moyen, une procédure.\\
(b) La réaction de l'organisateur est \eng{Certainly.} — réponse brève et positive qui montre la disponibilité.\\
(c) Formules de politesse : \eng{Good morning, Sir.} (ouverture formelle), \eng{I'd like some information...} (demande polie avec \eng{would like}), \eng{Thank you very much.} et \eng{You're welcome.} (remerciement et réponse).
\end{exoresolu}

\courssub{Comparer les deux registres}

\begin{popfigure}
\begin{tikzpicture}[
  box/.style={draw=popline, rounded corners=4pt, align=left, text width=6.2cm, inner sep=6pt, font=\small},
  head/.style={align=center, font=\bfseries\small}
]
\node[head, text=popblue, align=center] (h1) at (0,0){Registre informel\\(entre amis)};
\node[head, text=poporange, align=center] (h2) at (8,0){Registre semi-formel\\(avec un organisateur)};
\node[box, fill=popblueL, align=center] at (0,-2.2){\textbf{Ouverture :} Hi! Long time no see.\\[2pt]\textbf{Annonce :} Guess what?\\[2pt]\textbf{Réactions :} Really? Wow!\\[2pt]\textbf{Demande :} Who organized it?};
\node[box, fill=poporangeL, align=center] at (8,-2.2){\textbf{Ouverture :} Good morning, Sir.\\[2pt]\textbf{Annonce :} I'd like some information about...\\[2pt]\textbf{Réactions :} Certainly.\\[2pt]\textbf{Demande :} How can I register?};
\node[box, fill=popgreenL, text width=13.4cm] at (4,-4.6) {\textbf{Point commun :} dans les deux cas, on pose des questions en wh-, on réagit à l'information et on conclut poliment. Seul le \emph{niveau de langue} change.};
\end{tikzpicture}
\legende{Deux registres, une même fonction : échanger des informations sur une formation.}
\end{popfigure}

\begin{attention}[Ne mélange pas les registres !]
Au registre formel ou semi-formel, évite \eng{Hi} et \eng{Guess what} ; préfère \eng{Good morning} et \eng{I would like some information about...}. Inversement, dire \eng{Good morning, Sir} à ton meilleur ami sonnerait bizarre ! Autre piège : \eng{I would like} se contracte en \eng{I'd like} à l'oral, mais jamais \eng{I will like} pour exprimer un souhait poli.
\end{attention}

\begin{center}
\begin{tabularx}{\linewidth}{|X|X|X|}\hline
\rowcolor{popblueL} Fonction & Informel (amis) & Semi-formel (organisateur) \\ \hline
Saluer & Hi! / Hello! & Good morning, Sir/Madam. \\ \hline
Demander une information & Who organized it? & I'd like some information about... \\ \hline
Réagir & Really? Wow! & Certainly. / I see. \\ \hline
Remercier & Thanks! & Thank you very much. \\ \hline
Conclure & See you! & You're welcome. Goodbye. \\ \hline
\end{tabularx}
\end{center}

\courssub{Prononciation, accentuation et intonation}

\begin{propriete}[Accent de mot (Word Stress) — Mots clés du thème]
En anglais, chaque mot polysyllabique a une syllabe tonique (plus forte, plus longue, plus haute). Marque l'accent sur la syllabe en \textbf{gras} :
\begin{itemize}
\item \eng{de\textbf{moc}racy} (democracy) — 4 syllabes, accent sur la 2\up{ème}.
\item \eng{\textbf{cer}tificate} (certificate) — 4 syllabes, accent sur la 1\up{ère}.
\item \eng{\textbf{work}shop} (workshop) — 2 syllabes, accent sur la 1\up{ère} (nom composé).
\item \eng{or\textbf{ga}nize} (organize) — 3 syllabes, accent sur la 2\up{ème}.
\item \eng{re\textbf{gis}tration} (registration) — 4 syllabes, accent sur la 2\up{ème}.
\item \eng{\textbf{hu}man \textbf{rights}} (human rights) — accent tonique sur \eng{hu-} et \eng{rights}.
\end{itemize}
\emph{Astuce :} les suffixes \eng{-tion}, \eng{-sion}, \eng{-ic}, \eng{-ity} attirent l'accent sur la syllabe qui les précède (\eng{regis\underline{tra}tion}, \eng{demo\underline{cra}cy}).
\end{propriete}

\begin{propriete}[Intonation de la phrase (Sentence Intonation)]
\begin{itemize}
\item \textbf{Questions en wh- (Who, What, How, When...)} : intonation \textbf{descendante} (\searrow) à la fin. \eng{Who organized it?} (\eng{Who organized it\searrow})
\item \textbf{Questions fermées (Yes/No)} : intonation \textbf{montante} (\nearrow) à la fin. \eng{Did you get a certificate?} (\eng{Did you get a certificate\nearrow})
\item \textbf{Affirmations / Informations} : intonation \textbf{descendante} (\searrow). \eng{It will be held on 15th March.} (\eng{...March\searrow})
\item \textbf{Réactions / Exclamations (Really?, Wow!, That sounds great!)} : intonation \textbf{montante-descendante} (\nearrow\searrow) pour marquer l'enthousiasme ou la surprise.
\item \textbf{Listes (voting rights, peaceful elections, and...)} : intonation \textbf{montante} (\nearrow) sur les éléments non finaux, \textbf{descendante} (\searrow) sur le dernier.
\end{itemize}
\end{propriete}

\begin{experience}[Entraînement à la prononciation en binôme]
\begin{enumerate}
\item \textbf{Écoute et répète :} ton professeur (ou l'audio) lit les mots de la liste ci-dessus. Répète en tapant dans les mains sur la syllabe tonique.
\item \textbf{Flèches d'intonation :} reprends le Dialogue 1. Trace une flèche \searrow ou \nearrow au bout de chaque réplique selon la règle. Vérifie avec ton voisin.
\item \textbf{Lecture expressive :} relis le Dialogue 2 en exagérant l'intonation montante de \eng{How can I register?} et la chute sur \eng{Certainly.}
\end{enumerate}
\end{experience}

\courssub{Variantes : adapter les dialogues à d'autres situations}

Les deux dialogues sont des \cle{modèles} : on garde la structure (ouverture → annonce → questions → réactions → clôture) et on remplace le contenu. Voici des substitutions possibles :

\begin{itemize}
\item \eng{a training course on democracy} → \eng{a seminar on human rights} (un séminaire sur les droits de l'homme) ;
\item \eng{a local NGO, with ELECAM} → \eng{a youth association in Bamenda} ;
\item \eng{voting rights and peaceful elections} → \eng{children's rights and freedom of expression} ;
\item \eng{your next workshop on democracy} → \eng{your awareness campaign on citizens' rights} (votre campagne de sensibilisation) ;
\item \eng{on 15th March at the Community Hall} → \eng{on 2nd April at the school hall}.
\end{itemize}

\begin{exemplebox}[Variante du Dialogue 1 (extrait)]
\eng{Tabi: What did you learn? — Amina: Mainly children's rights and how to report abuses.}\\
« Tabi : Qu'as-tu appris ? — Amina : Surtout les droits des enfants et comment signaler les abus. »
\end{exemplebox}

\begin{methode}[Mémoriser un dialogue sans l'apprendre mot à mot]
\begin{enumerate}
\item \textbf{Apprends les formules fonctionnelles} (ouverture, questions en wh-, réactions, clôture) : ce sont les « piliers » du dialogue.
\item \textbf{Repère la structure} en cinq étapes : saluer → annoncer → questionner → réagir → conclure.
\item \textbf{Remplace le contenu} par ta propre expérience : ta ville, ta formation, tes dates.
\item \textbf{Répète à voix haute} en binôme, puis inverse les rôles : chacun doit jouer les deux personnages.
\end{enumerate}
\end{methode}

\courssub{Activités orales : lecture, inversion des rôles, dialogue à trous}

\begin{experience}[Lecture à voix haute et reconstruction du dialogue]
\textbf{Étape 1 — Lecture en binômes.} Lis le Dialogue 1 à voix haute avec ton voisin, en soignant l'intonation montante des questions fermées (\eng{Really?} \nearrow, \eng{Did you get a certificate?} \nearrow) et l'enthousiasme des réactions (\eng{That sounds great!} \nearrow\searrow).\\
\textbf{Étape 2 — Inversion des rôles.} Échangez les rôles : celui qui jouait Tabi joue maintenant Amina.\\
\textbf{Étape 3 — Dialogue à trous.} Reconstruisez le dialogue à partir de ces mots-clés seulement, sans regarder le modèle :\\
\emph{Hi — long time — guess — training course — democracy — who — NGO — how long — three days — what — learn — voting rights — certificate — I'd love — registration.}
\end{experience}

\begin{exoresolu}[Dialogue à trous : complète les répliques]
Énoncé : complète le dialogue avec les formules étudiées.\\
Tabi: Hi Amina! \ldots\ldots\ldots\ldots\ldots\ldots\ (ça fait longtemps)\\
Amina: \ldots\ldots\ldots\ldots\ldots\ldots\ ? I attended a seminar on human rights. (devine quoi)\\
Tabi: \ldots\ldots\ldots\ldots\ldots\ldots\ ? (qui l'a organisé)\\
Amina: A youth association in Bamenda.\\
Tabi: \ldots\ldots\ldots\ldots\ldots\ldots\ ! (ça a l'air génial)
\tcblower
Solution détaillée :\\
Tabi: Hi Amina! \eng{Long time no see.}\\
Amina: \eng{Guess what}? I attended a seminar on human rights.\\
Tabi: \eng{Who organized it}?\\
Amina: A youth association in Bamenda.\\
Tabi: \eng{That sounds great}!\\
Remarque : la question \eng{Who organized it?} est au prétérit (\eng{organized}) car l'événement est terminé ; la réaction \eng{That sounds great!} reste au présent car elle commente l'information au moment où on l'entend.
\end{exoresolu}

\begin{savaistu}[Contexte camerounais : ELECAM et la formation civique]
\cle{ELECAM} (Elections Cameroon) est l'organe chargé de l'organisation et de la supervision des élections au Cameroun. Il collabore fréquemment avec des ONG locales et des organisations internationales (comme le PNUD ou l'UNICEF) pour organiser des \eng{civic education workshops} (ateliers d'éducation civique) et des \eng{voter education campaigns} (campagnes d'éducation des électeurs), notamment à l'approche des échéances électorales. Participer à ces formations permet d'obtenir un \eng{certificate of participation} souvent valorisé pour l'engagement associatif ou le bénévolat.
\end{savaistu}

\begin{aretenir}[L'essentiel de la section]
\begin{itemize}
\item Un dialogue d'échange d'informations suit cinq étapes : \textbf{saluer → annoncer → questionner (wh-) → réagir → conclure}.
\item Le \textbf{registre} s'adapte à l'interlocuteur : informel entre amis (\eng{Hi}, \eng{Guess what?}), semi-formel avec un organisateur (\eng{Good morning}, \eng{I'd like...}).
\item L'\textbf{accent de mot} est fixe (dictionary) : \eng{de\textbf{moc}racy}, \eng{\textbf{cer}tificate}, \eng{\textbf{work}shop}.
\item L'\textbf{intonation} signale la fonction : \searrow (wh-, affirmation), \nearrow (Yes/No), \nearrow\searrow (exclamation).
\item On mémorise les \textbf{formules fonctionnelles}, puis on adapte le contenu à sa propre situation.
\end{itemize}
\end{aretenir}

\begin{exercice}[À toi de jouer]
Avec un camarade, écris puis joue un dialogue de huit répliques : l'un de vous a participé à une campagne de sensibilisation (\eng{awareness campaign}) sur les droits des citoyens à Douala ; l'autre pose au moins trois questions en wh- et réagit deux fois. Terminez par une formule de clôture adaptée au registre choisi.
\end{exercice}

\coursec{Prononciation, accentuation et intonation}

Après avoir analysé la structure des dialogues modèles de la section précédente, nous devons maintenant nous concentrer sur la \cle{forme orale} de ces échanges. En anglais, le sens ne réside pas seulement dans les mots choisis, mais dans la \cle{musique de la phrase} : l'accentuation tonique, les sons caractéristiques, l'intonation et les liaisons. Un énoncé grammaticalement parfait mais prononcé avec l'intonation du français ne sera pas compris — ou pire, sera interprété comme impoli ou hésitant. Cette section vous donne les clés pour \textbf{sonner anglais} dans le contexte des formations sur la démocratie.

\courssub{L'accentuation des mots-clés du thème}

\begin{textebox}[Script d'écoute : Mots-clés du thème « Democracy Training »]
Listen carefully to the stress placement in these key words. The stressed syllable is LOUDER, LONGER and HIGHER in pitch.
\begin{itemize}
    \item par\textbf{TI}cipant (not \textbf{PAR}ticipant)
    \item \textbf{OR}ganizer (not or-ga-NI-zer)
    \item de\textbf{MO}cracy (not de-mo-CRA-cy)
    \item cer\textbf{TI}ficate (not \textbf{CER}tificate)
    \item \textbf{REG}ister (not re-GIS-ter)
    \item \textbf{WORK}shop (not work-SHOP)
\end{itemize}
\textit{Traduction : Écoutez attentivement la place de l'accent dans ces mots-clés. La syllabe accentuée est plus FORTE, plus LONGUE et plus AIGÜE.}
\end{textebox}

\begin{propriete}[Règle de l'accentuation tonique (Word Stress)]
En anglais, \textbf{chaque mot polysyllabique a une syllabe accentuée principale} (primary stress). Cette syllabe porte l'information nouvelle. Contrairement au français où l'accent est toujours sur la dernière syllabe du groupe rythmique (accent d'intensité final), l'anglais place l'accent \textbf{à l'intérieur du mot}, selon des tendances :
\begin{itemize}
    \item \textbf{Noms et adjectifs à 3 syllabes et plus :} souvent sur l'antépénultième (3\up{e} en partant de la fin) : \cle{parTIcipant}, \cle{ORganizer}, \cle{cerTIficate}, \cle{deMOcracy}.
    \item \textbf{Verbes à 2 syllabes :} souvent sur la 2\up{e} (re\cle{CORD}, pre\cle{SENT}), \textbf{mais} les verbes en \textit{-er} / \textit{-ter} / \textit{-ster} de 3 syllabes accentuent souvent la 1\up{re} : \cle{REGister}, \cle{ENTer}, \cle{OFFer}.
    \item \textbf{Mots en \textit{-tion/-sion} :} toujours sur la syllabe avant le suffixe (par\cle{TI}cipation, or\cle{GA}ni\cle{za}tion).
    \item \textbf{Mots composés (noms) :} accent sur le premier élément (\cle{WORK}shop, \cle{TRAIN}ing course).
    \item \textbf{Familles de mots :} l'accent peut bouger : \textbf{deMO}cracy $\to$ demo\cle{CRA}tic ; \textbf{parTI}cipant $\to$ par\cle{TI}cipate.
\end{itemize}
\textbf{Erreur fatale :} Transférer l'accent final français (\textit{partici\cle{PANT}}, \textit{organi\cle{ZER}}, \textit{certifi\cle{CATE}}) rend le mot méconnaissable pour un anglophone.
\end{propriete}

\begin{center}
\begin{tabularx}{\linewidth}{|X|X|X|X|}
\hline
\rowcolor{popblueL} \textbf{Mot (English)} & \textbf{Syllabe accentuée} & \textbf{Prononciation approximative (FR)} & \textbf{Traduction} \\
\hline
Participant & par\textbf{TI}cipant & « par-\textbf{TI}-si-pente » & Participant \\
\hline
Organizer & \textbf{OR}ganizer & « \textbf{OR}-ga-naï-zeur » & Organisateur \\
\hline
Democracy & de\textbf{MO}cracy & « di-\textbf{MO}-kra-si » & Démocratie \\
\hline
Certificate & cer\textbf{TI}ficate & « sèr-\textbf{TI}-fi-kète » & Certificat \\
\hline
Register (v.) & \textbf{REG}ister & « \textbf{REDJ}-is-teur » & S'inscrire \\
\hline
Workshop & \textbf{WORK}shop & « \textbf{OUÈRK}-chop » & Atelier \\
\hline
\end{tabularx}
\end{center}

\begin{exoresolu}[Marquage de l'accentuation]
\textbf{Énoncé :} Écoutez l'énoncé suivant et soulignez la syllabe accentuée dans chaque mot en gras. Puis entraînez-vous à le lire à voix haute en tapant dans les mains sur la syllabe forte.

\textit{``The \textbf{organizer} asked the \textbf{participant} to \textbf{register} for the \textbf{workshop} on \textbf{democracy} to get a \textbf{certificate}.''}

\tcblower
\textbf{Solution détaillée :}
\begin{enumerate}
    \item \textbf{OR}ganizer (nom, 1\up{re} syllabe)
    \item par\textbf{TI}cipant (nom, 3 syllabes $\to$ antépénultième)
    \item \textbf{REG}ister (verbe 3 syllabes en -ter $\to$ 1\up{re} syllabe)
    \item \textbf{WORK}shop (nom composé $\to$ 1\up{er} élément)
    \item de\textbf{MO}cracy (exception famille : demo\cle{CRA}tic)
    \item cer\textbf{TI}ficate (nom 4 syllabes $\to$ antépénultième)
\end{enumerate}
\textbf{Conseil :} Dites la phrase en frappant des mains 6 fois, une par mot accentué. Le rythme doit être régulier : \textsc{TAC}... \textsc{TAC}... \textsc{TAC}... \textsc{TAC}... \textsc{TAC}... \textsc{TAC}.
\end{exoresolu}

\courssub{Les sons difficiles pour les francophones}

Trois sons du thème « Democracy Training » posent systématiquement problème aux locuteurs francophones camerounais. Maîtrisez-les pour gagner en intelligibilité immédiate.

\begin{propriete}[Le \textit{th} sonore /ð/ et sourd /θ/ : « La langue entre les dents »]
L'anglais possède deux sons \textit{th} absents du français. La pointe de la langue doit \textbf{sortir légèrement entre les incisives supérieures et inférieures}.
\begin{itemize}
    \item \textbf{/θ/ (sourd, sans vibration des cordes vocales) :} \cle{three} (« sri » — souffle seulement), \cle{think}, \cle{thirty}, \cle{Thursday}.
    \item \textbf{/ð/ (sonore, avec vibration) :} \cle{the} (« ze »), \cle{that} (« zate »), \cle{those} (« zoz »), \cle{other} (« aveur »), \cle{them} (« zème »).
\end{itemize}
\textbf{Astuce :} Placez un doigt verticalement sur vos lèvres. Si vos lèvres touchent le doigt en disant \textit{three}, c'est raté (vous faites /s/ ou /t/). La langue doit chatouiller le doigt.
\end{propriete}

\begin{propriete}[Le \textit{w} /w/ : « Arrondir les lèvres »]
Le \textit{w} anglais (\cle{workshop}, \cle{work}, \cle{will}, \cle{where}) n'est \textbf{pas} le \textit{ou} français (comme dans \textit{oui}). C'est une \textbf{semi-voyelle} : les lèvres sont \textbf{très arrondies et proéminentes} (comme pour un baiser exagéré), puis s'ouvrent vite vers la voyelle suivante. Ne dites pas « \textit{ou-ork-chop} » mais « \textbf{ouèrk}-chop » (transition rapide).
\end{propriete}

\begin{propriete}[Le \textit{r} anglais /r/ : « Non roulé, retroflexe »]
Le \textit{r} de \cle{rights} (« raïts »), \cle{training} (« trèï-ning »), \cle{organizer} ne se roule \textbf{jamais} (pas de vibration de la luette comme en français \textit{rouge} ou en bassa/fulfulde). La pointe de la langue se recourbe en arrière vers le palais dur (position rétroflexe) sans le toucher. Les lèvres sont légèrement arrondies. Prononcez « \textbf{raïts} » (lèvres en avant), pas « \textit{rrraits} ».
\end{propriete}

\begin{exemplebox}[Phrases d'entraînement ciblé (Minimal Pairs \& Context)]
\begin{enumerate}
    \item \eng{The \textbf{three} \textbf{workshops} start on \textbf{Thursday}.} (« Ze \textbf{sri} \textbf{ouèrk-chops} start on \textbf{sèrs-deï} » — \textit{th} sourd partout).
    \item \eng{The \textbf{organizer} \textbf{registered} \textbf{thirty} \textbf{participants}.} (« Zi \textbf{or-ga-naï-zeur} \textbf{redj-is-teureud} \textbf{sèr-ti} \textbf{par-ti-si-pentes} » — \textit{th} sourd dans \textit{thirty}, \textit{th} sonore dans \textit{the}, \textit{r} retroflexe partout).
    \item \eng{\textbf{Work} hard for your \textbf{rights}.} (« \textbf{Ouèrk} hard for your \textbf{raïts} » — \textit{w} arrondi, \textit{r} non roulé).
\end{enumerate}
\end{exemplebox}

\begin{attention}[Pièges classiques « Cameroun Spécial »]
\begin{itemize}
    \item \textbf{``H'' muet vs ``H'' aspiré :} N'oubliez pas : \cle{honest} /ˈɒnɪst/ (h muet $\to$ \textit{an honest man}) mais \cle{held} /held/ (\textbf{h aspiré} $\to$ \textit{the workshop was held} /ðə wɜːkʃɒp wɒz hɛld/). Beaucoup d'élèves disent « \textit{eld} » pour \textit{held} $\to$ incompréhensible.
    \item \textbf{``Democracy'' :} L'accent est sur \textbf{MO} (2\up{e} syllabe) : « di-\textbf{MO}-kra-si ». Erreur fréquente : *\textit{DE-mo-cra-cy} (accent français) ou *\textit{de-mo-CRA-cy} (confusion avec \textit{demo\cle{CRA}tic}).
    \item \textbf{``Certificate'' :} Ne dites pas *\textit{cer-ti-fi-CATE} (à la française) ni *\textit{CER-ti-fi-cate}. Dites \textbf{cer-TI-fi-cate} /səˈtɪfɪkət/ $\to$ « sèr-\textbf{TI}-fi-kète ».
    \item \textbf{``Participant'' :} Ne dites pas *\textit{PAR-ti-ci-pant}. Dites \textbf{par-TI-ci-pant} $\to$ « par-\textbf{TI}-si-pente ».
    \item \textbf{``Register'' (verbe) :} Ne dites pas *\textit{re-GIS-ter}. Dites \textbf{REG-is-ter} $\to$ « \textbf{REDJ}-is-teur ».
    \item \textbf{``NGO'' :} Se prononce lettre par lettre : /en dʒiː əʊ/ $\to$ « ène-dji-ou ». Pas *\textit{ène-go}.
\end{itemize}
\end{attention}

\courssub{L'intonation : questions et réactions}

L'intonation (la mélodie de la phrase) signale la \cle{fonction illocutoire} : est-ce une question fermée ? Une question ouverte ? Une surprise ? Un accord ? En français, l'intonation monte souvent à la fin pour toute question. En anglais, \textbf{c'est l'inverse pour les questions en \textit{wh-}}.

\begin{popfigure}
\centering
\begin{tikzpicture}[
    arrowstyle/.style={->, >=Stealth, thick},
    labelstyle/.style={font=\small\bfseries, text=popdark},
    axis/.style={->, >=Stealth, thin, popmuted}
]
% Axes communs
\def\xmax{6}
\def\ymax{3}

% Flèche MONTANTE (Yes/No)
\begin{scope}[xshift=0cm]
    \draw[axis] (0,0) -- (\xmax,0) node[right, popmuted] {Time};
    \draw[axis] (0,0) -- (0,\ymax) node[above, popmuted] {Pitch};
    % Courbe montante
    \draw[popblue, thick, decorate, decoration={snake, amplitude=1.5pt, segment length=10pt}] (0.5,1) -- (5.5,2.5);
    \node[labelstyle, popblue] at (3, 2.8) {Yes/No Question ↗};
    \node[labelstyle, popblue] at (3, 0.5) {\eng{Did you ENjoy it?}};
    \node[font=\tiny, popmuted] at (5.5, 2.5) {High};
    \node[font=\tiny, popmuted] at (0.5, 1) {Mid};
\end{scope}

% Flèche DESCENDANTE (Wh-)
\begin{scope}[xshift=9cm]
    \draw[axis] (0,0) -- (\xmax,0) node[right, popmuted] {Time};
    \draw[axis] (0,0) -- (0,\ymax) node[above, popmuted] {Pitch};
    % Courbe descendante
    \draw[poporange, thick, decorate, decoration={snake, amplitude=1.5pt, segment length=10pt}] (0.5,2.5) -- (5.5,1);
    \node[labelstyle, poporange] at (3, 2.8) {Wh- Question ↘};
    \node[labelstyle, poporange] at (3, 0.5) {\eng{WHERE was it HELD?}};
    \node[font=\tiny, popmuted] at (0.5, 2.5) {High};
    \node[font=\tiny, popmuted] at (5.5, 1) {Low};
\end{scope}
\end{tikzpicture}
\legende{Schéma d'intonation : montante (bleu) pour questions fermées, descendante (orange) pour questions ouvertes.}
\end{popfigure}

\begin{propriete}[Règle d'or de l'intonation interrogative]
\begin{itemize}
    \item \textbf{Questions fermées (Yes/No) :} Intonation \textbf{montante} ↗ en fin d'énoncé. Elles attendent \textit{Yes} ou \textit{No}. Le pic de hauteur est sur le dernier mot fort.
    \item \textbf{Questions ouvertes (Wh- : Who, What, Where, When, Why, How) :} Intonation \textbf{descendante} ↘ en fin d'énoncé. Elles attendent une information. Le pic est sur le mot interrogatif (ou le premier mot fort), puis ça redescend.
\end{itemize}
\textbf{Pourquoi ?} La question \textit{wh-} montre que le locuteur \textbf{sait déjà} qu'il y a une réponse (présupposition), donc il affirme sa demande (ton affirmatif $\to$ descendant). La question \textit{yes/no} exprime l'incertitude totale (ton interrogatif $\to$ montant).
\end{propriete}

\begin{exemplebox}[Exemples traduits — Intonation des questions]
\begin{center}
\begin{tabularx}{\linewidth}{|X|X|X|}
\hline
\rowcolor{popblueL} \textbf{English (avec flèche)} & \textbf{Prononciation approx. (FR)} & \textbf{Français} \\
\hline
\eng{Did you reCEIVE a cerTIficate? ↗} & « Di-dyou ri-\textbf{SEIVE} eu \textbf{SÈR}-ti-fi-kète ↗ » & As-tu reçu un certificat ? \\
\hline
\eng{Will the ORganizer be THERE? ↗} & « Ouile ze \textbf{OR}-ga-naï-zeur bi \textbf{ZÈRE} ↗ » & L'organisateur sera-t-il là ? \\
\hline
\eng{WHAT did you LEARN? ↘} & « \textbf{OUATE} di-dyou \textbf{LEURNE} ↘ » & Qu'as-tu appris ? \\
\hline
\eng{WHERE was the WORKshop HELD? ↘} & « \textbf{OUÈR} woz ze \textbf{OUÈRK}-chop \textbf{HÈLDE} ↘ » & Où l'atelier a-t-il eu lieu ? \\
\hline
\eng{HOW much does it COST? ↘} & « \textbf{HAO} matche daz it \textbf{KOST} ↘ » & Combien ça coûte ? \\
\hline
\end{tabularx}
\end{center}
\end{exemplebox}

\begin{propriete}[Intonation des réactions (Backchanneling)]
Dans un échange oral, les réactions courtes maintiennent la conversation. Leur intonation change le sens :
\begin{itemize}
    \item \eng{REAlly? ↗} (« \textbf{RI}-li ↗ ») : \textbf{Surprise, incrédulité, intérêt vif}. Montée haute.
    \item \eng{Really. ↘} (« \textbf{RI}-li ↘ ») : \textbf{Doute, scepticisme, ironie, lassitude}. Descente plate.
    \item \eng{That's GREAT! ↘} (« Dats \textbf{GRÈITE} ↘ ») : \textbf{Enthousiasme sincère, accord fort}. Descente énergique depuis un pic haut.
    \item \eng{That's great... ↗} (« Dats grèite ↗ ») : \textbf{Politesse distante, ironie possible}. Montée finale = « et alors ? ».
\end{itemize}
\end{propriete}

\begin{savaistu}[Le pouvoir de ``Really?``]
Saviez-vous que \eng{Really?} est un caméléon ? Avec une intonation \textbf{montante ↗} (\textit{« Ri-li ? »} voix haute), vous dites : « Incroyable ! Dis-moi plus ! ». Avec une intonation \textbf{descendante ↘} (\textit{« Ri-li. »} voix basse, plate), vous dites : « Mouais... j'en doute » ou « C'est ennuyeux ». En contexte camerounais, entre amis, un \eng{Really? ↗} sonore après « I met the Minister! » montre l'admiration ; un \eng{Really. ↘} sec peut clore la discussion. Maîtrisez cette nuance pour ne pas passer pour impoli !
\end{savaistu}

\courssub{Les liaisons naturelles (Enchaînements) : la fluidité}

L'anglais parlé n'est pas une suite de mots isolés. Les mots s'enchaînent par des \cle{liaisons} (linking) pour former des groupes rythmiques. Trois types sont essentiels ici :

\begin{definition}[Types de liaisons dans le thème]
\begin{enumerate}
    \item \textbf{Consonne $\to$ Voyelle (C-V) :} La consonne finale « saute » sur la voyelle initiale du mot suivant.
    \item \textbf{Voyelle $\to$ Voyelle (V-V) :} Insertion d'un /j/ (y), /w/ ou /r/ (intrusif) pour éviter le hiatus.
    \item \textbf{Liaison \textit{r} (Linking R) :} Un \textit{r} final muet (devant consonne ou pause) se prononce si le mot suivant commence par une voyelle.
    \item \textbf{Élision / Assimilation :} Disparition ou modification d'un son pour faciliter l'enchaînement (ex. \textit{last time} $\to$ « lasse taïme »).
\end{enumerate}
\end{definition}

\begin{exemplebox}[Liaisons clés du thème « Democracy Training »]
\begin{center}
\begin{tabular}{|l|l|l|}
\hline
\rowcolor{popblueL} \textbf{Écrit (Citation form)} & \textbf{Parlé (Connected speech)} & \textbf{Prononciation approx. (FR)} \\
\hline
an NGO & an\_NGO & « è-\textbf{ne}-dji-\textbf{ou} » \\
\hline
took part in & took\_part\_in & « \textbf{TOUK}-\textbf{PAR}-\textbf{TIN} » (liaison \textit{r}) \\
\hline
tell me about it & tell\_me\_a\_bout\_it & « \textbf{TEL}-\textbf{MI}-\textbf{YA}-\textbf{BOU}-\textbf{TIT} » \\
\hline
register for & register\_for & « \textbf{REDJ}-is-te-\textbf{RE}-\textbf{FORE} » (liaison \textit{r}) \\
\hline
participant in & participant\_in & « par-\textbf{TI}-si-\textbf{PA}-\textbf{TIN} » \\
\hline
\end{tabular}
\end{center}
\textbf{Règle pratique :} Ne faites \textbf{jamais de pause} entre ces mots. Imaginez un seul gros mot : \textit{anNGO}, \textit{tookpartin}, \textit{tellmeaboutit}.
\end{exemplebox}

\begin{methode}[S'entraîner à la fluidité : La technique du ``Shadowing`` (Ombre)]
\begin{enumerate}
    \item \textbf{Écoutez} le modèle (professeur ou audio) une fois sans parler.
    \item \textbf{Murmurez} en même temps (décalage de 0,5 seconde), en imitant \textbf{exactement} la mélodie, les liaisons, les accents.
    \item \textbf{Enregistrez-vous} sur votre téléphone (mémo vocal).
    \item \textbf{Comparez} : votre courbe d'intonation monte-t-elle aux bons endroits ? Vos \textit{th} sortent-ils ? Vos \textit{r} sont-ils non roulés ?
    \item \textbf{Répétez} jusqu'à ce que votre version se superpose au modèle.
\end{enumerate}
\textbf{Application immédiate :} Prenez le dialogue de la Section 4 (Dialogues modèles). Appliquez le \textit{shadowing} sur 2 minutes par jour pendant une semaine.
\end{methode}

\begin{exoresolu}[Diagnostic global : Prononciation \& Intonation]
\textbf{Énoncé :} L'élève lit l'échange suivant. Le professeur note les erreurs. Corrigez les 6 erreurs majeures ci-dessous.

\textit{Élève A :} ``Did you go to the workshop? ↗'' \\
\textit{Élève B :} ``Yes. It was good. ↘'' \\
\textit{Élève A :} ``What did you learn? ↗'' \textbf{(Erreur 1)} \\
\textit{Élève B :} ``I learned about deMOcracy. ↘'' \textbf{(Erreur 2)} \\
\textit{Élève A :} ``Realli? ↘'' \textbf{(Erreur 3)} \\
\textit{Élève B :} ``Yes. The organiZER gave me a certifiCATE. ↘'' \textbf{(Erreur 4, 5)} \\
\textit{Élève A :} ``Tell me about it. ↗'' \textbf{(Erreur 6 : intonation affirmation)} \\

\tcblower
\textbf{Solution détaillée :}
\begin{enumerate}
    \item \textbf{Erreur 1 :} \eng{What did you learn?} est une question \textit{wh-} $\to$ intonation \textbf{descendante ↘} (pas montante). Correction : \eng{What did you LEARN? ↘}
    \item \textbf{Erreur 2 :} \eng{deMOcracy} : accent sur \textbf{MO} (2\up{e} syllabe). L'élève a mis l'accent sur \textit{DE} ou \textit{CRA}. Correction : \eng{de\cle{MO}cracy}.
    \item \textbf{Erreur 3 :} \eng{Really?} avec ↘ exprime le scepticisme. Pour montrer l'intérêt : \textbf{montante ↗}. Correction : \eng{REAlly? ↗}
    \item \textbf{Erreur 4 :} \eng{organiZER} : accent français sur dernière syllabe. Correction : \eng{\cle{OR}ganizer}.
    \item \textbf{Erreur 5 :} \eng{certifiCATE} : accent français. Correction : \eng{cer\cle{TI}ficate}.
    \item \textbf{Erreur 6 :} \eng{Tell me about it} est une injonction/affirmation (impératif doux) $\to$ intonation \textbf{descendante ↘} (pas montante). Liaison : \eng{tell\_me\_a\_bout\_it}.
\end{enumerate}
\textbf{Bilan :} 3 erreurs d'accentuation lexicale (ORganizer, cerTIficate, deMOcracy), 3 erreurs d'intonation de phrase. Priorité : corriger l'accent des mots (ORganizer, cerTIficate, deMOcracy, parTIcipant, REGister) car ils reviennent sans cesse.
\end{exoresolu}

\courssub{Vers la production orale : Jeux de rôle guidés}

Vous maîtrisez maintenant les \textbf{outils phonologiques} (sons, accents, intonation, liaisons) pour incarner les \textbf{fonctions langagières} vues en Section 3 (demander, donner, relancer, conclure). La Section 6 vous proposera des \cle{jeux de rôle} où vous devrez mobiliser tout cela en situation quasi-réelle. Rappelez-vous : \textbf{la prononciation n'est pas un ornement, c'est le véhicule du sens}. Un \eng{REAlly? ↗} bien placé vaut mieux qu'un long discours mal accentué.

\begin{aretenir}[Check-list phonologique pour le jeu de rôle (Section 6)]
Avant de parler, cochez mentalement :
\begin{itemize}
    \item [ ] Mes mots-clés sont accentués : par\textbf{TI}cipant, \textbf{OR}ganizer, de\textbf{MO}cracy, cer\textbf{TI}ficate, \textbf{REG}ister, \textbf{WORK}shop.
    \item [ ] Mes \textit{th} sortent la langue (\textit{three, the, that, thirty, Thursday}).
    \item [ ] Mes \textit{w} arrondissent les lèvres (\textit{workshop, work, will, where}).
    \item [ ] Mes \textit{r} ne roulent pas (\textit{rights, training, organizer}).
    \item [ ] Questions Yes/No $\to$ ↗ ; Questions Wh- $\to$ ↘.
    \item [ ] Réactions : \eng{Really? ↗} (intérêt), \eng{That's great! ↘} (enthousiasme).
    \item [ ] Liaisons : \eng{an\_NGO, took\_part\_in, tell\_me\_about\_it, register\_for}.
\end{itemize}
\end{aretenir}

\coursec{Jeux de rôle guidés et production orale}

Après avoir travaillé la \textbf{prononciation}, l'\textbf{accentuation} et l'\textbf{intonation} dans la section précédente, nous passons maintenant à la mise en pratique communicative. Connaître les sons et le rythme ne suffit pas : il faut savoir \textbf{construire un échange vivant}, \textbf{réagir à l'imprévu} et \textbf{structurer son propos} pour être compris et évalué positivement. C'est l'objectif de cette dernière section : transformer vos connaissances en \textbf{compétence interactionnelle} grâce à trois jeux de rôle à autonomie progressive.

\courssub{De la prononciation à l'interaction : l'escalier de l'autonomie}

Pour réussir l'épreuve orale du Baccalauréat, il faut gravir les marches de l'autonomie. Le schéma ci-dessous illustre la progression que nous venons de suivre et que vous devez reproduire lors de vos révisions personnelles.

\begin{popfigure}
\begin{tikzpicture}[
    node distance=1.3cm and 0.5cm,
    box/.style={draw=popblue, thick, rounded corners, fill=popblueL, minimum width=4.2cm, minimum height=1.1cm, align=center, font=\small},
    arrow/.style={-Stealth, thick, popdark}
]
\node[box, align=center] (n1){Dialogue lu\\\emph{(Modèle audio / texte)}};
\node[box, below=of n1, align=center] (n2){Dialogue à trous\\\emph{(Complétion guidée)}};
\node[box, below=of n2, align=center] (n3){Jeu de rôle guidé\\\emph{(Fiches A / B imposées)}};
\node[box, below=of n3, align=center] (n4){Jeu de rôle semi-guidé\\\emph{(Rôle + objectif, contenu libre)}};
\node[box, below=of n4, align=center] (n5){Jeu de rôle libre\\\emph{(Situation ouverte, improvisation)}};

\draw[arrow] (n1) -- (n2) node[midway, right=0.2cm, font=\tiny, text=popmuted] {Repérage};
\draw[arrow] (n2) -- (n3) node[midway, right=0.2cm, font=\tiny, text=popmuted] {Appropriation};
\draw[arrow] (n3) -- (n4) node[midway, right=0.2cm, font=\tiny, text=popmuted] {Transfert};
\draw[arrow] (n4) -- (n5) node[midway, right=0.2cm, font=\tiny, text=popmuted] {Autonomie};
\end{tikzpicture}
\legende{L'escalier de l'autonomie orale : du modèle imité à l'improvisation maîtrisée.}
\end{popfigure}

\begin{aretenir}[Objectif de la section]
À l'issue de cette section, vous devez être capable de :
\begin{itemize}
    \item Mener un échange d'informations complet sur une formation civique (droits de l'homme, processus électoral, etc.).
    \item Utiliser des \textbf{formules de relance}, de \textbf{réaction} et de \textbf{conclusion} naturelles.
    \item Gérer vos hésitations avec des \textbf{bouche-trous} (fillers) anglais, sans recourir au français « euh ».
    \item Vous auto-évaluer sur la grille interactive (questions, réactions, vocabulaire, politesse).
\end{itemize}
\end{aretenir}

\courssub{Jeu de rôle 1 : Guidé — « Le workshop sur les droits de l'homme » (Fiches A / B)}

C'est l'exercice type de l'épreuve : vous avez une \textbf{fiche de rôle} contraignante. L'élève A \textbf{a vécu l'expérience} (données factuelles imposées), l'élève B \textbf{doit la découvrir} en posant des questions ciblées. La réussite repose sur la \textbf{précision des questions} (B) et la \textbf{cohérence des réponses} (A).

\begin{textebox}[Fiches de rôle — Jeu de rôle 1]
\textbf{FICHE A (Le participant)}\\
You attended a two-day workshop on human rights in Douala last month. It was organised by a local NGO called “Youth for Democracy”. There were 30 participants from different schools. You received a certificate of participation. You found it very useful because you learnt about the Universal Declaration of Human Rights and how to report violations. You met interesting people and made new friends.

\textbf{FICHE B (L'ami curieux)}\\
Your friend attended a workshop recently. You want to know everything about it. Find out:
\begin{enumerate}
    \item What the workshop was about (topic).
    \item Where and when it took place.
    \item How long it lasted.
    \item How many participants there were.
    \item What your friend thought about it (opinion, details).
    \item If he/she received anything (certificate, documents).
\end{enumerate}
\emph{Start the conversation naturally. React to the answers. Ask follow-up questions. Conclude politely.}
\end{textebox}

\begin{methode}[Réussir un jeu de rôle guidé — 5 étapes]
\begin{enumerate}
    \item \textbf{Lis ta fiche activement :} Surligne les \textbf{mots-clés} (topic, place, time, number, opinion, certificate). Si tu es B, transforme chaque point en \textbf{question ouverte} (Wh-) ou fermée (Yes/No).
    \item \textbf{Prépare 3 questions « de secours » :} Au cas où ton partenaire serait bref. Ex. \eng{Can you give me an example?} / \eng{What was the best moment?}
    \item \textbf{Écoute VRAIMENT la réponse :} Ne prépare pas ta question suivante pendant que l'autre parle. \textbf{Réagis} (\eng{Really?}, \eng{That sounds great!}, \eng{Oh, I see.}) \textbf{avant} de relancer.
    \item \textbf{Utilise le vocabulaire du thème :} \cle{human rights}, \cle{NGO}, \cle{certificate}, \cle{violation}, \cle{Universal Declaration}, \cle{participant}, \cle{organiser}.
    \item \textbf{Conclus explicitement :} \eng{Thanks for sharing!}, \eng{It was nice talking to you.}, \eng{Good luck for the next one!}
\end{enumerate}
\end{methode}

\begin{exemplebox}[Dialogue modèle commenté (Extrait)]
\textbf{B :} Hi! You look happy. Did something good happen? \emph{(Ouverture naturelle, pas « Question 1 » direct.)}\\
\textbf{A :} Yes, actually. I attended a workshop on human rights in Douala.\\
\textbf{B :} \textbf{Really? That sounds interesting.} \emph{(Réaction obligatoire.)} \textbf{Who organised it?} \emph{(Question Wh- précise.)}\\
\textbf{A :} An NGO called “Youth for Democracy”. It lasted two days.\\
\textbf{B :} \textbf{Two days?} \emph{(Écho / relance.)} \textbf{How many participants were there?}\\
\textbf{A :} We were about thirty. From different schools.\\
\textbf{B :} \textbf{Thirty!} \emph{(Réaction.)} \textbf{And... did you get a certificate?}\\
\textbf{A :} Yes, I did. I found it very useful. We studied the Universal Declaration.\\
\textbf{B :} \textbf{That's great.} \emph{(Validation.)} \textbf{Thanks for telling me. I might join next time.} \emph{(Conclusion + projection future.)}
\end{exemplebox}

\begin{exoresolu}[Analyse d'une production élève — Corriger les erreurs d'interaction]
\textbf{Énoncé :} Voici un extrait d'un élève B. Identifiez 3 défauts d'interaction et proposez une version corrigée.
\begin{quote}
\textbf{B :} What was the workshop about? \\
\textbf{A :} Human rights in Douala. \\
\textbf{B :} Where was it? \\
\textbf{A :} In Douala, at the city hall. \\
\textbf{B :} How long? \\
\textbf{A :} Two days. \\
\textbf{B :} How many people? \\
\textbf{A :} Thirty. \\
\textbf{B :} OK. Bye.
\end{quote}

\tcblower
\textbf{Solution détaillée :}
\begin{enumerate}
    \item \textbf{Aucune réaction (backchanneling) :} L'élève B enchaîne les questions comme un interrogatoire de police. \textbf{Correction :} Insérer \eng{Really? / I see. / Interesting.} après chaque réponse.
    \item \textbf{Questions trop courtes / télégraphiques :} \eng{Where was it?} au lieu de \eng{Where did it take place?} ou \eng{Where was it held?}. \textbf{Correction :} Phrases complètes, formes variées (passif \eng{was it held}, préposition \eng{take place}).
    \item \textbf{Conclusion absente / brutale :} \eng{OK. Bye.} est impoli. \textbf{Correction :} \eng{Thanks for the info. It sounds really useful. See you later!}
\end{enumerate}
\textbf{Version corrigée :}
\eng{B: What was the workshop about? / A: Human rights in Douala. / B: \textbf{Really? That's a crucial topic.} Where was it held? / A: In Douala, at the city hall. / B: \textbf{I see.} How long did it last? / A: Two days. / B: \textbf{Two days, okay.} And how many participants were there? / A: Thirty. / B: \textbf{Quite a group!} Thanks for sharing. See you!}
\end{exoresolu}

\courssub{Jeu de rôle 2 : Semi-guidé — « Interview pour le journal du lycée : Formation ELECAM »}

Ici, les \textbf{rôles sont définis} (Journaliste / Participant) et l'\textbf{objectif est fixé} (article pour le journal), mais le \textbf{contenu est libre}. Vous devez mobiliser votre \textbf{connaissance du monde} (culture générale sur ELECAM, élections au Cameroun) et votre \textbf{lexique personnel}.

\begin{experience}[Mise en place — Jeu de rôle 2]
\begin{itemize}
    \item \textbf{Contexte :} Le journal scolaire \eng{The School Echo} prépare un numéro spécial « Citoyenneté ».
    \item \textbf{Rôle A (Journaliste) :} Tu dois écrire un article. Tu as 2 min pour préparer 6-7 questions couvrant : motivation, déroulement, apprentissage clé, anecdote, conseil aux autres élèves.
    \item \textbf{Rôle B (Participant ELECAM) :} Tu as suivi une formation de 3 jours à Yaoundé (bureau ELECAM) sur le processus électoral : inscription, vote, dépouillement, rôle des observateurs. Tu as simulé un vote. Prépare tes idées clés (vocabulaire : \cle{ballot box}, \cle{polling station}, \cle{observer}, \cle{turnout}, \cle{civic duty}).
    \item \textbf{Déroulement :} 3 min d'interview. Le journaliste prend des notes. Inversion des rôles possible.
    \item \textbf{Production écrite post-tâche (optionnel) :} Rédiger le lead (chapeau) de l'article en 3 phrases.
\end{itemize}
\end{experience}

\begin{center}
\begin{tabularx}{\linewidth}{|X|X|X|}
\hline
\rowcolor{popblueL} \textbf{Fonction communicative} & \textbf{Exponent (Journaliste)} & \textbf{Exponent (Participant)} \\
\hline
Demander la motivation & \eng{What motivated you to apply?} & \eng{I wanted to understand how elections work...} \\
\hline
Demander le déroulement & \eng{Walk me through the three days.} & \eng{Day 1 was theory... Day 2, we visited a polling station...} \\
\hline
Demander l'apprentissage clé & \eng{What was the main takeaway?} & \eng{The importance of transparency and the role of observers.} \\
\hline
Demander une anecdote & \eng{Any memorable moment?} & \eng{The simulation! I was the presiding officer. It was stressful but fun.} \\
\hline
Demander un conseil & \eng{Would you recommend it?} & \eng{Absolutely. Every young Cameroonian should know their civic rights.} \\
\hline
Conclure l'interview & \eng{Thank you for your time.} & \eng{My pleasure. Good luck with the article!} \\
\hline
\end{tabularx}
\end{center}

\begin{attention}[Piège : Le monologue du participant]
En tant que Participant (B), ne récitez pas un exposé de 2 minutes. \textbf{Attendez la question.} Si le journaliste ne pose pas la « bonne » question, \textbf{orientez-le} par une relance : \eng{That reminds me of the simulation we did...} ou \eng{What I found most surprising was...}. L'examinateur note l'\textbf{interaction}, pas le monologue.
\end{attention}

\courssub{Jeu de rôle 3 : Libre — « Deux amis comparent leurs expériences civiques »}

Situation totalement ouverte. Aucune fiche imposée. Vous devez \textbf{négocier le sens}, \textbf{comparer}, \textbf{nuancer} et \textbf{conclure ensemble}. C'est le niveau le plus haut de l'escalier (cf. figure 1).

\begin{textebox}[Consigne — Jeu de rôle 3]
\textbf{Situation :} Vous êtes deux amis (élèves en Terminale). L'un a fait un \textbf{workshop ONG « Droits de l'enfant »} (3 jours, Buea, jeux de rôle, plaidoyer). L'autre a suivi une \textbf{formation « Jeunes leaders » au Conseil de la Jeunesse} (week-end, Yaoundé, discours public, gestion de projet).
\textbf{Objectif :} Échanger sur vos expériences, \textbf{comparer} les apports (connaissances, compétences, ambiance), et \textbf{décider ensemble} laquelle vous recommanderiez à un ami qui veut s'engager, en justifiant votre choix.
\textbf{Durée :} 4 à 5 minutes.
\end{textebox}

\begin{aretenir}[Structures pour COMPARER et NUANCER (Niveau Terminale)]
\begin{center}
\begin{tabular}{|l|l|l|}
\hline
\rowcolor{popblueL} \textbf{Fonction} & \textbf{Structure anglaise} & \textbf{Équivalent français} \\
\hline
Similitude & \eng{Both experiences were...} / \eng{Similarly, ...} & Les deux expériences étaient... / De même, ... \\
\hline
Différence & \eng{While the workshop focused on..., the training...} & Tandis que le workshop portait sur..., la formation... \\
\hline
Contraste fort & \eng{Unlike the workshop, the training allowed me to...} & Contrairement au workshop, la formation m'a permis de... \\
\hline
Préférence argumentée & \eng{I'd recommend the training because...} / \eng{The workshop was more hands-on.} & Je recommanderais la formation car... / Le workshop était plus pratique. \\
\hline
Nuance & \eng{On the one hand... On the other hand...} / \eng{However,...} / \eng{That said,...} & D'un côté... De l'autre... / Cependant,... / Cela dit,... \\
\hline
\end{tabular}
\end{center}
\end{aretenir}

\begin{exemplebox}[Extrait de dialogue libre — Comparaison]
\textbf{A :} ...so the NGO workshop was great for \textbf{legal knowledge} — the UN Convention on the Rights of the Child. But \textbf{it was mostly listening to experts}.\\
\textbf{B :} \textbf{Interesting.} \textbf{In contrast}, the Youth Council training was \textbf{skills-based}. We had to \textbf{deliver a speech} and \textbf{write a project proposal}.\\
\textbf{A :} \textbf{I see.} \textbf{So, one is theoretical, the other practical.} \eng{On the one hand}, you know your rights; \eng{on the other hand}, you can act.\\
\textbf{B :} \textbf{Exactly.} \textbf{If I had to choose for a friend who wants to start an association...} \eng{I'd say do the Youth Council training first.} \textbf{You get the tools to act.} \eng{What do you think?}\\
\textbf{A :} \textbf{I agree.} \textbf{But ideally, do both!} \eng{Knowing your rights gives you the legitimacy; the skills give you the power.} \eng{Shall we suggest a combined programme to the principal?}\\
\textbf{B :} \textbf{Great idea! Let's draft a proposal together.}
\end{exemplebox}

\courssub{Gestion de l'oral : hésitations, réactions et fluidité}

L'oral n'est pas un texte écrit lu à voix haute. La \textbf{fluidité} (fluency) chez un locuteur non-natif, c'est la capacité à \textbf{gagner du temps pour penser} en anglais, sans casser le flux.

\begin{attention}[Le piège « Euh » — Utilisez des fillers anglais !]
\begin{itemize}
    \item \textbf{Interdit :} « Euh... », « Ben... », « Heu... » (Marqueurs francophones, pénalisés).
    \item \textbf{Adoptez :} \eng{Well...} (neutre, très fréquent), \eng{Let me see...} / \eng{Let me think...} (recherche d'info), \eng{You know...} (recherche de mot / vérification compréhension), \eng{I mean...} (auto-correction / précision), \eng{How shall I put it...} (recherche formulation précise), \eng{Actually...} (correction / précision / contraste).
\end{itemize}
\end{attention}

\begin{exemplebox}[Fillers en contexte — Traduction]
\begin{itemize}
    \item \eng{Well... it was in July, I think.} = « Voyons... c'était en juillet, je crois. »
    \item \eng{Let me think... there were about thirty participants.} = « Laisse-moi réfléchir... il y avait environ trente participants. »
    \item \eng{You know... the Universal Declaration thing.} = « Tu sais... le truc de la Déclaration universelle. »
    \item \eng{I mean, it wasn't boring, just very theoretical.} = « Je veux dire, ce n'était pas ennuyeux, juste très théorique. »
    \item \eng{How shall I put it... it was life-changing.} = « Comment dire... ça a changé ma vie. »
\end{itemize}
\end{exemplebox}

\begin{savaistu}[Le secret de l'examinateur : INTERACTION > GRAMMAIRE PARFAIT]
À l'épreuve orale du Baccalauréat, la grille d'évaluation pondère souvent \textbf{l'Interaction (Réagir, Relancer, Prendre la parole, Céder la parole)} autant, voire plus, que la \textbf{Grammaire} ou le \textbf{Vocabulaire} pris séparément.
\begin{itemize}
    \item Un élève qui fait 2 fautes de grammaire mais qui \textbf{regarde son partenaire}, \textbf{sourit}, \textbf{réagit} (\eng{Really?}), \textbf{relance} (\eng{And then?}) et \textbf{conclut} poliment aura une \textbf{meilleure note} qu'un élève qui récite un texte parfait sans regarder personne, sans réagir, et qui s'arrête net à la fin.
    \item \textbf{Conseil :} Entraînez-vous à \textbf{regarder votre partenaire} (ou l'examinateur) dans les yeux 80\% du temps. Votre carnet de notes ne doit servir qu'à jeter un coup d'œil rapide (mots-clés), pas à être lu.
\end{itemize}
\end{savaistu}

\courssub{Grille d'auto-évaluation : Suis-je prêt pour le Bac ?}

Utilisez cette grille \textbf{après chaque entraînement} (enregistrez-vous si possible). Soyez honnête : l'objectif n'est pas d'avoir « Oui » partout tout de suite, mais d'identifier \textbf{un point précis à améliorer} pour la prochaine fois.

\begin{center}
\begin{tabularx}{\linewidth}{|X|c|c|c|}
\hline
\rowcolor{popblueL} \textbf{Critère (Compétence interactionnelle)} & \textbf{Oui} & \textbf{En partie} & \textbf{Non} \\
\hline
1. \textbf{Questions variées :} J'ai utilisé Wh- (What, Where, How long, Why) ET Yes/No. Pas seulement « And...? » & & & \\
\hline
2. \textbf{Réactions authentiques :} J'ai réagi à \textbf{chaque} réponse partenaire (Really? / That's great / I see / Oh!). & & & \\
\hline
3. \textbf{Relances / Follow-up :} J'ai posé au moins 2 questions \textbf{non prévues} nées de l'écoute (ex. \eng{Why did you like that?}). & & & \\
\hline
4. \textbf{Vocabulaire thématique :} J'ai utilisé 5+ mots du thème (workshop, certificate, NGO, human rights, civic duty, polling station, turnout...). & & & \\
\hline
5. \textbf{Gestion hésitations :} J'ai utilisé \textbf{fillers anglais} (Well, Let me see) et \textbf{pas} « euh ». & & & \\
\hline
6. \textbf{Structures de comparaison :} (Jeu 3) J'ai utilisé \eng{While / Unlike / On the one hand / I'd recommend}. & & & \\
\hline
7. \textbf{Conclusion polie :} J'ai clos l'échange explicitement (\eng{Thanks for sharing / Good luck / See you}). & & & \\
\hline
8. \textbf{Prononciation / Intonation :} J'ai été compréhensible ; accentuation tonique correcte ; intonation montante (questions) / descendante (affirmations). & & & \\
\hline
\end{tabularx}
\end{center}

\begin{methode}[Plan d'action si « En partie » ou « Non »]
\begin{enumerate}
    \item Ciblez \textbf{UN seul critère} à travailler (ex. Critère 2 : Réactions).
    \item Écrivez \textbf{3 formules types} sur une fiche Bristol (ex. \eng{Really? / That's interesting! / I didn't know that.}).
    \item Faites un \textbf{mini-jeu de rôle de 2 min} en ne vous concentrant \textbf{que sur ce critère} (le contenu n'importe pas, seul le réflexe de réaction compte).
    \item Re-testez la grille complète la semaine suivante.
\end{enumerate}
\end{methode}

\courssub{Complément Bac : Décrire une expérience personnelle en 5-6 phrases fluides (Compétence transverse)}

\emph{Signalé comme « Complément utile au Bac » :} Cette compétence (raconter une expérience vécue : \cle{Past Simple} + \cle{Present Perfect} + \cle{adverbes de temps} + \cle{connecteurs logiques} + \cle{opinion}) est \textbf{réutilisable dans TOUS les sujets d'expression orale} (Description de photo, Récit, Exposé, Jeu de rôle). Maîtrisez ce « module » une fois pour toutes.

\begin{methode}[La méthode « STAR » adaptée à l'oral (5-6 phrases)]
Structurez votre récit comme une mini-histoire cohérente :
\begin{enumerate}
    \item \textbf{Situation (Contexte) :} \eng{Last year / In 2023, I took part in a...} (Past Simple \cle{date passée}).
    \item \textbf{Task (Objectif) :} \eng{The aim was to learn about... / We had to...} (Infinitif de but / \cle{had to}).
    \item \textbf{Action (Ce que VOUS avez fait) :} \eng{I worked in groups / I presented a project / We visited...} (Past Simple \textbf{verbes d'action}, \textbf{je} sujet).
    \item \textbf{Result (Résultat concret) :} \eng{We received a certificate / I improved my public speaking / We wrote a report.} (Past Simple / Present Perfect si conséquence actuelle).
    \item \textbf{Reflection (Bilan personnel / Opinion) :} \eng{It was really useful because... / I now feel more confident about... / I'd recommend it to anyone who...} (\cle{was/were} + adj + \cle{because} / \cle{Present Perfect} + \cle{now}).
    \item \textbf{(Optionnel) Opening :} \eng{Since then, I have joined... / I'm planning to...} (Present Perfect / \cle{be planning to} = projection).
\end{enumerate}
\end{methode}

\begin{exemplebox}[Modèle prêt à l'emploi — « Mon expérience ELECAM »]
\eng{In March 2024, I participated in a three-day training session organised by ELECAM at their Yaoundé headquarters. \textbf{(Situation)} The goal was to train young observers for the upcoming elections. \textbf{(Task)} I learnt how to fill out observation checklists and how to report irregularities at polling stations. \textbf{(Action)} We also ran a full voting simulation where I acted as a polling officer. \textbf{(Action)} At the end, I received an official observer badge and a certificate. \textbf{(Result)} \textbf{It was an eye-opening experience} because I realised how important transparency is for democracy. \textbf{(Reflection)} \textbf{Since then, I have encouraged all my friends to register to vote.} \textbf{(Opening / Present Perfect)}}
\end{exemplebox}

\begin{exercice}[Entraînement « STAR » — À vous de jouer !]
Préparez votre propre description « STAR » (5-6 phrases) pour l'une de ces situations. Enregistrez-vous. Visez 45-60 secondes.
\begin{enumerate}
    \item Le workshop « Droits de l'enfant » à Buea (Jeu de rôle 3, Rôle A).
    \item La formation « Jeunes leaders » au Conseil de la Jeunesse (Jeu de rôle 3, Rôle B).
    \item Une action de nettoyage / plantation d'arbres avec votre club environnement.
    \item Votre stage d'observation en entreprise / administration.
\end{enumerate}
\textbf{Check-list avant enregistrement :} \square Past Simple (dates passées) \square Connecteurs (\eng{First, Then, Finally, Because}) \square Au moins 3 mots « thème citoyenneté » \square 1 filler anglais naturel \square Intonation descendante sur les affirmations.
\end{exercice}

\begin{corrige}[Exercice « STAR » — Exemple de production attendue (Situation 3 : Club environnement)]
\eng{Last June, I joined a tree-planting campaign organised by our school's Green Club in the Mfoundi neighbourhood. \textbf{(Situation)} We aimed to plant 200 trees to fight erosion and provide shade. \textbf{(Task)} I was in charge of the logistics team: distributing seedlings and tools to the volunteers. \textbf{(Action)} We worked all Saturday morning and exceeded the target with 250 trees planted. \textbf{(Result)} \textbf{I felt really proud} because we transformed a bare space into a future green lung for the community. \textbf{(Reflection)} \textbf{Since then, I have become the club's vice-president} and I'm planning a clean-up walk for next month. \textbf{(Opening)}}
\end{corrige}

\begin{aretenir}[Récapitulatif final de la Leçon 40 — Speaking : Democracy Training]
\begin{itemize}
    \item \textbf{Vocabulaire clé :} \cle{workshop}, \cle{training}, \cle{certificate}, \cle{NGO}, \cle{human rights}, \cle{civic duty}, \cle{ELECAM}, \cle{polling station}, \cle{ballot box}, \cle{observer}, \cle{turnout}, \cle{advocacy}, \cle{empowerment}.
    \item \textbf{Fonctions :} Demander info (Wh-/Yes-No), Donner info (Past Simple/Present Perfect), Relancer (\eng{Tell me more...}), Réagir (\eng{Really?}), Comparer (\eng{While/Unlike}), Conclure (\eng{Thanks for sharing}).
    \item \textbf{Prononciation :} Accent tonique (\cle{workshop} « worc-chope », \cle{certificate} « ser-ti-fi-kate », \cle{democracy} « di-mo-cra-si »), Intonation (Montante Y/N, Descendante Wh-), Liaisons (\cle{did you} « di-djou »).
    \item \textbf{Stratégie orale :} Fillers anglais (\eng{Well, Let me see}), Contact visuel, Structure STAR pour récit, Grille auto-évaluation.
\end{itemize}
\emph{Bon courage pour vos entraînements et pour l'épreuve orale du Baccalauréat ! Votre voix compte.}
\end{aretenir}

\newpage

\coursec{Activité d'intégration}

\coursec{Lesson 40 — Speaking: Exchanges information on participating in training activities/courses on democracy}

\courssub{Mise en situation et vocabulaire thématique}

\begin{textebox}[CRTV Youth Programme — Call for young guests]
\textbf{CRTV Youth Programme — “Young Citizens Speak!”}

Do you live in Yaoundé, Douala, Buea or Bamenda? Have you ever taken part in a training course on democracy, elections or human rights — with ELECAM, an NGO, or your school civic club?

Our radio programme “Young Citizens Speak!” is looking for young guests for next Saturday's show. Come to our studio with a partner: one of you will be the journalist, the other will be the guest. You will record a three-minute interview in English.

The journalist must ask at least six different questions, react to the answers and close the interview politely. The guest must give full answers: where and when the training took place, what you learnt, and your opinion about it.

The best interviews will be broadcast on Saturday at 5 p.m. Don't miss this chance to make young voices heard!
\end{textebox}

\begin{definition}[Vocabulaire clé : Democracy and Civic Training]
\begin{itemize}
\item \textbf{training course} (n.) / workshop (n.) / seminar (n.) : stage de formation / atelier / séminaire.
\item \textbf{democracy} (n.) / \textbf{elections} (n. pl.) / \textbf{human rights} (n. pl.) : démocratie / élections / droits de l'homme.
\item \textbf{NGO} (n.) = \emph{Non-Governmental Organisation} : ONG.
\item \textbf{civic club} (n.) / \textbf{youth organisation} (n.) : club civique / organisation de jeunesse.
\item \textbf{ELECAM} (n.) = \emph{Elections Cameroon} : organe gestionnaire des élections au Cameroun.
\item \textbf{to register} (v.) / \textbf{registration} (n.) : s'inscrire / inscription (sur les listes électorales).
\item \textbf{ballot box} (n.) / \textbf{polling station} (n.) : urne / bureau de vote.
\item \textbf{to broadcast} (v.) / \textbf{broadcast} (n.) : diffuser / diffusion ; \emph{guest} (n.) : invité ; \emph{host} (n.) / \emph{presenter} (n.) : animateur ; \emph{listener} (n.) : auditeur.
\item \textbf{to take part in} / \textbf{to attend} / \textbf{to participate in} : participer à / assister à.
\item \textbf{to make one's voice heard} : faire entendre sa voix.
\end{itemize}
\end{definition}

\courssub{Demander et donner une information : former des questions}

\begin{propriete}[Questions en \cle{wh-} (questions ouvertes)]
Pour demander une information précise, on place le mot interrogatif (\cle{wh-}) en tête, suivi de l'auxiliaire (do/does/did/have/be/modal) + sujet + verbe (forme de base).
\begin{center}
\begin{tabularx}{\linewidth}{|X|X|X|}\hline
\rowcolor{popblueL} \textbf{Mot interrogatif} & \textbf{Structure} & \textbf{Exemple (contexte démocratie)} \\ \hline
\eng{Where} (Où) & \eng{Wh- + auxiliaire + sujet + V ?} & \eng{Where did the training take place?} \\ \hline
\eng{When} (Quand) & \eng{Wh- + auxiliaire + sujet + V ?} & \eng{When did you attend the workshop?} \\ \hline
\eng{What} (Quoi / Quel) & \eng{Wh- + auxiliaire + sujet + V ?} & \eng{What did you learn about voting?} \\ \hline
\eng{Who} (Qui - sujet) & \eng{Who + verbe + ... ?} & \eng{Who organised the seminar?} \\ \hline
\eng{How long} (Combien de temps) & \eng{Wh- + auxiliaire + sujet + V ?} & \eng{How long did the course last?} \\ \hline
\eng{How many} (Combien - dénombrable) & \eng{Wh- + auxiliaire + sujet + V ?} & \eng{How many young people attended?} \\ \hline
\end{tabularx}
\end{center}
\eng{Who} peut être sujet (\eng{Who called you?}) ou objet (\eng{Who did you meet?}). En anglais courant, \eng{who} remplace \eng{whom} à l'objet.
\end{propriete}

\begin{propriete}[Questions fermées \cle{Yes/No} (auxiliaire en tête)]
Pour une réponse par \eng{Yes} ou \eng{No}, l'auxiliaire précède le sujet. Pas de \cle{wh-}.
\begin{center}
\begin{tabular}{|l|l|l|}\hline
\rowcolor{popblueL} \textbf{Temps / Auxiliaire} & \textbf{Structure} & \textbf{Exemple} \\ \hline
\cle{Present Simple} & \eng{Do/Does + sujet + V ?} & \eng{Do you know ELECAM? / Does he vote?} \\ \hline
\cle{Past Simple} & \eng{Did + sujet + V (base) ?} & \eng{Did you enjoy the training?} \\ \hline
\cle{Present Perfect} & \eng{Have/Has + sujet + V (participe passé) ?} & \eng{Have you ever voted?} \\ \hline
\cle{Modaux (can, will, should...)} & \eng{Modal + sujet + V (base) ?} & \eng{Can young people change things?} \\ \hline
\cle{BE (présent / passé)} & \eng{Am/Is/Are/Was/Were + sujet + ... ?} & \eng{Was the workshop interesting?} \\ \hline
\end{tabular}
\end{center}
\emph{Règle d'or :} En question, le verbe principal ne porte jamais la marque du temps (pas de \eng{-s}, pas de \eng{-ed}) ; c'est l'auxiliaire qui la porte.
\end{propriete}

\courssub{Choisir le bon temps : \cle{Past Simple} vs \cle{Present Perfect}}

\begin{propriete}[Past Simple : faits datés, terminés, coupés du présent]
\begin{itemize}
\item \textbf{Emploi :} Action terminée dans le passé, souvent avec un repère de temps \textbf{terminé} (\eng{yesterday, last week/month/year, in 2023, two days ago, when I was...}).
\item \textbf{Forme :} \eng{V+ed} (réguliers) ou 2\textsuperscript{e} colonne (irréguliers). Négation : \eng{did not (didn't) + V base}. Question : \eng{Did + sujet + V base ?}
\end{itemize}
\eng{I attended a workshop last month.} (\"at-tèn-did\") \\
\eng{We learnt how elections are organised.} (\"lèrnt\") \\
\eng{The training took place in Buea.}
\end{propriete}

\begin{propriete}[Present Perfect : lien avec le présent, expérience de vie, résultat actuel]
\begin{itemize}
\item \textbf{Emploi :} 1. Expérience de vie (jamais / déjà / déjà) sans date précise. 2. Action passée avec conséquence visible maintenant. 3. Durée depuis le passé (\eng{for/since}) — non requis ici.
\item \textbf{Forme :} \eng{have/has + V-en (participe passé)}. Négation : \eng{have not (haven't) + V-en}. Question : \eng{Have/Has + sujet + V-en ?}
\item \textbf{Marqueurs fréquents :} \eng{ever, never, already, yet, just, recently, so far, up to now}.
\end{itemize}
\eng{I have never taken part in an election.} \\
\eng{She has already registered on the electoral list.} \\
\eng{We have learnt a lot about our rights.}
\end{propriete}

\begin{attention}[Piège majeur : Repère de temps terminé = Past Simple obligatoire]
Ne dites jamais : \eng{*I have attended the workshop last month.} \\
\eng{Last month} est un repère terminé $\rightarrow$ \eng{I attended the workshop last month.}
\par
Autre piège : \eng{*Where it took place?} (ordre affirmatif). Question = \textbf{inversion} : \eng{Where did it take place?}
\end{attention}

\courssub{Relancer, réagir, ouvrir et conclure : la boîte à outils fonctionnelle}

\begin{aretenir}[Formules de l'interview radiophonique]
\textbf{1. Ouvrir (Journaliste)} \\
\eng{Good afternoon, dear listeners, and welcome to “Young Citizens Speak!”} \\
\eng{Today, our guest is [Name], a student from [School]. Welcome, [Name]!} \\

\textbf{2. Demander (Journaliste)} \\
\eng{Where did the training take place?} \quad \eng{When did you attend it?} \\
\eng{What did you learn there?} \quad \eng{How long did it last?} \\
\eng{Did you enjoy the course?} \quad \eng{Have you ever taken part in an election?} \\
\eng{Who organised the workshop?} \quad \eng{How many participants were there?} \\

\textbf{3. Informer / Répondre (Invité)} \\
\eng{I attended a two-day workshop organised by ELECAM at their headquarters in Yaoundé.} \\
\eng{It took place last month, on a Saturday and a Sunday.} \\
\eng{We learnt how votes are counted and why registration is crucial.} \\
\eng{I have met many committed young people.} \\
\eng{In my opinion, every young citizen should register and vote.} \\

\textbf{4. Réagir et relancer (Journaliste — \emph{backchannelling})} \\
\eng{Really? That sounds great / interesting / important!} \\
\eng{How interesting! / I see. / That's impressive.} \\
\eng{Can you tell me more about that?} \quad \eng{And what happened next?} \\
\eng{Why do you think that is important?} \\

\textbf{5. Conclure (Journaliste)} \\
\eng{Thank you very much for sharing your experience with us, [Name].} \\
\eng{That was very interesting / inspiring.} \\
\eng{Thank you, dear listeners. See you next Saturday for another episode!} \\
\eng{Goodbye!}
\end{aretenir}

\courssub{Prononciation, accentuation et intonation}

\begin{propriete}[Intonation des questions]
\begin{itemize}
\item \textbf{Questions en \cle{wh-} (Where, What, When, How...)} : \textbf{intonation descendante} à la fin (voix qui tombe). \eng{Where did it take place? \textsearrow}
\item \textbf{Questions fermées \cle{Yes/No} (Did, Have, Do, Is...)} : \textbf{intonation montante} à la fin (voix qui monte). \eng{Did you enjoy it? \textnearrow}
\item \textbf{Listes / Choix} : montante sur les éléments non finaux, descendante sur le dernier. \eng{Was it in Yaoundé \textnearrow, Douala \textnearrow or Buea \textsearrow?}
\end{itemize}
\end{propriete}

\begin{propriete}[Prononciation du \cle{-ed} au Past Simple (réguliers)]
Trois réalisations selon le son final du verbe (hors \eng{-e} muet) :
\begin{center}
\begin{tabular}{|c|c|l|}\hline
\rowcolor{popblueL} \textbf{Son final du radical} & \textbf{Prononciation -ed} & \textbf{Exemples (contexte)} \\ \hline
Sourds : /p/, /k/, /f/, /s/, /$\theta$/, /t$\int$/ & \textbf{/t/} (\"t\" sec) & \eng{worked} (\"wèrkt\"), \eng{asked} (\"æskt\"), \eng{finished} (\"finicht\") \\ \hline
Sonores : voyelles, /b/, /d/, /g/, /v/, /z/, /d$\int$/, /m/, /n/, /ŋ/, /l/, /r/ & \textbf{/d/} (\"d\" vibrant) & \eng{attended} (\"atèn-did\"), \eng{organised} (\"organ-aïzd\"), \eng{learnt} (\"lèrnt\" irrégulier), \eng{changed} (\"tchéindjd\") \\ \hline
/t/ ou /d/ & \textbf{/id/} (\"id\" syllabe à part) & \eng{wanted} (\"wont-id\"), \eng{needed} (\"ni-did\"), \eng{decided} (\"di-saï-did\") \\ \hline
\end{tabular}
\end{center}
\emph{Astuce :} Posez deux doigts sur la gorge. Si ça vibre en disant le radical (ex. \eng{organise} $\rightarrow$ vibration sur /z/), le \eng{-ed} se prononce /d/. Si pas de vibration (ex. \eng{ask} $\rightarrow$ /k/), c'est /t/.
\end{propriete}

\begin{aretenir}[Accentuation de phrase (Stress-timing)]
Les \textbf{mots de contenu} (noms, verbes principaux, adjectifs, adverbes, interrogatifs, \eng{not}) sont accentués (plus forts, plus longs, plus aigus). Les \textbf{mots grammaticaux} (articles, auxiliaires, prépositions, pronoms, conjonctions) sont faibles, souvent réduits (\eng{schwa /ə/}).
\par
\eng{\textbf{WHERE} did the \textbf{TRAINING} \textbf{TAKE} \textbf{PLACE}?} (Wh- : accent sur \eng{Where}, \eng{training}, \eng{take}, \eng{place} ; \eng{did}, \eng{the} faibles).
\par
\eng{\textbf{DID} you \textbf{ENJOY} the \textbf{COURSE}?} (Yes/No : accent sur \eng{Did}, \eng{enjoy}, \eng{course} ; \eng{you}, \eng{the} faibles).
\end{aretenir}

\courssub{Mise en œuvre : Situation d'intégration — Jeu de rôle « Young Citizens Speak! »}

\begin{experience}[Consignes de l'activité d'intégration (2 h)]
\textbf{Objectif final :} En binôme, préparer et jouer une interview radiophonique de 3 minutes en anglais sur une formation à la démocratie, devant la classe ou en enregistrement audio.

\textbf{Étape 1 — Compréhension du déclencheur (10 min).} Lisez l'annonce CRTV (texte ci-dessus). Répondez en anglais :
\begin{enumerate}
\item Who is looking for young guests?
\item What is the topic of the interview?
\item How long must the interview last?
\item When will the best interviews be broadcast?
\end{enumerate}
\textit{Corrigé attendu :} 1. The CRTV Youth Programme “Young Citizens Speak!”. 2. Participating in a training course on democracy, elections or human rights. 3. Three minutes. 4. Next Saturday at 5 p.m.

\textbf{Étape 2 — Choix des rôles et du contexte (10 min).}
\begin{itemize}
\item Décidez : \textbf{Journaliste} (anime, questionne, réagit, conclut) / \textbf{Invité} (a suivi la formation, informe, donne son avis).
\item L'invité choisit un cadre réaliste : \eng{ELECAM workshop on voting} (Yaoundé/Douala), \eng{NGO seminar on human rights} (Buea/Bamenda), \eng{School civic club training}.
\item L'invité note 3 infos clés : \textbf{Lieu}, \textbf{Date/Durée}, \textbf{Contenu principal} (3 points), \textbf{Avis personnel}.
\end{itemize}

\textbf{Étape 3 — Préparation linguistique guidée (25 min).}
\begin{itemize}
\item \textbf{Journaliste :} Rédigez \textbf{6 questions minimum} : 3 \cle{wh-} + 3 \cle{Yes/No}. Vérifiez l'inversion (Aux + S + V). Préparez 3 formules de réaction/relance.
\item \textbf{Invité :} Rédigez des réponses complètes (phrases complètes). Utilisez \cle{Past Simple} pour les faits datés (\eng{It took place... / We learnt...}) et \cle{Present Perfect} pour l'expérience/résultat (\eng{I have met... / I have understood...}). Préparez une phrase d'opinion avec \eng{In my opinion / I think / I believe + should/must}.
\end{itemize}

\textbf{Étape 4 — Répétition avec focus phonologique (25 min).}
Jouez l'interview à voix basse. Journaliste : travaillez l'intonation \textsearrow (Wh-) / \textnearrow (Yes/No). Invité : soignez la prononciation des \eng{-ed} (\eng{attended, organised, learnt, registered}). Chronométrez : viser 3 minutes.

\textbf{Étape 5 — Performance et évaluation par les pairs (40 min).}
Chaque binôme passe devant la classe (ou enregistre). Les camarades remplissent la grille d'écoute ci-dessous. L'enseignant note avec la grille détaillée (critères : interaction, grammaire, vocabulaire, phonologie, cohérence).

\begin{center}
\begin{tabularx}{\linewidth}{|X|X|}\hline
\rowcolor{popblueL} \textbf{Critère d'écoute (Grille pair)} & \textbf{Oui / Non} \\ \hline
At least six questions were asked (wh- and yes/no). & \\ \hline
The journalist reacted to the answers (backchannelling). & \\ \hline
The guest used the present perfect and the past simple correctly. & \\ \hline
The interview was opened and closed politely (radio style). & \\ \hline
The intonation was correct (falling wh- / rising yes-no) and speech clear. & \\ \hline
\end{tabularx}
\end{center}
\end{experience}

\courssub{Production modèle commentée}

\begin{exoresolu}[The Youth Civic Radio Interview — Modèle complet (3 min)]
Énoncé : Produisez l'interview complète entre la journaliste Amina et l'invité Samuel, élève au Lycée de Nkolbisson (Yaoundé), ayant suivi un atelier ELECAM.
\tcblower
\textbf{Production modèle :}

\eng{Amina: Good afternoon, dear listeners, and welcome to “Young Citizens Speak!” Today, our guest is Samuel, a student from Lycée de Nkolbisson. Good afternoon, Samuel!}

\eng{Samuel: Good afternoon, Amina. Thank you for inviting me.}

\eng{Amina: Samuel, you have recently taken part in a training course on democracy. Where did it take place?}

\eng{Samuel: It took place at the ELECAM headquarters in Yaoundé, last month.}

\eng{Amina: Really? And how long did the training last?}

\eng{Samuel: It lasted two days, a Saturday and a Sunday.}

\eng{Amina: That sounds interesting! What did you learn there?}

\eng{Samuel: We learnt how elections are organised, how votes are counted, and why every citizen should register on the electoral list.}

\eng{Amina: I see. Did you meet other young people there?}

\eng{Samuel: Yes, I did. There were students from many schools in Yaoundé. I have made new friends and I have understood that voting is a real responsibility.}

\eng{Amina: Wonderful! Have you ever voted in a school election yourself?}

\eng{Samuel: Yes, I have. I voted for our class delegate last year, and now I want to encourage all my friends to vote too.}

\eng{Amina: That is a great message! Thank you very much, Samuel, for sharing your experience with us.}

\eng{Samuel: You are welcome. It was a pleasure.}

\eng{Amina: And thank you, dear listeners. See you next Saturday for another episode of “Young Citizens Speak!”}

\textbf{Commentaire linguistique détaillé :}
\begin{enumerate}
\item \textbf{Opening} : Formule radiophonique authentique (\eng{dear listeners, welcome to...}), présentation de l'invité (nom, école).
\item \textbf{Questions (6+)} : Alternance \cle{wh-} (\eng{Where, How long, What}) $\rightarrow$ intonation \textsearrow ; et \cle{Yes/No} (\eng{Did you meet..., Have you ever voted...}) $\rightarrow$ intonation \textnearrow. Inversion auxiliaire-sujet respectée partout.
\item \textbf{Temps verbaux (Invité)} : \cle{Past Simple} pour faits datés (\eng{took place, lasted, learnt, met, voted}) — repères \eng{last month, last year, two days}. \cle{Present Perfect} pour expérience vie/résultat présent (\eng{have taken part, have made, have understood, have ever voted}).
\item \textbf{Réactions (Journaliste)} : \eng{Really? That sounds interesting! I see. Wonderful!} (Backchannelling naturel).
\item \textbf{Relance} : \eng{And how long...?} (suite logique) ; \eng{Have you ever voted...?} (ouverture sur expérience personnelle).
\item \textbf{Closing} : Double clôture standard radio : remerciement invité (\eng{Thank you... for sharing...}) + salut auditeurs (\eng{Thank you, dear listeners. See you next Saturday...}).
\item \textbf{Prononciation ciblée} : \eng{attended} (/t/), \eng{organised} (/d/), \eng{learnt} (/nt/), \eng{registered} (/d/), \eng{encourage} (accent sur 2\textsuperscript{e} syllabe /ɪnˈkʌrɪdʒ/).
\end{enumerate}
\end{exoresolu}

\courssub{Bilan de la leçon et prolongements}

\begin{aretenir}[Ce qu'il faut retenir de la Leçon 40]
\begin{itemize}
\item \textbf{Savoir-faire speaking :} Mener un échange d'information structuré (interview) : ouvrir $\rightarrow$ questionner (Wh- + Yes/No) $\rightarrow$ réagir/relancer $\rightarrow$ conclure.
\item \textbf{Grammaire clé :} Distinction stricte \cle{Past Simple} (passé daté, terminé) vs \cle{Present Perfect} (expérience, lien présent). Formation correcte des questions (inversion auxiliaire).
\item \textbf{Phonologie clé :} Intonation descendante (Wh-) / montante (Yes/No). Prononciation correcte des \eng{-ed} (/t/, /d/, /id/).
\item \textbf{Lexique thème :} Democracy, elections, human rights, civic engagement, training, workshop, ELECAM, NGO, register, vote, ballot, broadcast.
\end{itemize}
\end{aretenir}

\begin{savaistu}[Citoyenneté jeune au Cameroun : la Semaine de la Jeunesse et les clubs civiques]
Chaque année, la \textbf{Semaine nationale de la Jeunesse} (février) précède la Fête de la Jeunesse (11 février). C'est l'occasion pour les lycées (Yaoundé, Douala, Buea, Bamenda, Maroua...) d'organiser des \textbf{clubs civiques}, des simulations d'élections (vote des délégués), des ateliers avec \textbf{ELECAM} ou des ONG comme \textbf{Plan International}, \textbf{UNICEF Cameroun} ou \textbf{Transparency International}. Participer à ces formations permet d'obtenir des attestations valorisantes pour le dossier d'orientation post-bac (universités, grandes écoles, concours administratifs). L'anglais est la langue de travail de nombreuses ONG internationales : maîtriser ce vocabulaire et ces échanges oraux est un atout professionnel réel.
\end{savaistu}

\begin{methode}[Pour réussir l'épreuve orale du Baccalauréat (Interaction)]
\begin{enumerate}
\item \textbf{Écoutez la consigne / le partenaire} : ne récitez pas un texte appris par cœur ; répondez à ce qui est dit.
\item \textbf{Structurez} : Ouvrez $\rightarrow$ Développez (arguments + exemples) $\rightarrow$ Concluez.
\item \textbf{Variez la grammaire} : Montrez que vous maîtrisez au moins 3 temps (Present, Past, Perfect, Future/Modaux) et 2 types de questions.
\item \textbf{Soignez la prononciation} : Intonation questions, accentuation mots forts, \eng{-ed} distincts. Parlez assez fort, regardez l'examinateur/partenaire.
\item \textbf{Gérez l'erreur} : Si vous bloquez, utilisez \eng{Let me think... / What I mean is... / In other words...} pour gagner du temps sans silence long.
\end{enumerate}
\end{methode}

\newpage

\coursec{Méthodes et exercices résolus}

\coursec{Lesson 40 — Speaking: Exchanges information on participating in training activities/courses on democracy}

\courssub{Mise en situation et vocabulaire du thème}

\begin{textebox}{Script d'écoute : Une formation à la démocratie à Bamenda}
\eng{Ngozi:} “Hello Paul! You look very happy today.”
\eng{Paul:} “Hi Ngozi! I am. I just came back from a two-day training course on democracy in Bamenda.”
\eng{Ngozi:} “That sounds interesting. Who organised it?”
\eng{Paul:} “It was organised by a local NGO called ‘Youth for Change’, with support from the town council.”
\eng{Ngozi:} “Where did it take place exactly?”
\eng{Paul:} “At the community hall, near the main market. About sixty young people attended.”
\eng{Ngozi:} “What topics did you cover?”
\eng{Paul:} “We discussed human rights, the electoral process, civic duties and how to write a petition to local authorities.”
\eng{Ngozi:} “Did you find it useful?”
\eng{Paul:} “Absolutely! The trainers were excellent and the debates were lively. I learnt a lot about my rights.”
\eng{Ngozi:} “Will there be another session?”
\eng{Paul:} “Yes, next month there is a workshop on leadership. I plan to register.”
\eng{Ngozi:} “Great! Thanks for sharing. I might join the next one.”
\end{textebox}

\begin{definition}[Vocabulaire clé : Training Courses on Democracy]
\begin{center}
\begin{tabularx}{\linewidth}{|X|X|X|}
\hline
\rowcolor{popblueL} \textbf{English} & \textbf{Français} & \textbf{Exemple / Contexte} \\
\hline
\eng{a training course} & une formation, un stage & \eng{attend a training course on democracy} \\
\eng{a workshop} & un atelier (pratique) & \eng{a workshop on human rights} \\
\eng{a seminar} & un séminaire & \eng{a seminar on civic engagement} \\
\eng{the organiser} & l'organisateur & \eng{The course was organised by an NGO.} \\
\eng{a participant / a trainee} & un participant / un stagiaire & \eng{Sixty participants attended.} \\
\eng{a trainer / a facilitator} & un formateur / un animateur & \eng{The facilitator explained the voting process.} \\
\eng{the venue} & le lieu (de l'événement) & \eng{The venue was the community hall.} \\
\eng{the duration} & la durée & \eng{The duration was two days.} \\
\eng{the agenda / the schedule} & l'ordre du jour / le programme & \eng{The agenda included elections and rights.} \\
\eng{a module} & un module & \eng{Module 1: Human Rights.} \\
\eng{civic duties} & les devoirs civiques & \eng{We learnt about our civic duties.} \\
\eng{voting rights} & le droit de vote & \eng{Voting rights are fundamental.} \\
\eng{a petition} & une pétition & \eng{write a petition to the mayor} \\
\eng{local authorities} & les autorités locales & \eng{the town council, the mayor} \\
\eng{an NGO (Non-Governmental Organisation)} & une ONG & \eng{a local NGO} \\
\eng{the town council} & la mairie / le conseil municipal & \eng{with the help of the town council} \\
\eng{a youth association} & une association de jeunes & \eng{members of a youth association} \\
\eng{a certificate} & un certificat, une attestation & \eng{receive a certificate of attendance} \\
\hline
\end{tabularx}
\end{center}
\end{definition}

\begin{popfigure}
\begin{tikzpicture}[mindmap, grow cyclic, every node/.style=concept, concept color=popblue!60,
    level 1/.append style={level distance=4.5cm, sibling angle=90},
    level 2/.append style={level distance=3cm, sibling angle=45}]
\node {Training Course}
    child[concept color=popgreen!60] {node {People}
        child {node {Organiser}}
        child {node {Trainer / Facilitator}}
        child {node {Participant / Trainee}}
    }
    child[concept color=poporange!60] {node {Logistics}
        child {node {Venue}}
        child {node {Duration / Schedule}}
        child {node {Agenda / Modules}}
    }
    child[concept color=poppurple!60] {node {Content}
        child {node {Human Rights}}
        child {node {Elections / Voting}}
        child {node {Civic Duties}}
        child {node {Petition Writing}}
    }
    child[concept color=poppink!60] {node {Outcome}
        child {node {Certificate}}
        child {node {Skills / Knowledge}}
        child {node {Civic Engagement}}
    };
\end{tikzpicture}
\legende{Carte mentale : lexique d'une formation à la démocratie}
\end{popfigure}

\courssub{Demander et donner une information}

\begin{methode}[Échanger des informations sur une formation à la démocratie]
Pour mener un échange oral sur la participation à un \emph{training course on democracy}, suis cette démarche en cinq étapes :
\begin{enumerate}
\item \textbf{Prépare ton sujet} : identifie les informations clés — le nom de la formation, l'organisme (\eng{an NGO, the town council, a youth association}), le lieu (\eng{at the community hall in Bamenda}), la durée (\eng{a two-week course}), le contenu (\eng{elections, human rights, civic duties}).
\item \textbf{Ouvre l'échange poliment} : \eng{Excuse me, may I ask you a few questions?} ou \eng{Hello! I heard you attended a training course on democracy.}
\item \textbf{Pose des questions ouvertes} avec les mots interrogatifs : \eng{Where did the training take place?}, \eng{What did you learn?}, \eng{Who organised it?}
\item \textbf{Donne des informations précises} en répondant avec des phrases complètes, pas seulement \eng{yes} ou \eng{no} : \eng{The course was organised by a local NGO and lasted three days.}
\item \textbf{Conclus l'échange} : \eng{Thank you very much for the information.}, \eng{It was nice talking to you.}
\end{enumerate}
\textbf{Réflexes à retenir :} inversion du sujet et de l'auxiliaire dans les questions (\eng{did you attend}, jamais \eng{you did attend?}) ; prétérit pour raconter une formation terminée ; présent pour parler d'une formation en cours ou à venir.
\end{methode}

\begin{exoresolu}[Échanger des informations — type examen]
\textbf{Énoncé.} Your friend attended a training course on democracy organised in Buea. Write five questions you would ask him or her to get information about: (a) the organiser, (b) the duration, (c) the topics, (d) the number of participants, (e) the usefulness of the course.
\tcblower
\textbf{Solution détaillée.}
\begin{enumerate}
\item \textbf{Raisonnement grammatical.} La formation est terminée : on utilise le \textbf{prétérit} avec l'auxiliaire \eng{did} suivi de la base verbale. Pour une question fermée (réponse oui/non), on commence par l'auxiliaire ; pour une question ouverte, on place le mot interrogatif (\eng{who, how long, what, how many}) devant l'auxiliaire.
\item \textbf{Attention à l'ordre des mots} : mot interrogatif + auxiliaire + sujet + base verbale.
\item \textbf{Réponses rédigées :}
\begin{itemize}
\item (a) \eng{Who organised the training course?} — « Qui a organisé la formation ? » (\eng{who} est ici sujet : pas d'auxiliaire \eng{did}.)
\item (b) \eng{How long did the course last?} — « Combien de temps a duré la formation ? »
\item (c) \eng{What topics did you study during the training?} — « Quels sujets avez-vous étudiés ? »
\item (d) \eng{How many participants attended the course?} — « Combien de participants ont suivi la formation ? » (\eng{how many} + nom pluriel ; sujet = pas d'auxiliaire \eng{did}.)
\item (e) \eng{Did you find the course useful?} — « As-tu trouvé la formation utile ? » (question fermée : auxiliaire en tête.)
\end{itemize}
\item \textbf{Conclusion.} Cinq questions correctement construites, avec prétérit et inversion, couvrant toutes les informations demandées.
\end{enumerate}
\end{exoresolu}

\courssub{Relancer, réagir, ouvrir et conclure}

\begin{methode}[Les formules fonctionnelles de l'échange oral]
Un bon dialogue ne se limite pas à des questions-réponses : il faut \textbf{ouvrir}, \textbf{relancer}, \textbf{réagir} et \textbf{conclure}. Mémorise ces formules par fonction :
\begin{enumerate}
\item \textbf{Ouvrir la conversation} : \eng{Excuse me, can I ask you something?} ; \eng{I heard you took part in a training course. Is that true?}
\item \textbf{Demander une information} : \eng{Could you tell me more about the programme?} ; \eng{What exactly did you learn?}
\item \textbf{Relancer pour obtenir des précisions} : \eng{Really? Tell me more!} ; \eng{What do you mean by that?} ; \eng{Could you give me an example?} ; \eng{And what happened next?}
\item \textbf{Réagir à une information} : \eng{That's wonderful!} ; \eng{How interesting!} ; \eng{Oh, I see.} ; \eng{What a pity!} ; \eng{Congratulations!}
\item \textbf{Donner son avis} : \eng{In my opinion, such courses are very useful.} ; \eng{I think every citizen should learn about democracy.}
\item \textbf{Conclure} : \eng{Thank you for sharing this with me.} ; \eng{I must go now. See you soon!} ; \eng{It was a pleasure talking to you.}
\end{enumerate}
\textbf{Réflexe clé :} enchaîne les tours de parole naturellement — après chaque réponse, réagis d'abord (\eng{Oh, I see!}), puis relance (\eng{And what did you do after that?}).
\end{methode}

\begin{exoresolu}[Relancer et réagir — type examen]
\textbf{Énoncé.} Complete this dialogue with suitable expressions to react, ask for more details and conclude.

A: \eng{Hi, Peter! I heard you attended a workshop on human rights.}

B: \eng{Yes, I did. It was organised by a youth association in Yaoundé.}

A: (1) ............................................................ (react)

B: \eng{Yes! We discussed the right to education and freedom of expression.}

A: (2) ............................................................ (ask for more details)

B: \eng{For example, we learnt how to write a petition to the local authorities.}

A: (3) ............................................................ (give your opinion)

B: \eng{You are right. You should join the next session.}

A: (4) ............................................................ (conclude)
\tcblower
\textbf{Solution détaillée.}
\begin{enumerate}
\item \textbf{(1) Réagir} à une bonne nouvelle : \eng{That's wonderful!} ou \eng{How interesting!} « C'est formidable ! » — exclamation courte avec \eng{how} + adjectif ou \eng{what} + nom.
\item \textbf{(2) Relancer} pour obtenir un exemple : \eng{Could you give me an example?} « Pourrais-tu me donner un exemple ? » — formule polie avec \eng{could you} + base verbale.
\item \textbf{(3) Donner son avis} : \eng{In my opinion, such workshops are very useful for young citizens.} « À mon avis, de tels ateliers sont très utiles pour les jeunes citoyens. » — \eng{in my opinion} + phrase affirmative ; \eng{such} + nom pluriel sans article.
\item \textbf{(4) Conclure} : \eng{Thank you very much for the information. See you soon!} « Merci beaucoup pour ces informations. À bientôt ! »
\item \textbf{Conclusion.} Le dialogue complet est naturel et fluide : chaque tour de parole remplit une fonction précise (réaction, relance, opinion, clôture), ce qui est exactement ce que l'examinateur évalue à l'oral.
\end{enumerate}
\end{exoresolu}

\courssub{Dialogues modèles commentés}

\begin{methode}[Questionner et répondre : les structures à maîtriser]
\begin{enumerate}
\item \textbf{Distingue les deux types de questions} :
\begin{itemize}
\item \textbf{Questions fermées} (yes/no) : auxiliaire + sujet + verbe. \eng{Did you enjoy the workshop?} — \eng{Yes, I did. / No, I didn't.}
\item \textbf{Questions ouvertes} (informations) : wh-word + auxiliaire + sujet + verbe. \eng{When did the seminar start?}
\end{itemize}
\item \textbf{Choisis le bon mot interrogatif} : \eng{who} (personne), \eng{what} (chose), \eng{where} (lieu), \eng{when} (moment), \eng{why} (raison), \eng{how} (manière), \eng{how long} (durée), \eng{how many/much} (quantité).
\item \textbf{Réponds de façon adaptée et complète} : à une question fermée, réponds d'abord \eng{yes} ou \eng{no} avec l'auxiliaire, puis ajoute une information : \eng{Yes, I did. I learnt a lot about voting rights.}
\item \textbf{Soigne la politesse} : ajoute \eng{please}, \eng{could you tell me...?}, \eng{I'd like to know...} pour adoucir la demande : \eng{Could you tell me where the course took place?} (dans une question indirecte, pas d'inversion : \eng{where the course took place}, jamais \eng{where did the course take place} après \eng{could you tell me}).
\item \textbf{Vérifie l'accord des temps} : question au prétérit → réponse au prétérit ; question au futur (\eng{will you attend...?}) → réponse avec \eng{will} ou \eng{going to}.
\end{enumerate}
\end{methode}

\begin{exoresolu}[Questions et réponses adaptées — type examen]
\textbf{Énoncé.} Complete the following dialogue with appropriate questions or answers.

A: \eng{Hello, Mary. I heard you attended a training course on democracy last month.}

B: \eng{Yes, I did. It was very interesting.}

A: (1) ............................................................?

B: \eng{It took place at the City Hall in Douala.}

A: (2) ............................................................?

B: \eng{About fifty young people participated.}

A: (3) ............................................................?

B: \eng{Because I wanted to learn about my rights and duties as a citizen.}

A: (4) ............................................................?

B: \eng{Yes, I will. Next month, there is another course on elections.}
\tcblower
\textbf{Solution détaillée.}
\begin{enumerate}
\item \textbf{Méthode} : chaque question doit correspondre exactement à la réponse qui la suit. On repère dans la réponse le type d'information (lieu, quantité, raison, intention future), ce qui détermine le mot interrogatif et le temps.
\item \textbf{(1)} La réponse indique un \textbf{lieu} (\eng{at the City Hall in Douala}) au passé → \eng{Where did the training course take place?} « Où la formation a-t-elle eu lieu ? » (prétérit, inversion avec \eng{did}.)
\item \textbf{(2)} La réponse donne un \textbf{nombre de personnes} → \eng{How many young people participated?} « Combien de jeunes ont participé ? » (\eng{how many} + pluriel ; ici le groupe interrogatif est sujet, donc pas d'auxiliaire \eng{did}.)
\item \textbf{(3)} La réponse commence par \eng{because} → question de \textbf{raison} : \eng{Why did you attend the course?} « Pourquoi as-tu suivi cette formation ? »
\item \textbf{(4)} La réponse \eng{Yes, I will} annonce une \textbf{intention future} → question fermée avec \eng{will} : \eng{Will you attend another course?} « Assisteras-tu à une autre formation ? »
\item \textbf{Conclusion.} Le dialogue est cohérent : chaque question respecte l'ordre des mots anglais (mot interrogatif + auxiliaire + sujet + base verbale) et l'accord des temps avec la réponse.
\end{enumerate}
\end{exoresolu}

\courssub{Prononciation, accentuation et intonation}

\begin{propriete}[Accentuation des mots clés du thème (Word Stress)]
En anglais, l'accent de mot (word stress) est fixe et change le sens ou la compréhension. Marque l'accent sur la syllabe en \textbf{MAJUSCULES} (notation française entre guillemets).
\begin{center}
\begin{tabular}{|l|l|l|}
\hline
\rowcolor{popblueL} \textbf{Mot} & \textbf{Accentuation (notation FR)} & \textbf{Règle / Note} \\
\hline
\eng{democracy} & « dé-\textbf{MO}-kra-si » & Accent sur l'antépénultième (règle des mots en \eng{-cy}). \\
\eng{participant} & « par-\textbf{TI}-si-pant » & Accent sur l'antépénultième (règle des mots en \eng{-ant}). \\
\eng{training} & « \textbf{TRÉ}-ning » & Accent sur la 1ère syllabe (nom/verbe bisyllabique). \\
\eng{citizen} & « \textbf{SI}-ti-zen » & Accent sur la 1ère syllabe. \\
\eng{organise} & « \textbf{OR}-ga-naïz » & Verbe : accent sur la 1ère syllabe. \\
\eng{organisation} & « or-ga-ni-\textbf{ZA}-tion » & Nom en \eng{-tion} : accent sur la syllabe avant \eng{-tion}. \\
\eng{facilitator} & « fa-si-li-\textbf{TA}-teur » & Accent sur l'antépénultième. \\
\eng{certificate} & « \textbf{SER}-ti-fi-ket » & Accent sur la 1ère syllabe (attention : pas « cer-ti-\textbf{FI}-cate »). \\
\hline
\end{tabular}
\end{center}
\end{propriete}

\begin{propriete}[Intonation des phrases (Sentence Intonation)]
\begin{itemize}
\item \textbf{Questions fermées (Yes/No) :} intonation \textbf{montante} ↗ à la fin. \eng{Did you enjoy the course? ↗}
\item \textbf{Questions ouvertes (Wh- questions) :} intonation \textbf{descendante} ↘ à la fin. \eng{Where did it take place? ↘}
\item \textbf{Affirmations / Réponses :} intonation \textbf{descendante} ↘. \eng{It took place in Bamenda. ↘}
\item \textbf{Listes / Énumérations :} montante ↗ sur les éléments non finaux, descendante ↘ sur le dernier. \eng{We studied rights ↗, duties ↗, and elections ↘.}
\item \textbf{Questions indirectes polies :} intonation souvent descendante ↘ (comme une affirmation). \eng{Could you tell me when it starts? ↘}
\end{itemize}
\end{propriete}

\begin{experience}[Exercice de prononciation guidée]
En binôme, lisez le dialogue du \textbf{Script d'écoute} (Mise en situation) en respectant :
\begin{enumerate}
\item L'accent de mot sur \eng{de-MO-cra-cy}, \eng{par-TI-ci-pant}, \eng{OR-ga-nised}, \eng{com-mu-NI-ty}, \eng{HU-man}, \eng{e-LEC-tion}.
\item L'intonation montante ↗ pour \eng{Who organised it?}, \eng{Did you find it useful?}, \eng{Will there be another session?}
\item L'intonation descendante ↘ pour \eng{It was organised by an NGO.}, \eng{We discussed human rights.}, \eng{The trainers were excellent.}
\item La liaison (linking) : \eng{attended a} (« a-tten-de-da »), \eng{training on} (« trai-ni-non »), \eng{about sixty} (« a-bou-tix-ty »).
\end{enumerate}
\end{experience}

\courssub{Jeux de rôle guidés et production orale}

\begin{methode}[Construire et jouer un dialogue complet]
\begin{enumerate}
\item \textbf{Lis attentivement la consigne} : identifie les rôles (\eng{Student A / Student B}), le lieu, le sujet et le nombre de répliques attendu (généralement 6 à 10 échanges).
\item \textbf{Planifie la structure du dialogue} : ouverture → 2 ou 3 questions-réponses → réactions et relances → conclusion. Respecte l'alternance des tours de parole.
\item \textbf{Rédige des phrases simples et correctes} : mieux vaut une phrase courte juste qu'une phrase longue fautive. Utilise le prétérit pour le passé, \eng{will / going to} pour les projets.
\item \textbf{Varie les formules} : ne répète pas trois fois la même question ; alterne \eng{What about...?}, \eng{Could you tell me...?}, \eng{Did you...?}
\item \textbf{À l'oral, soigne la prononciation et l'intonation} : voix montante ↗ pour les questions fermées (\eng{Did you enjoy it? ↗}), voix descendante ↘ pour les questions en \eng{wh-} (\eng{Where did it take place? ↘}) et pour les affirmations.
\item \textbf{Accentue correctement les mots clés} : \eng{de-MO-cra-cy} (« dé-MO-kra-si »), \eng{PAR-ti-ci-pant} (« par-TI-si-pant »), \eng{TRAI-ning} (« TRÉ-ning »), \eng{CIT-i-zen} (« SI-ti-zen »).
\end{enumerate}
\end{methode}

\begin{exoresolu}[Jeu de rôle — production d'un dialogue, type examen]
\textbf{Énoncé.} \emph{Role play.} Student A is a journalist for the school magazine. Student B attended a two-day training course on democracy organised in Bamenda. Write a dialogue of at least eight exchanges in which the journalist interviews the participant.
\tcblower
\textbf{Solution détaillée.}
\begin{enumerate}
\item \textbf{Analyse de la tâche} : deux rôles (journaliste / participant), sujet imposé (formation de deux jours à Bamenda), au moins huit répliques. Structure : salutation → questions (organisateur, contenu, appréciation, projets) → remerciements.
\item \textbf{Dialogue rédigé :}

\eng{A: Good morning! I am a journalist for the school magazine. May I ask you a few questions about the training course you attended?}

\eng{B: Good morning! Yes, of course.}

\eng{A: Where and when did the course take place?}

\eng{B: It took place last week at the community hall in Bamenda.}

\eng{A: Who organised it?}

\eng{B: A local NGO organised it, with the help of the town council.}

\eng{A: What did you learn during those two days?}

\eng{B: We learnt about elections, the rights and duties of citizens, and how to vote responsibly.}

\eng{A: That sounds very interesting! Did you enjoy the course?}

\eng{B: Yes, I did. The trainers were excellent and the discussions were lively.}

\eng{A: Will you attend another course in the future?}

\eng{B: Certainly! Next month, there is a workshop on human rights.}

\eng{A: Thank you very much for your time.}

\eng{B: You are welcome.}
\item \textbf{Justifications grammaticales} : prétérit avec \eng{did} + base verbale pour les questions sur le passé (\eng{did the course take place}, \eng{did you learn}) ; \eng{who organised it?} sans auxiliaire car \eng{who} est sujet ; futur avec \eng{will} pour le projet ; réaction \eng{That sounds very interesting!} pour relancer ; conclusion polie \eng{Thank you for your time.}
\item \textbf{Conclusion.} Le dialogue compte huit échanges complets, respecte l'alternance des rôles, utilise les formules fonctionnelles étudiées et reste grammaticalement correct : c'est la production attendue au Bac.
\end{enumerate}
\end{exoresolu}

\begin{exercice}[Entraînement à l'oral — Jeux de rôle]
\begin{enumerate}
\item \textbf{Role Play 1.} Student A: You are a member of a youth association in Douala. Student B: You attended a seminar on \eng{civic engagement} in Yaoundé. Student A interviews Student B for the association's newsletter (8 exchanges minimum).
\item \textbf{Role Play 2.} Student A: You want to organise a training course on \eng{human rights} in your school. Student B: You are a trainer from an NGO. Discuss the venue, duration, topics and target participants (8 exchanges minimum).
\item \textbf{Role Play 3.} Two friends. One attended a workshop on \eng{the electoral process} in Buea. The other wants details to decide whether to attend the next session (8 exchanges minimum).
\end{enumerate}
\end{exercice}

\begin{attention}[Les 5 erreurs qui coûtent des points au Bac]
\begin{enumerate}
\item \textbf{Oublier l'inversion dans les questions} : ne dis jamais \eng{Where the course took place?} ; dis \eng{Where did the course take place?} Attention : après \eng{Could you tell me...}, il n'y a \emph{pas} d'inversion : \eng{Could you tell me where the course took place?}
\item \textbf{Mal conjuguer le prétérit avec l'auxiliaire} : après \eng{did}, le verbe reste à la base verbale : \eng{Did you attend?} et non \eng{Did you attended?} De même : \eng{I didn't go}, jamais \eng{I didn't went}.
\item \textbf{Traduire mot à mot le français} : « assister à une formation » se dit \eng{attend a training course} (sans \eng{to} après \eng{attend}) ; « participer à » se dit \eng{take part in} ou \eng{participate in} (la préposition \eng{in} est obligatoire).
\item \textbf{Confondre les faux amis} : \eng{formation} (français) ≠ \eng{formation} (anglais = « formation géologique ») ; dis \eng{training course}. \eng{Actually} signifie « en fait », pas « actuellement » ; pour « actuellement », dis \eng{currently} ou \eng{at present}.
\item \textbf{Négliger la prononciation et l'intonation à l'oral} : accentue \eng{de-MO-cra-cy} sur la deuxième syllabe (et non sur la première comme en français) ; monte la voix ↗ dans les questions fermées (\eng{Did you enjoy it? ↗}) et descends-la ↘ dans les questions en \eng{wh-}. Une bonne intonation rapporte des points à l'épreuve orale.
\end{enumerate}
\end{attention}

\newpage

\coursec{Exercices}

\courssub{Je vérifie mes connaissances}

\begin{exercice}[QCM : les formules fonctionnelles]
Pour chaque situation, choisis la formule la plus adaptée (a, b ou c).
\begin{enumerate}
\item You want to know the duration of the training course.
\begin{itemize}
\item[a)] How long did it last?
\item[b)] How far is it?
\item[c)] How many times did you go?
\end{itemize}
\item Your friend says the course was free. You react with surprise.
\begin{itemize}
\item[a)] What a pity!
\item[b)] Really? That's interesting!
\item[c)] Never mind.
\end{itemize}
\item You want to end the conversation politely.
\begin{itemize}
\item[a)] Shut up now.
\item[b)] Thanks for telling me all about it.
\item[c)] Go away, please.
\end{itemize}
\item You want to ask for information in a formal way.
\begin{itemize}
\item[a)] Tell me everything, now!
\item[b)] Could you tell me how I can register?
\item[c)] Registration. How?
\end{itemize}
\end{enumerate}
\end{exercice}

\begin{exercice}[Vrai ou faux ?]
Réponds par \textbf{Vrai} ou \textbf{Faux} et justifie ta réponse en une phrase en anglais.
\begin{enumerate}
\item In English, we say \eng{“I have participated to a training course.”}
\item \eng{“How long did it last?”} is a question about duration.
\item To keep a conversation going, you can say \eng{“And then what happened?”}
\item \eng{“Could you tell me where the centre is?”} is more polite than \eng{“Where is the centre?”}
\end{enumerate}
\end{exercice}

\begin{exercice}[Vocabulaire : les mots de la formation civique]
Complète chaque phrase avec le mot anglais qui convient, choisi dans la liste : \emph{training course, workshop, certificate, participant, to register, democracy}.
\begin{enumerate}
\item Every \dots\ received a \dots\ at the end of the course.
\item You must \dots\ online before the 15th of May.
\item The \dots\ on human rights lasted three days.
\item During the \dots\ , we worked in small groups on voting rights.
\item \dots\ means that citizens can choose their leaders freely.
\end{enumerate}
\end{exercice}

\begin{exercice}[Grammaire : present perfect ou past simple ?]
Mets le verbe entre parenthèses au temps qui convient (\emph{present perfect} ou \emph{past simple}).
\begin{enumerate}
\item I \dots\ (attend) a training course on democracy last month.
\item She \dots\ (never / take part) in a workshop before.
\item They \dots\ (learn) a lot about elections during the session.
\item \dots\ you \dots\ (ever / receive) a certificate?
\item We \dots\ (meet) the trainer in Yaoundé two weeks ago.
\end{enumerate}
\end{exercice}

\courssub{Je m'entraîne}

\begin{exercice}[Dialogue à trous]
Complète ce dialogue avec les formules de la liste : \emph{Really? — How long did it last? — Thanks for telling me — Could you tell me how to register? — What did you learn?}
\begin{textebox}[A dialogue to complete]
Amina: Hi, Paul! I heard you attended a training course on democracy. \dots\ ?
Paul: Yes, I did! It was organised by the town council.
Amina: \dots\ ?
Paul: It lasted two days, Saturday and Sunday.
Amina: \dots\ ?
Paul: We learned about citizens' rights and how elections work.
Amina: \dots\ ? That sounds really useful!
Paul: It was! Do you want to join the next session?
Amina: Yes! \dots\ ?
Paul: Just go to the community centre with your ID card.
Amina: Great. \dots\ , Paul!
\end{textebox}
\end{exercice}

\begin{exercice}[Questions directes ou indirectes ?]
Transforme ces cinq questions directes en questions polies et formelles, en commençant par \emph{Could you tell me…} ou \emph{I'd like to know…}.
\begin{enumerate}
\item How can I register?
\item Where does the course take place?
\item How much does it cost?
\item Who is the trainer?
\item When does the next session start?
\end{enumerate}
\end{exercice}

\begin{exercice}[Appariement : questions et réponses]
Associe chaque question (1--5) à la réponse qui convient (a--e).
\begin{enumerate}
\item How long did the training last?
\item What did you learn there?
\item Was the course expensive?
\item Who organised the workshop?
\item Would you recommend it?
\end{enumerate}
\begin{itemize}
\item[a)] A local NGO called “Youth for Democracy”.
\item[b)] Absolutely! You should attend the next one.
\item[c)] Three full days, from Monday to Wednesday.
\item[d)] No, it was completely free.
\item[e)] How to debate and respect other people's opinions.
\end{itemize}
\end{exercice}

\begin{exercice}[Prononciation : intonation montante ou descendante ?]
Classe les dix questions suivantes selon leur intonation : montante ($\nearrow$) ou descendante ($\searrow$). Puis lis-les à voix haute avec la bonne mélodie.
\begin{enumerate}
\item Did you enjoy the course?
\item How long did it last?
\item Was it free?
\item Where was the training centre?
\item Have you ever attended a workshop?
\item What did you learn?
\item Can I register online?
\item Who was the trainer?
\item Did you get a certificate?
\item When does the next session begin?
\end{enumerate}
\end{exercice}

\courssub{Je me prépare au Bac}

\begin{exercice}[Compréhension écrite et étude de la langue]
Lis le texte suivant, puis réponds aux questions en anglais, par des phrases complètes.
\begin{textebox}[A training course that changed my mind]
Last holidays, I attended a five-day training course on democracy at the community centre of Buea. At first, I was not very enthusiastic: I thought it would be long speeches and boring documents. I was completely wrong.

On the first day, the trainer, a young lawyer from Douala, asked us a simple question: “What can you do for your country?” Nobody answered. Then she explained that democracy is not only about voting every few years; it is also about listening to others, respecting different opinions and taking part in the life of our community.

During the workshops, we organised a mock election. I was a candidate, and I had to present a programme in front of fifty participants. My hands were shaking, but I did it! We also learned how to check information before sharing it, because false news can divide people.

At the end of the course, each participant received a certificate. More importantly, I left with a new idea: young people are not too young to act. Next month, our group will organise a debate on human rights in our school. I have already registered, and this time, I will not be afraid to speak.
\end{textebox}
\textbf{Comprehension questions}
\begin{enumerate}
\item Where and when did the narrator attend the training course? (2 marks)
\item What did the narrator think about the course before attending it? (2 marks)
\item According to the trainer, what is democracy, besides voting? (2 marks)
\item What activity did the participants do during the workshops? (2 marks)
\item What did the narrator decide to do after the course? (2 marks)
\end{enumerate}
\textbf{Language work}
\begin{enumerate}
\item Find in the text two verbs in the past simple and two verbs in the present perfect. (2 marks)
\item Turn into an indirect question: \eng{“What can you do for your country?”} (Begin: \emph{She asked us…}) (2 marks)
\item Find in the text the opposite of: \emph{interesting} (paragraph 1) and \emph{unite} (paragraph 3). (2 marks)
\end{enumerate}
\end{exercice}

\begin{exercice}[Traduction]
Traduis en anglais le passage suivant (environ 60 mots).
\begin{textebox}[Texte à traduire]
Le mois dernier, mon amie Clarisse a participé à une formation sur la démocratie à Yaoundé. Elle m'a dit que les ateliers étaient très intéressants. Elle a appris à débattre et à respecter les opinions des autres. À la fin de la formation, elle a reçu un certificat. Maintenant, elle veut organiser un club de débat dans notre lycée. Je me suis déjà inscrit !
\end{textebox}
\emph{Conseil : fais attention au choix des temps (past simple / present perfect) et aux prépositions.}
\end{exercice}

\begin{exercice}[Expression orale : jeu de rôle type épreuve du Bac]
\textbf{Situation :} \emph{Your friend attended a training course on democracy. Interview him/her about it.}
\begin{enumerate}
\item \textbf{Préparation (10 min) :} en binôme, préparez un dialogue de 10 à 12 répliques. L'élève A pose au moins cinq questions (durée, lieu, contenu, coût, avis personnel) ; l'élève B répond avec des informations précises et utilise le past simple et le present perfect.
\item \textbf{Production (5 min) :} jouez la scène devant la classe. Utilisez au moins deux formules de réaction (\emph{Really? That's interesting!}) et une formule de conclusion (\emph{Thanks for telling me…}).
\item \textbf{Compte rendu écrit :} rédigez ensuite un court report (5--6 phrases, 60--80 mots) sur la participation de votre camarade à la formation.
\end{enumerate}
\textbf{Grille d'évaluation (sur 20) :} vocabulaire du thème (5) ; formules fonctionnelles (5) ; interaction et relances (5) ; prononciation et intonation (5).
\end{exercice}

\newpage

\coursec{Corrigés des exercices}

\begin{corrige}[Exercice 1 — QCM : les formules fonctionnelles]
\begin{enumerate}
\item \textbf{b) est faux ; la bonne réponse est a)} \eng{How long did it last?} — « Combien de temps cela a-t-il duré ? » C'est la seule question qui porte sur la \emph{durée}. \eng{How far is it?} interroge sur la distance ; \eng{How many times did you go?} sur la fréquence.
\item \textbf{b)} \eng{Really? That's interesting!} — « Vraiment ? C'est intéressant ! » Formule de réaction (surprise positive). \eng{What a pity!} exprime le regret ; \eng{Never mind.} signifie « tant pis ».
\item \textbf{b)} \eng{Thanks for telling me all about it.} — « Merci de m'avoir tout raconté. » C'est la formule de conclusion polie. Les options a) et c) sont grossières et à bannir absolument.
\item \textbf{b)} \eng{Could you tell me how I can register?} — « Pourriez-vous me dire comment je peux m'inscrire ? » Question indirecte, registre formel. L'option a) est un ordre brutal ; l'option c) est télégraphique et impolie.
\end{enumerate}
\textbf{Point de méthode :} à l'oral comme à l'écrit, le choix de la formule dépend du \emph{registre} (familier / neutre / formel) et de la \emph{fonction} visée (demander, réagir, conclure).
\end{corrige}

\begin{corrige}[Exercice 2 — Vrai ou faux ?]
\begin{enumerate}
\item \textbf{Faux.} \eng{In English, we say “I have participated in a training course”, because the verb “to participate” is always followed by the preposition “in”, never “to”.} (Faux ami avec le français « participer à ».)
\item \textbf{Vrai.} \eng{“How long did it last?” asks about duration, because “how long” is used to ask about the length of time of an event.}
\item \textbf{Vrai.} \eng{“And then what happened?” is a follow-up question that encourages your partner to continue speaking, so it keeps the conversation going.}
\item \textbf{Vrai.} \eng{“Could you tell me where the centre is?” is an indirect question, and indirect questions with “Could you tell me…” are more polite and formal than direct questions.}
\end{enumerate}
\textbf{Points de vigilance :} attention à l'ordre des mots dans la question indirecte : \eng{Could you tell me where the centre \textbf{is}} (sujet + verbe), et non \emph{where is the centre}.
\end{corrige}

\begin{corrige}[Exercice 3 — Vocabulaire : les mots de la formation civique]
\begin{enumerate}
\item \eng{Every \textbf{participant} received a \textbf{certificate} at the end of the course.} « Chaque participant a reçu un certificat à la fin de la formation. »
\item \eng{You must \textbf{register} online before the 15th of May.} « Tu dois t'inscrire en ligne avant le 15 mai. » (Après \eng{must}, base verbale : \eng{register}, sans \emph{to}.)
\item \eng{The \textbf{training course} on human rights lasted three days.} « La formation sur les droits de l'homme a duré trois jours. »
\item \eng{During the \textbf{workshop}, we worked in small groups on voting rights.} « Pendant l'atelier, nous avons travaillé en petits groupes sur le droit de vote. »
\item \eng{\textbf{Democracy} means that citizens can choose their leaders freely.} « La démocratie signifie que les citoyens peuvent choisir librement leurs dirigeants. »
\end{enumerate}
\textbf{Astuce :} repère les indices grammaticaux : \eng{every} + nom singulier ; \eng{must} + verbe ; \eng{the \dots\ on human rights} appelle un nom d'événement (\eng{training course}) ; \eng{during the \dots} un nom d'activité (\eng{workshop}).
\end{corrige}

\begin{corrige}[Exercice 4 — Grammaire : present perfect ou past simple ?]
\begin{enumerate}
\item \eng{I \textbf{attended} a training course on democracy last month.} — \emph{Past simple} : repère temporel précis et terminé (\eng{last month}).
\item \eng{She \textbf{has never taken part} in a workshop before.} — \emph{Present perfect} : expérience de vie sans date précise (\eng{never}, \eng{before}). Participe passé irrégulier : \emph{take -- took -- taken}.
\item \eng{They \textbf{learnt} (ou \textbf{learned}) a lot about elections during the session.} — \emph{Past simple} : l'action est située dans un moment passé délimité (\eng{during the session}).
\item \eng{\textbf{Have} you \textbf{ever received} a certificate?} — \emph{Present perfect} : question sur une expérience (\eng{ever}).
\item \eng{We \textbf{met} the trainer in Yaoundé two weeks ago.} — \emph{Past simple} : repère \eng{ago}. Participe/passé irrégulier : \emph{meet -- met -- met}.
\end{enumerate}
\textbf{Règle à retenir :} date ou repère passé précis (\eng{last month}, \eng{ago}, \eng{yesterday}, \eng{in 2023}) $\Rightarrow$ past simple. Expérience sans date (\eng{ever}, \eng{never}, \eng{already}, \eng{before}) ou lien avec le présent $\Rightarrow$ present perfect.
\end{corrige}

\begin{corrige}[Exercice 5 — Dialogue à trous]
\begin{textebox}[Corrigé du dialogue]
Amina: Hi, Paul! I heard you attended a training course on democracy. \textbf{Really?}
Paul: Yes, I did! It was organised by the town council.
Amina: \textbf{How long did it last?}
Paul: It lasted two days, Saturday and Sunday.
Amina: \textbf{What did you learn?}
Paul: We learned about citizens' rights and how elections work.
Amina: \textbf{Really?} That sounds really useful! \emph{(ou : réaction libre)}
Paul: It was! Do you want to join the next session?
Amina: Yes! \textbf{Could you tell me how to register?}
Paul: Just go to the community centre with your ID card.
Amina: Great. \textbf{Thanks for telling me}, Paul!
\end{textebox}
\textbf{Justification :} la réponse \eng{It lasted two days} impose la question sur la durée ; \eng{We learned about\ldots} impose \eng{What did you learn?} ; la réponse \eng{Just go to the community centre\ldots} répond à une demande d'information sur l'inscription ; la dernière réplique clôt poliment l'échange. \emph{Remarque :} \eng{Really?} peut figurer à la première ou à la quatrième ligne ; les deux positions sont acceptables si la logique du dialogue est respectée.
\end{corrige}

\begin{corrige}[Exercice 6 — Questions directes ou indirectes ?]
\begin{enumerate}
\item \eng{Could you tell me how I can register?}
\item \eng{Could you tell me where the course takes place?}
\item \eng{I'd like to know how much it costs.}
\item \eng{Could you tell me who the trainer is?}
\item \eng{I'd like to know when the next session starts.}
\end{enumerate}
\textbf{Justification grammaticale :} dans une question indirecte introduite par \eng{Could you tell me\ldots} ou \eng{I'd like to know\ldots}, on retrouve l'ordre \emph{sujet + verbe} (comme dans une phrase affirmative) et il n'y a \textbf{pas d'inversion} ni d'auxiliaire \eng{do/does/did} : \eng{where the course \textbf{takes place}}, \eng{who the trainer \textbf{is}}. Le point d'interrogation se maintient après \eng{Could you tell me\ldots?} ; après \eng{I'd like to know\ldots} on met un point.
\end{corrige}

\begin{corrige}[Exercice 7 — Appariement : questions et réponses]
\begin{center}
\begin{tabular}{|l|l|}\hline
\rowcolor{popblueL} Question & Réponse \\ \hline
1. How long did the training last? & \textbf{c)} Three full days, from Monday to Wednesday. \\ \hline
2. What did you learn there? & \textbf{e)} How to debate and respect other people's opinions. \\ \hline
3. Was the course expensive? & \textbf{d)} No, it was completely free. \\ \hline
4. Who organised the workshop? & \textbf{a)} A local NGO called “Youth for Democracy”. \\ \hline
5. Would you recommend it? & \textbf{b)} Absolutely! You should attend the next one. \\ \hline
\end{tabular}
\end{center}
\textbf{Justification :} \eng{How long\ldots?} appelle une durée (c) ; \eng{What did you learn\ldots?} un contenu appris (e) ; \eng{Was it expensive?} une réponse sur le prix (d) ; \eng{Who\ldots?} un organisme ou une personne (a) ; \eng{Would you recommend it?} une opinion et un conseil (b).
\end{corrige}

\begin{corrige}[Exercice 8 — Prononciation : intonation montante ou descendante ?]
\textbf{Règle :} les questions fermées (réponse \emph{yes/no}, commençant par un auxiliaire : \eng{Did\ldots?}, \eng{Was\ldots?}, \eng{Have\ldots?}, \eng{Can\ldots?}) ont une intonation \textbf{montante} ($\nearrow$). Les questions ouvertes en \emph{wh-} (\eng{How}, \eng{Where}, \eng{What}, \eng{Who}, \eng{When}) ont une intonation \textbf{descendante} ($\searrow$).
\begin{center}
\begin{tabular}{|l|l|}\hline
\rowcolor{popblueL} Montante ($\nearrow$) & Descendante ($\searrow$) \\ \hline
1. Did you enjoy the course? & 2. How long did it last? \\ \hline
3. Was it free? & 4. Where was the training centre? \\ \hline
5. Have you ever attended a workshop? & 6. What did you learn? \\ \hline
7. Can I register online? & 8. Who was the trainer? \\ \hline
9. Did you get a certificate? & 10. When does the next session begin? \\ \hline
\end{tabular}
\end{center}
\textbf{Entraînement oral :} pour la question 1, monter la voix sur \eng{course} (« cour-ssse » en montant) ; pour la question 2, descendre sur \eng{last}. Faire répéter chaque question en chœur, puis individuellement.
\end{corrige}

\begin{corrige}[Exercice 9 — Compréhension écrite et étude de la langue]
\textbf{Comprehension questions (réponses modèles)}
\begin{enumerate}
\item \eng{The narrator attended the training course at the community centre of Buea, during the last holidays.} (2 marks : lieu 1 + temps 1)
\item \eng{Before attending it, the narrator thought the course would be boring, with long speeches and boring documents, so he/she was not very enthusiastic.} (2 marks)
\item \eng{According to the trainer, democracy is not only about voting; it is also about listening to others, respecting different opinions and taking part in the life of the community.} (2 marks)
\item \eng{During the workshops, the participants organised a mock election, and the narrator had to present a programme in front of fifty participants.} (2 marks)
\item \eng{After the course, the narrator decided to organise a debate on human rights at school with his/her group, and he/she has already registered.} (2 marks)
\end{enumerate}
\textbf{Language work}
\begin{enumerate}
\item Past simple : \eng{attended}, \eng{was}, \eng{explained}, \eng{received} (au choix). Present perfect : \eng{have already registered} ; on peut aussi citer \eng{has} s'il apparaît — ici la forme certaine du texte est \eng{I have already registered}. (2 marks : 0,5 par verbe correctement identifié)
\item \eng{She asked us what we could do for our country.} — Question indirecte au passé : \eng{can} $\rightarrow$ \eng{could}, \eng{you} $\rightarrow$ \eng{we}, \eng{your} $\rightarrow$ \eng{our}, ordre sujet + verbe, sans point d'interrogation. (2 marks)
\item Opposite of \emph{interesting} : \eng{boring} (paragraph 1). Opposite of \emph{unite} : \eng{divide} (paragraph 3). (2 marks)
\end{enumerate}
\textbf{Barème indicatif :} compréhension 10 pts ; étude de la langue 6 pts ; qualité de l'expression anglaise 4 pts (total 20).
\textbf{Points de vigilance :} répondre par des phrases complètes avec les mots de l'élève (éviter le recopiage intégral) ; dans la question indirecte, ne jamais écrire \emph{what can we do} après \eng{She asked us}.
\end{corrige}

\begin{corrige}[Exercice 10 — Traduction]
\textbf{Proposition de traduction :}
\begin{textebox}[Model translation]
Last month, my friend Clarisse took part in a training course on democracy in Yaoundé. She told me that the workshops were very interesting. She learnt how to debate and how to respect other people's opinions. At the end of the course, she received a certificate. Now, she wants to organise a debate club in our school. I have already registered!
\end{textebox}
\textbf{Points de vigilance et barème indicatif (sur 10) :}
\begin{itemize}
\item « a participé à » $\rightarrow$ \eng{took part in} ou \eng{participated in} (jamais \emph{participated to}) — 1,5 pt.
\item « une formation » $\rightarrow$ \eng{a training course} (faux ami : \emph{formation} seul est rarement correct) — 1 pt.
\item « Elle m'a dit que\ldots » $\rightarrow$ \eng{She told me that\ldots} avec concordance des temps : \eng{the workshops \textbf{were}} — 1,5 pt.
\item « elle a appris à débattre » $\rightarrow$ \eng{she learnt how to debate} — 1 pt.
\item « les opinions des autres » $\rightarrow$ \eng{other people's opinions} (génitif) — 1 pt.
\item « un club de débat » $\rightarrow$ \eng{a debate club} ; « notre lycée » $\rightarrow$ \eng{our school} — 1 pt.
\item « Je me suis déjà inscrit ! » $\rightarrow$ \eng{I have already registered!} (present perfect avec \eng{already}, résultat présent) — 1,5 pt.
\item Fluidité générale, orthographe, ponctuation — 1,5 pt.
\end{itemize}
\end{corrige}

\begin{corrige}[Exercice 11 — Expression orale : jeu de rôle type épreuve du Bac]
\textbf{Dialogue modèle (10--12 répliques) :}
\begin{textebox}[Model role-play]
A: Hi, Sandra! I heard you attended a training course on democracy. Is that true?
B: Yes, it is! I took part in it last month, in Bamenda.
A: Really? That's interesting! How long did it last?
B: It lasted three days, from Friday to Sunday.
A: And who organised it?
B: A local NGO, together with our school. It was completely free.
A: Free? Wonderful! What did you learn there?
B: We learned about citizens' rights, and we organised a mock election.
A: Have you ever spoken in public before?
B: No, I had never done it before, but now I feel more confident.
A: Would you recommend the course to other students?
B: Absolutely! I have already registered for the next session.
A: Great! Could you tell me how I can register too?
B: Of course. Just fill in the form at the students' office.
A: Thanks for telling me all about it, Sandra!
B: You're welcome!
\end{textebox}
\textbf{Report modèle (60--80 mots) :}
\begin{textebox}[Model report]
My classmate Sandra attended a three-day training course on democracy in Bamenda last month. It was organised by a local NGO and it was completely free. During the workshops, she learned about citizens' rights and took part in a mock election. She told me she had never spoken in public before, but now she feels more confident. She has already registered for the next session, and she advised me to register too.
\end{textebox}
\textbf{Barème indicatif (sur 20) :} vocabulaire du thème (\eng{training course}, \eng{workshop}, \eng{to register}, \eng{certificate}, \eng{citizens' rights}) — 5 pts ; formules fonctionnelles (demander, réagir, conclure) — 5 pts ; interaction (relances, écoute, réponses adaptées) — 5 pts ; prononciation, accentuation et intonation (montante/descendante) — 5 pts.
\textbf{Points de vigilance :} alterner past simple (faits datés) et present perfect (expériences) ; employer au moins une question indirecte polie ; éviter \emph{participate to} et l'intonation montante sur les questions en \emph{wh-} ; conclure toujours poliment (\eng{Thanks for telling me\ldots}).
\end{corrige}

\newpage

\coursec{Fiche bilan}

\begin{popfigure}
\begin{tikzpicture}[
  mindmap,
  grow cyclic,
  every node/.style={concept, execute at begin node=\hskip0pt},
  concept color=popblue,
  level 1/.append style={level distance=4.5cm, sibling angle=360/6},
  level 2/.append style={level distance=3cm, sibling angle=360/4},
  text=white,
  font=\small\bfseries
]
\node[concept, scale=1.2, align=center]{Lesson 40\\Speaking:\\Democracy Training}
  child[concept color=poporange] { node[concept, align=center]{Vocabulary\\Democracy \& Training} }
  child[concept color=popgreen] { node[concept, align=center]{Asking for\\Information} }
  child[concept color=poppurple] { node[concept, align=center]{Giving\\Information} }
  child[concept color=poppink] { node[concept, align=center]{Reacting\\ \& Follow-up} }
  child[concept color=popgold] { node[concept, align=center]{Opening\\ \& Closing} }
  child[concept color=popdark] { node[concept, align=center]{Pronunciation\\ \& Intonation} };
\end{tikzpicture}
\legende{Carte mentale : Lesson 40 — Exchanges on Democracy Training}
\end{popfigure}

\begin{bilan}
\end{bilan}

\courssub{Lexique clé (Vocabulary)}

\begin{center}
\begin{tabularx}{\linewidth}{|X|X|X|}
\hline
\rowcolor{popblueL} \textbf{English} & \textbf{Français} & \textbf{Contexte / Collocation} \\
\hline
a training course & une formation / un stage & attend a training course on democracy \\
a workshop & un atelier & participate in a workshop \\
democracy & la démocratie & democratic values, democratic process \\
civic education & l'éducation civique & civic education programme \\
human rights & les droits de l'homme & human rights advocacy \\
rule of law & l'État de droit & respect for the rule of law \\
good governance & la bonne gouvernance & principles of good governance \\
a facilitator / a trainer & un formateur / un animateur & The facilitator explained the module. \\
a participant / a trainee & un participant / un stagiaire & Participants received certificates. \\
a module / a session & un module / une session & The first module covers elections. \\
to enrol (in) / to register (for) & s'inscrire (à) & I want to enrol in this course. \\
to empower & autonomiser / donner les moyens & Empowering youth through education. \\
to raise awareness & sensibiliser & The NGO raises awareness on voting. \\
a certificate & un certificat / une attestation & receive a certificate of completion \\
funding / sponsorship & un financement / un parrainage & The training is funded by the EU. \\
\hline
\end{tabularx}
\end{center}

\courssub{Fonctions langagières \& Grammaire essentielle}

\begin{center}
\begin{tabularx}{\linewidth}{|p{3.5cm}|X|X|}
\hline
\rowcolor{popblueL} \textbf{Fonction} & \textbf{Formules types (English)} & \textbf{Traduction / Note} \\
\hline
\cle{Demander une information} & \eng{Can you tell me more about...?} & « Peux-tu m'en dire plus sur... ? » (standard) \\
& \eng{I'd like to know...} & « J'aimerais savoir... » (plus formel) \\
& \eng{What does the course involve?} & « En quoi consiste la formation ? » \\
& \eng{How long does it last?} & « Combien de temps dure-t-elle ? » \\
& \eng{Who is eligible to apply?} & « Qui est éligible pour postuler ? » \\
\hline
\cle{Donner une information} & \eng{The course covers...} & « La formation porte sur... » \\
& \eng{It lasts for two weeks.} & « Elle dure deux semaines. » \\
& \eng{Applicants must have...} & « Les candidats doivent avoir... » \\
& \eng{The deadline is June 15th.} & « La date limite est le 15 juin. » \\
& \eng{It is organised by [ONG].} & « Elle est organisée par [ONG]. » (passif) \\
\hline
\cle{Relancer / Réagir} & \eng{Really? / Is that so?} & « Vraiment ? / Ah bon ? » (surprise/intérêt) \\
& \eng{That sounds interesting.} & « Cela a l'air intéressant. » \\
& \eng{Could you clarify...?} & « Pourriez-vous préciser... ? » \\
& \eng{What about the fees?} & « Et pour les frais ? » (relance) \\
\hline
\cle{Ouvrir / Conclure} & \eng{Hello, I'm calling about...} & Ouverture téléphonique/formelle. \\
& \eng{Thank you for your help.} & Remerciement de clôture. \\
& \eng{I'll think about it and get back to you.} & « Je réfléchis et je vous recontacte. » \\
& \eng{Have a nice day!} & Formule de politesse finale. \\
\hline
\end{tabularx}
\end{center}

\courssub{Points de grammaire à maîtriser}
\begin{itemize}
\item \textbf{Modaux de politesse :} \eng{Can/Could you...?}, \eng{Would you mind...?}, \eng{I would like to...} (\cle{would like} + infinitif, pas \eng{like} seul).
\item \textbf{Passif (Present/Past Simple) :} \eng{The course is organised by... / It was funded by...} (Focus sur l'action, pas l'agent).
\item \textbf{Questions indirectes :} \eng{Could you tell me \textbf{when it starts}?} (Ordre sujet-verbe, pas d'auxiliaire \eng{do/does}).
\item \textbf{Prépositions :} \eng{enrol \textbf{in}}, \eng{register \textbf{for}}, \eng{participate \textbf{in}}, \eng{eligible \textbf{for}}.
\item \textbf{Futur proche (intention) :} \eng{I am going to apply / I am planning to attend.}
\end{itemize}

\courssub{Méthode express : Réussir un échange oral}
\begin{enumerate}
\item \textbf{Analyser le rôle :} Qui suis-je ? (étudiant, parent, journaliste) À qui je parle ? (secrétaire, formateur, ami) $\rightarrow$ Niveau de langue (formel/informel).
\item \textbf{Préparer les mots-clés :} Noter 5--7 mots anglais indispensables (dates, noms, verbes d'action).
\item \textbf{Structurer l'échange :} Ouverture $\rightarrow$ Questions/Infos $\rightarrow$ Relances $\rightarrow$ Conclusion.
\item \textbf{Jouer le jeu :} Regarder son partenaire, sourire, utiliser les « fillers » (\eng{Well...}, \eng{Let me see...}) pour gagner du temps.
\item \textbf{Vérifier :} Ai-je donné/obtenu les infos clés ? (Quoi, Qui, Quand, Où, Combien).
\end{enumerate}


\begin{aretenir}[Checklist avant le Bac -- Lesson 40]
$\square$ Je connais le vocabulaire de la démocratie et de la formation (15 mots minimum).\\
$\square$ Je sais demander une info poliment (\eng{Could you tell me...? / I'd like to know...}).\\
$\square$ Je sais donner une info claire (sujet + verbe + complément, passif si besoin).\\
$\square$ Je maîtrise 3 formules pour relancer / montrer l'intérêt (\eng{Really?}, \eng{That sounds...}, \eng{What about...?}).\\
$\square$ Je sais ouvrir (\eng{Hello, I'm calling about...}) et fermer (\eng{Thank you, goodbye}) un échange.\\
$\square$ Je place correctement les prépositions : \eng{enrol in}, \eng{register for}, \eng{eligible for}.\\
$\square$ Je forme correctement les questions indirectes (ordre affirmatif : \eng{...when it starts}).\\
$\square$ Je prononce correctement les mots-clés : \eng{democracy} (« di-mo-kra-si »), \eng{workshop} (« wor-kchop »), \eng{certificate} (« ser-ti-fi-kat »).\\
$\square$ J'utilise l'intonation montante pour les questions ouvertes, descendante pour les questions fermées/affirmations.\\
$\square$ Je suis capable de jouer un rôle complet (3--4 minutes) sans lire mes notes.\\
$\square$ Je gère mon stress : je respire, j'écoute, je reformule si je ne comprends pas (\eng{Sorry, could you repeat?}).
\end{aretenir}

\findecours

\end{document}
